% Embedded Linux Top 20, audited kit cookbook: illustrated edition
% Build:  python build.py            (PDF + HTML)
%         python build.py --check sections/p07.tex   (compile one lab alone)
\documentclass[11pt,a4paper]{article}
\usepackage[utf8]{inputenc}
\usepackage[T1]{fontenc}
\usepackage{lmodern}
\usepackage{microtype}
\usepackage[margin=21mm,top=24mm,bottom=24mm]{geometry}
\usepackage{xcolor}
\usepackage{graphicx}
\usepackage{booktabs}
\usepackage{tabularx}
\usepackage{array}
\usepackage{enumitem}
\usepackage{float}
\usepackage{caption}
\usepackage{listings}
\usepackage{tcolorbox}
\usepackage{tikz}
\usepackage[european]{circuitikz}
% Shared TikZ and circuitikz setup, used by main.tex and by the standalone
% figure wrapper that build.py generates for the HTML/SVG output.
\usetikzlibrary{arrows.meta,positioning,calc,shapes.geometric,shapes.misc,shapes.symbols,fit,backgrounds,chains,decorations.pathreplacing,decorations.markings,patterns,shadows,matrix}

% ---------------------------------------------------------------- colours
\definecolor{hwcol}{RGB}{255,236,179}   % hardware, boards, sensors
\definecolor{kercol}{RGB}{197,222,255}  % kernel space
\definecolor{usrcol}{RGB}{205,240,205}  % user space
\definecolor{extcol}{RGB}{228,228,228}  % external systems, cloud, PC
\definecolor{fwcol}{RGB}{245,210,235}   % MCU firmware
\definecolor{notecol}{RGB}{255,250,205} % notes
\definecolor{linecol}{RGB}{60,60,60}

% ---------------------------------------------------------------- block diagrams
\tikzset{
  every picture/.style={font=\sffamily\small, line width=0.6pt, >=Latex},
  blk/.style={draw=linecol, rounded corners=2pt, align=center, minimum height=8mm,
              inner sep=4pt, text width=28mm, fill=white},
  hw/.style={blk, fill=hwcol},
  kern/.style={blk, fill=kercol},
  user/.style={blk, fill=usrcol},
  ext/.style={blk, fill=extcol},
  fw/.style={blk, fill=fwcol},
  wide/.style={text width=44mm},
  narrow/.style={text width=18mm},
  layer/.style={draw=linecol, dashed, rounded corners=3pt, inner sep=6pt, fill=none},
  layerlabel/.style={anchor=north west, font=\sffamily\scriptsize\bfseries, inner sep=2pt},
  flow/.style={->, thick, draw=linecol},
  bus/.style={<->, thick, draw=linecol},
  lbl/.style={font=\sffamily\scriptsize, fill=white, inner sep=1.5pt, midway},
  irq/.style={->, thick, draw=red!70!black, dashed},
  % ---------------------------------------------------------------- UML
  umlclass/.style={draw=linecol, fill=white, align=left, inner sep=0pt, font=\sffamily\scriptsize},
  lifeline/.style={draw=linecol, dashed},
  actor/.style={draw=linecol, fill=extcol, rounded corners=2pt, align=center, minimum height=7mm, inner sep=4pt, font=\sffamily\small},
  msg/.style={->, draw=linecol, thick, font=\sffamily\scriptsize},
  reply/.style={->, draw=linecol, dashed, font=\sffamily\scriptsize},
  activation/.style={draw=linecol, fill=kercol, inner sep=0pt, minimum width=2.4mm},
  state/.style={draw=linecol, rounded corners=3mm, fill=usrcol, align=center, minimum height=8mm, inner sep=5pt},
  initial/.style={circle, fill=linecol, inner sep=0pt, minimum size=2.4mm},
  final/.style={circle, draw=linecol, fill=white, inner sep=0pt, minimum size=3.4mm,
                path picture={\fill[linecol] (path picture bounding box.center) circle (1mm);}},
  trans/.style={->, draw=linecol, thick, font=\sffamily\scriptsize},
  umlnote/.style={draw=linecol, fill=notecol, align=left, font=\sffamily\scriptsize, inner sep=3pt},
  % ---------------------------------------------------------------- bench art
  board/.style={draw=linecol, fill=green!25!black!60, rounded corners=1.5mm, minimum height=20mm, minimum width=32mm,
                text=white, font=\sffamily\bfseries\small, align=center},
  hat/.style={board, fill=blue!40!black!60},
  breadboard/.style={draw=linecol, fill=white, minimum height=12mm, minimum width=40mm, font=\sffamily\scriptsize},
  pinrow/.style={draw=linecol, fill=black, minimum height=1.6mm, inner sep=0pt},
  cable/.style={line width=1.2pt},
  usbcable/.style={line width=1.4pt, draw=black!70},
  ledsym/.style={circle, draw=linecol, minimum size=3mm, inner sep=0pt},
}

\usepackage{fancyhdr}
\usepackage[hidelinks,pdftitle={Embedded Linux Top 20: audited kit cookbook, illustrated edition},pdfauthor={Joseph Ambrose Pagaran}]{hyperref}

% ---------------------------------------------------------------- unicode helpers
\DeclareUnicodeCharacter{2192}{\ensuremath{\rightarrow}}
\DeclareUnicodeCharacter{2190}{\ensuremath{\leftarrow}}
\DeclareUnicodeCharacter{2264}{\ensuremath{\leq}}
\DeclareUnicodeCharacter{2265}{\ensuremath{\geq}}
\DeclareUnicodeCharacter{2248}{\ensuremath{\approx}}
\DeclareUnicodeCharacter{03A9}{\ensuremath{\Omega}}
\DeclareUnicodeCharacter{2126}{\ensuremath{\Omega}}
\DeclareUnicodeCharacter{2026}{\ldots}
\DeclareUnicodeCharacter{2013}{-}
\DeclareUnicodeCharacter{2014}{-}
\DeclareUnicodeCharacter{00D7}{\ensuremath{\times}}
\DeclareUnicodeCharacter{00B1}{\ensuremath{\pm}}

% ---------------------------------------------------------------- layout
\setlength{\emergencystretch}{3em}
\hyphenpenalty=800
\tolerance=1500
\setlength{\parskip}{4pt}
\setlength{\parindent}{0pt}
\captionsetup{font=small,labelfont=bf,skip=4pt}
\setlist{itemsep=2pt,topsep=3pt,parsep=0pt}
\pagestyle{fancy}
\fancyhf{}
\fancyhead[L]{\sffamily\small Embedded Linux Top 20, audited kit-only}
\fancyhead[R]{\sffamily\small\nouppercase{\benchmark}}
\newcommand{\benchmark}{\leftmark}
\fancyfoot[C]{\sffamily\small\thepage}
\renewcommand{\headrulewidth}{0.3pt}

% ---------------------------------------------------------------- listings
\lstset{basicstyle=\ttfamily\footnotesize,columns=fixed,basewidth=0.48em,keepspaces=true,
  breaklines=true,frame=single,rulecolor=\color{gray!60},backgroundcolor=\color{gray!6},
  xleftmargin=3pt,framexleftmargin=3pt,showstringspaces=false,tabsize=4,
  commentstyle=\color{green!40!black},keywordstyle=\color{blue!60!black}\bfseries,
  stringstyle=\color{red!50!black},extendedchars=false,aboveskip=6pt,belowskip=6pt}
\lstnewenvironment{asciiart}{\lstset{language={},frame=none,backgroundcolor=\color{white},
  breaklines=false,basicstyle=\ttfamily\scriptsize,keywordstyle=,commentstyle=,stringstyle=}}{}
\lstnewenvironment{ccode}{\lstset{language=C}}{}
\lstnewenvironment{shellcode}{\lstset{language=bash}}{}
\lstnewenvironment{pycode}{\lstset{language=Python}}{}
\lstnewenvironment{makecode}{\lstset{language=make}}{}
\lstnewenvironment{plaincode}{\lstset{language={}}}{}
\lstnewenvironment{dtscode}{\lstset{language={},morecomment=[l]{//},morecomment=[s]{/*}{*/}}}{}
\lstnewenvironment{yamlcode}{\lstset{language={},morecomment=[l]{\#}}}{}
% AT command transcripts: no language, but comments after # stay grey
\lstnewenvironment{atcode}{\lstset{language={},morecomment=[l]{\#}}}{}

% ---------------------------------------------------------------- authoring macros
% \lab{NN}{Title}{Host and exclusive part}{What it owns}
\newcommand{\project}[4]{\clearpage\section[P#1: #2]{#2}\label{sec:p#1}%
  \renewcommand{\benchmark}{Lab P#1}%
  {\color{black!60}\sffamily\small Host: #3 \quad|\quad Owns: #4}\par\medskip}
% the three coloured boxes of the source book
\newtcolorbox{unique}{colback=green!5,colframe=green!35!black,title=Why this lab is unique,
  fonttitle=\sffamily\bfseries\small,left=3pt,right=3pt,top=2pt,bottom=2pt}
\newtcolorbox{blueprint}{colback=blue!5,colframe=blue!45!black,title=What the 198-page blueprint got wrong,
  fonttitle=\sffamily\bfseries\small,left=3pt,right=3pt,top=2pt,bottom=2pt}
\newtcolorbox{kit}{colback=yellow!8,colframe=orange!75!black,title=Kit from the bin,
  fonttitle=\sffamily\bfseries\small,left=3pt,right=3pt,top=2pt,bottom=2pt}
% a red rule box for the non-negotiable rules and teardown
\newtcolorbox{rulebox}[1]{colback=red!4,colframe=red!55!black,title=#1,
  fonttitle=\sffamily\bfseries\small,left=3pt,right=3pt,top=2pt,bottom=2pt}
% note box
\newtcolorbox{note}[1]{colback=yellow!8,colframe=orange!75!black,title=#1,fonttitle=\sffamily\bfseries,
  left=3pt,right=3pt,top=2pt,bottom=2pt}
% numbered steps
\newenvironment{steps}{\begin{enumerate}[label=\textbf{Step \arabic*.},leftmargin=*,itemsep=5pt]}{\end{enumerate}}
% \diagram{figure-name}{Caption}: inputs figures/<name>.tex, scaled down only when too wide
\newsavebox{\diagbox}
\newcommand{\diagram}[2]{\begin{figure}[H]\centering
  \sbox{\diagbox}{\input{figures/#1.tex}}%
  \ifdim\wd\diagbox>\linewidth\resizebox{\linewidth}{!}{\usebox{\diagbox}}\else\usebox{\diagbox}\fi
  \caption{#2}\label{fig:#1}\end{figure}}

\title{\sffamily\bfseries Embedded Linux Top 20\\[2pt]\large Audited kit cookbook, illustrated edition}
\author{Joseph Ambrose Pagaran}
\date{September 2026}

\begin{document}
\maketitle
\thispagestyle{empty}
\section*{About this edition}
\addcontentsline{toc}{section}{About this edition}

This volume re-publishes the audited kit cookbook with the diagrams it never had. The text you are reading is the cookbook: its twenty labs, its verdicts, its corrections and its refusals. Nothing in the audit has been softened, and no hardware appears here that is not in the bin the cookbook was written against.

What this edition adds is drawing. Every lab now carries an architecture diagram, a schematic with a wiring table of real pin numbers, a bench layout, a UML view of the software or the procedure, an ASCII data flow and a short pitfalls list. Those six additions are the work of this edition. The content, the verdicts and the corrections are the cookbook's own.

The cookbook itself starts from two sources: the kit-only 20-lab split, and a 198-page collection of advanced architectural blueprints. The draft had the right ambition, including QMI instead of PPP, PSM and eDRX, sensor-hub modes, the PPK2 on the modem rail and store-and-forward. It also contained pairings that cannot be built from this bin. This book keeps the ambition and throws out the impossible wiring.

Hosts run Raspberry Pi OS Bookworm; the NanoPi runs FriendlyElec Ubuntu.

\diagram{front_map}{The twenty labs, grouped by the five parts of this book, with the exclusive part each one occupies. The eight labs drawn with a heavy border seat a 40-pin HAT, and only one of those can sit on a given Pi at a time.}

\section*{Non-negotiable lab rules}
\addcontentsline{toc}{section}{Non-negotiable lab rules}

\begin{rulebox}{Non-negotiable lab rules}
\begin{enumerate}
\item \textbf{No new parts.} HAT antennas are in-scope. A data SIM is assumed only for P01 to P03 and P13. Without a SIM those labs stop at AT bring-up, and that is still a pass.
\item \textbf{HAT exclusivity.} SIM7600E-H, SIM7070G, SIM7020E, MCC 118, JOY-iT Explorer700 and the Waveshare 3.5 inch LCD (A) each own the 40-pin header. Fit one. The official DSI touchscreen does not consume the header.
\item \textbf{NanoPi is 24-pin.} It cannot accept any 40-pin HAT. Do not seat SIM7020, SIM7070, SIM7600, MCC 118, Explorer700 or the 3.5 inch LCD on the NEO Air.
\item \textbf{3.3 V GPIO.} Pi, NanoPi and Nucleo I/O are 3.3 V. POW-BB is a labelled rail splitter.
\item \textbf{One physics per lab.} IKS4A1, IKS5A1, STWIN, ADXL345, VL53L8CX and MCC 118 are six different instruments, never interchangeable.
\item \textbf{Linux owns the system.} Nucleo, STWIN, ESP32 and ESP8266 speak USB-CDC, UART or I2C to a host that owns storage, time, networking and the UI. A bare-metal sketch is allowed only when a Linux host still records the stream (P18).
\end{enumerate}
\end{rulebox}

Build one lab at a time. Tear the 40-pin HAT off before the next lab.

\section*{The audit}
\addcontentsline{toc}{section}{The audit}

Every lab in this book was checked against the parts that actually exist in the bin. The verdict table is the result. A status of Buildable means the wiring closes with the kit in hand, not that the lab is easy.

\begin{table}[H]\centering\small
\begin{tabular}{p{8mm}p{44mm}p{40mm}p{52mm}}
\toprule
Lab & Primary capability & Host and exclusive part & Status \\
\midrule
P01 & LTE Cat-4 default route plus shock events & Pi 4 + SIM7600 + ADXL & Buildable \\
P02 & NB-IoT PSM/eDRX CoAP & Pi 3 + SIM7020 & Buildable, SIM required for RF \\
P03 & Cat-M MQTT plus GNSS stamp & Pi 3B+ + SIM7070 & Buildable; lock AT+CMNB=1 \\
P04 & 6 kHz vibration and ultrasound CDC & Pi 4 + STWIN + DSI & Buildable; firmware chooses CDC \\
P05 & 10 V 100 kS/s analog & Pi 4 + MCC 118 & Buildable via daqhats, not IIO \\
P06 & 8x8 ToF occupancy kiosk & Pi 3 + LCD + Nucleo + 53L8 & Buildable; two hosts \\
P07 & Consumer 9-DoF plus RH and Qvar & Nucleo + IKS4A1 + Pi CDC & Buildable \\
P08 & Industrial high-g and 4060 hPa & Nucleo + IKS5A1 + Pi CDC & Buildable; not ISM330DHCX \\
P09 & Allwinner headless I2C hub & NanoPi + ADXL + UART0 & Buildable; 24-pin only \\
P10 & Source-measure energy figures & Pi 4 + PPK2 + ESP32 DUT & Buildable; isolate DUT USB \\
P11 & RTC, OLED, IR and DTO console & Pi 3B+ + Explorer700 & Buildable \\
P12 & Operator glass plus Mosquitto & Pi 4 + official DSI + keyboard & Buildable; no 40-pin HAT \\
P13 & Signalling failover 4G to NB & Pi 4+7600 and Pi 3+7020 & Buildable; two hosts \\
P14 & Serial provisioning jig & Pi 3 + USB-TTL + ESP32/8266 & Buildable \\
P15 & Mixed-voltage GPIO discipline & POW-BB + LEDs + any host & Buildable \\
P16 & Wi-Fi/BLE sidecar & Pi 4 + ESP32 & Buildable \\
P17 & AT-command air-gap modem & NanoPi + ESP8266, onboard Wi-Fi off & Buildable \\
P18 & High-rate USB-CDC recorder & Pi 4 + bare Nucleo, EKF optional & Buildable; Linux still records \\
P19 & Offline SPI-LCD tilt meter & Pi 3 + 3.5 inch LCD + ADXL & Buildable; LCD owns header \\
P20 & Four-host fleet health & All Linux boards + SSH & Buildable \\
\bottomrule
\end{tabular}
\caption{The audit verdict for all twenty labs.}
\end{table}

\begin{rulebox}{Do not build these pairings}
Errors found in the 198-page blueprint.
\begin{itemize}
\item \textbf{NanoPi plus any 40-pin HAT} (draft P4, P12). NEO Air is a 24-pin header. SIM7020E does not seat. NB-IoT moves to Pi 3, our P02.
\item \textbf{Explorer700 plus SIM7600 on one Pi} (draft P18). Two 40-pin HATs. Split across hosts, or drop one HAT.
\item \textbf{IKS5A1 hosting an ISM330DHCX MLC} (draft P13). That IMU is on the STWIN.box. IKS5A1 has ISM330IS with the ISPU, and ISM6HG256X. Our P08 uses the chips that are actually on IKS5A1.
\item \textbf{MCC 118 as a mainline IIO plus DMABUF and io\_uring device} (draft P2, P17). The supported path is SPI plus the daqhats userspace library. There is no IIO\_BUFFER\_DMABUF\_ATTACH\_IOCTL for this HAT.
\item \textbf{Intel MKL on a Pi 4} (draft P1). Use NumPy, FFTW or OpenBLAS.
\item \textbf{SIM7070 stacked on the Nucleo as a HAT} (draft P3). The modem is a Pi HAT. The Nucleo talks UART to a Pi that owns the modem, or to breakout pins with flying leads, never as a stacked Arduino HAT.
\item \textbf{Yocto-from-scratch as a week-one lab} (draft P12). A valid stretch goal on a separate x86 build machine. The kit lab is FriendlyElec Ubuntu on eMMC, our P09.
\item \textbf{VCP as a property of USB-C} (draft P1). USB-C is the connector. CDC-ACM appears only after the STWIN firmware declares it.
\item \textbf{stsw-img042 V4L2 ToF node on Raspberry Pi OS} (draft P8). ST's Linux ToF stack is not a drop-in Bookworm camera device. Stream 64 ranges over the Nucleo USB-CDC instead, our P06.
\item \textbf{Draft P1 and P20 both owning STWIN vibration uplink.} Keep P04 as USB raw plus FFT and do not add a second STWIN product lab.
\end{itemize}
\end{rulebox}

Ideas from the draft that are kept: QMI and RNDIS instead of PPP for Cat-4; PSM and eDRX on NB-IoT; GNSS powered down during LTE-M transmission; PPK2 in the DUT power rail; DS3231 as hwclock; Explorer700 overlays; anomaly-only uplink as a software policy on P04, not a second project; isolcpus for the FFT thread.

\section*{Inventory and interface map}
\addcontentsline{toc}{section}{Inventory and interface map}

\begin{table}[H]\centering\small
\begin{tabular}{p{36mm}p{24mm}p{88mm}}
\toprule
Item & PN & Role \\
\midrule
RPi 4 / 3B+ / 3 & RPi & Bookworm hosts \\
NanoPi NEO Air & FriendlyElec & H3, 512 MB, 8 GB eMMC, Wi-Fi and BT, \textbf{24-pin} \\
NUCLEO-H7A3ZI-Q & ST & STM32H7A3 USB-CDC sensor brain \\
X-NUCLEO-IKS4A1 & ST & LSM6DSO16IS, LSM6DSV16X, LIS2DUXS12, LIS2MDL, LPS22DF, SHT40, STTS22H, Qvar \\
X-NUCLEO-IKS5A1 & ST & ISM6HG256X, ISM330IS, IIS2DULPX, IIS2MDC, ILPS22QS \\
X-NUCLEO-53L8A1 & ST & VL53L8CX 8x8 ToF \\
STEVAL-STWINBX1 & ST & IIS3DWB 6 kHz, ISM330DHCX with MLC, IMP34DT05, IMP23ABSU \\
nRF PPK2 & Nordic & 100 kS/s current source-measure \\
SEN0032 ADXL345 & DFRobot & I2C 0x53 3-axis, INT1 available \\
SIM7600E-H 4G HAT & Waveshare & LTE Cat-4 plus GNSS, USB RNDIS or QMI \\
SIM7070G HAT & Waveshare & Cat-M, NB-IoT or GPRS, plus GNSS \\
SIM7020E HAT & Waveshare & NB-IoT only \\
MCC 118 & Digilent/MCC & 8 ch, 10 V, 12-bit, 100 kS/s, SPI \\
3.5 inch RPi LCD (A) & Waveshare & 480x320 SPI plus resistive touch \\
RB-Explorer700 & JOY-iT & OLED, DS3231, BMP280, PCF8591, PCF8574, joystick, IR, buzzer \\
SBC-NodeMCU-ESP32 & JOY-iT & Wi-Fi plus BLE coprocessor \\
SBC-ESP8266-PROG & JOY-iT & AT-modem sidecar \\
USB/TTL cable & Renkforce & 3.3 V UART \\
SBC-POW-BB and LEDs & JOY-iT & Rails and semaphores \\
Touchscreen and keyboard & RPi & Operator HMI, DSI, not a HAT \\
\bottomrule
\end{tabular}
\caption{The bin. No lab in this book adds a part to it.}
\end{table}

\begin{table}[H]\centering\small
\begin{tabular}{p{50mm}p{50mm}p{40mm}}
\toprule
Device & Board & Address \\
\midrule
ADXL345 & SEN0032 & 0x53 \\
LIS2MDL / IIS2MDC & IKS4 / IKS5 / STWIN & 0x1E \\
SHT40AD1B & IKS4A1 & 0x44 \\
STTS22H & IKS4 / STWIN & 0x38 \\
LPS22DF & IKS4A1 & 0x5C \\
LSM6DSO16IS / LSM6DSV16X & IKS4A1 & 0x6A \\
ISM6HG256X & IKS5A1 & 0x6A \\
ISM330IS & IKS5A1 & 0x6B \\
ILPS22QS & IKS5 / STWIN & 0x5C \\
IIS2DULPX & IKS5A1 & 0x19 \\
DS3231 & Explorer700 & 0x68 \\
BMP280 & Explorer700 & 0x76 or 0x77 \\
PCF8591 & Explorer700 & 0x48 \\
PCF8574 & Explorer700 & 0x20 to 0x27 \\
VL53L8CX & 53L8A1 I2C mode & 0x29 \\
\bottomrule
\end{tabular}
\caption{Default I2C map, 7-bit addresses.}
\end{table}

\begin{asciiart}
40-pin XOR (one per Pi):  SIM7600 | SIM7070 | SIM7020 | MCC118 | Explorer700 | 3.5" LCD
May coexist with a HAT:   official DSI touchscreen, USB gadgets
USB-host gadgets:         STWIN, Nucleo, PPK2, ESP32, ESP8266, Renkforce TTL
NanoPi 24-pin only:       I2C0 / UART1 / SPI0 / GPIO --- no 40-pin HAT
Arduino-stack XOR:        IKS4A1 | IKS5A1 | 53L8A1  (one shield on the Nucleo)
\end{asciiart}

\section*{How this book is built}
\addcontentsline{toc}{section}{How this book is built}

The whole volume is one command:

\begin{shellcode}
python build.py
\end{shellcode}

\texttt{build.py} renders each figure on its own through \texttt{latex} and \texttt{dvisvgm}, then runs \texttt{pdflatex} over \texttt{main.tex} to make the PDF. The same sources are parsed a second time and converted to one self-contained HTML file, which is why the LaTeX in the section files is a small, fixed subset: the converter understands exactly that subset and nothing else. Figures reach the HTML as inline SVG, so the web version carries the same drawings as the page.

\texttt{lint.py} enforces the house rules on the section sources: no em or en dashes, ASCII only inside code blocks, no line too long to print, and the required subsection skeleton in every lab. A single lab can be compiled and checked alone with \texttt{python build.py --check sections/pNN.tex}, which is how each chapter in this edition was proofed before it joined the book.

\clearpage
\tableofcontents
\clearpage

\part{Cellular labs}
\project{01}{LTE Cat-4 motion-triggered gateway}{Raspberry Pi 4 + SIM7600E-H + ADXL345}{the only Cat-4 default route, plus shock-event compression}

\begin{unique}
Owns the only Cat-4 \emph{default route} in the book, plus ADXL345 event compression. Cat-M and NB-IoT are P03 and P02.
\end{unique}

\begin{blueprint}
Draft P5 is the same radio used as a NAT router; draft P1 used Intel MKL. We use QMI/RNDIS for the WAN and NumPy for nothing heavier than a magnitude check. Do not also put GNSS tracking here, that is P03.
\end{blueprint}

\subsection*{Intent}

The Pi 4 gets a default IPv4 route through the SIM7600E-H over USB, as RNDIS or as QMI-WWAN. The ADXL345 publishes a shock event only when the acceleration magnitude leaves a dead-band around 1 g, so the uplink carries events rather than a sample stream. Three LEDs report the three states that matter on a bench: the modem registered, the default route moved to the modem, and an event fired.

The lab is a gateway, not a tracker and not a router. Two things are deliberately out of scope. GNSS belongs to P03, because a Cat-4 module that also chases a fix hides which subsystem failed. Network address translation for other hosts belongs to nothing in this book, because one SIM and one bench do not need it.

\begin{kit}
Pi 4, SIM7600E-H with LTE and GNSS antennas and a micro-SIM, SEN0032 ADXL345, three LK-LED10 modules on BCM 16/20/21, a 3 A USB-C supply, keyboard.
\end{kit}

\subsection*{System architecture}

\diagram{p01_arch}{The two planes of the lab. The user plane is the modem's USB network interface; the control plane is an AT character device on the same cable. The accelerometer never touches either.}

The SIM7600E-H presents several interfaces on one USB cable. One of them is a network device that carries the user plane, and another is a character device that accepts AT commands. Keeping those apart is the whole discipline of the lab: the WAN comes up through \texttt{dhcpcd} or ModemManager on the network device, and diagnostics go to the AT device. Point-to-point protocol is not used at all, because running it alongside QMI or RNDIS gives two things one modem cannot serve at once.

The event path is independent of the network path. A Python loop reads six registers over I2C at 20 Hz, computes a magnitude, and calls \texttt{mosquitto\_pub} when the magnitude leaves the dead-band. If the WAN is down the publish fails and the loop continues, which is the behaviour you want from a field gateway.

\subsection*{Wiring and schematic}

\diagram{p01_schematic}{The accelerometer on the HAT pass-through header and the three status LEDs. SDO left floating selects address 0x53.}

\begin{table}[H]\centering\small
\begin{tabular}{p{42mm}p{26mm}p{22mm}p{58mm}}
\toprule
Signal & Pi header pin & BCM GPIO & Note \\
\midrule
ADXL345 VCC & 1 & 3V3 & Never 5 V \\
ADXL345 SDA & 3 & GPIO2 & I2C1 data \\
ADXL345 SCL & 5 & GPIO3 & I2C1 clock \\
ADXL345 GND & 6 & GND & \\
ADXL345 SDO & & & Left floating, selects 0x53 \\
ADXL345 INT1 & 11 & GPIO17 & Optional, not needed for the magnitude poll \\
Green module, S1 & 36 & GPIO16 & CEREG registered. Lit when driven high \\
Yellow module, S1 & 38 & GPIO20 & Default route on usb0 or wwan0 \\
Red module, S1 & 40 & GPIO21 & Last shock event \\
Module grounds & 34 or 39 & GND & The breadboard ground rail, as in P15 \\
Pins 13, 15, 16 & & GPIO27, 22, 23 & Not free on this HAT. See the note below \\
\bottomrule
\end{tabular}
\caption{Wiring. The HAT carries a full 40-pin pass-through, so the accelerometer and the modem coexist on one header and pins 3 and 5 are untouched by the HAT. The indicator pins are not the ones the draft used; the note below says why.}
\end{table}

\begin{note}{The status triple cannot sit on BCM 27, 22 and 23 under this HAT}
The draft put the three indicators on header pins 13, 15 and 16. On the SIM7600E-H those are the modem's own control lines: GPIO27 is the ring indicator, GPIO22 is data terminal ready and GPIO23 is clear to send. An LED hung on any of them is a second driver on a line the modem is already using, and the failure it produces is not a dark LED but a modem that behaves oddly for reasons nothing in the log explains.

This is not peculiar to one board. On the SIM7020E of P02, GPIO27 is the module's power key, which the vendor scripts pulse to boot it. An indicator there would switch the modem on and off every time it reported a change of state, so the lab would interfere with the thing it is reporting.

So the triple is per lab rather than per book. The pins above are chosen from the high end of the header, where neither modem HAT documents a connection, with the module grounds going to pin 34 or 39. \textbf{Confirm them against the pinout of the HAT in front of you before wiring}, because this is a fact about one board and not about the Raspberry Pi. Some Waveshare HATs carry a jumper block that selects which Pi GPIO reaches each control line, and on a board that has one, removing three jumpers is the other way to free the original pins.
\end{note}

\begin{asciiart}
ADXL345                 Pi 40-pin (HAT pass-through pins 1/3/5/6 still available)
VCC ------------------- pin 1  3V3      NEVER 5V
GND ------------------- pin 6  GND
SDA ------------------- pin 3  GPIO2
SCL ------------------- pin 5  GPIO3
SDO floating => 0x53
INT1 optional --------- BCM17  (not required for the magnitude poll)

SIM7600 jumpers:
  VCCIO = 3V3
  PWR   = 3V3   (auto-on). Use PWR=D6 only if you want a GPIO power-key.
  UART  = B if you need ttyAMA0; prefer the HAT USB cable for AT + data.
USB cable: HAT USB-A/micro -> Pi USB-A  (user plane + AT channels)
Antennas: LTE main + GNSS. No antenna = no attach, and can stress the PA.
\end{asciiart}

\subsection*{Bench layout}

\diagram{p01_bench}{Bench layout. The short USB cable between the HAT and the Pi carries the user plane, so it is not optional. The 3 A supply is not optional either: a transmit burst on Cat-4 draws current the Pi cannot lend.}

Seat the HAT, fit both antennas and insert the SIM before applying power. Connect the HAT USB cable in the same pass. A modem that is powered with no antenna fitted can stress its own power amplifier, and a SIM inserted under power is not detected until the next boot.

\subsection*{Software design (UML)}

\diagram{p01_uml}{Bring-up sequence. The AT channel confirms registration, the network interface carries traffic, and the accelerometer loop is a separate process that survives a WAN outage.}

The publisher is a single loop with one piece of state, a cool-off timestamp. It is a token bucket in its simplest form: at most one event every two seconds, no matter how long the table keeps ringing. Without it, one fist produces a burst of twenty publishes and the value of the compression is lost.

\subsection*{Data flow (ASCII)}

\begin{asciiart}
  Pi 4                                            SIM7600E-H HAT
  +---------------------------------+             +------------------------+
  | p01_shock.py                    |             |  LTE Cat-4 radio       |
  |   i2c-1 -> 0x53 ADXL345 @20 Hz  |             |    |                   |
  |   |a| - 1g > 0.35 ?             |             |    v                   |
  |   |  yes -> mosquitto_pub ------------------> | usb0 / wwan0  (user)   |---> broker
  |   |  cool-off 2 s               |             |                        |
  |   v                             |   USB       | /dev/ttyUSB2  (AT) <--------+
  | gpio: GRN 16  YEL 20  RED 21    |<===========>|                        |     |
  +---------------------------------+             +------------------------+     |
        ^                                                                        |
        +--- AT+CSQ every 60 s, AT+CEREG? for the green LED --------------------+
\end{asciiart}

\subsection*{Steps}

\begin{steps}
\item \textbf{Prepare the host.} Flash Raspberry Pi OS Bookworm 64-bit. Enable I2C under \texttt{raspi-config}, Interface options. Use a 3 A supply.

\item \textbf{Assemble before power.} Seat the HAT, fit both antennas, insert the SIM, and connect the HAT USB cable, all before applying power.

\item \textbf{Identify the modem and the accelerometer.}
\begin{shellcode}
lsusb | grep -iE "sim|1e0e"
dmesg | tail -n 40
ls /dev/ttyUSB* /dev/ttyACM* 2>/dev/null
sudo apt-get update
sudo apt-get install -y minicom i2c-tools mosquitto-clients python3-smbus python3-rpi.gpio
sudo i2cdetect -y 1          # must show 53
\end{shellcode}

\item \textbf{Open the AT port.} It is commonly \texttt{/dev/ttyUSB2} on the SIM7600. If minicom shows nothing, try USB3 then USB1.
\begin{atcode}
sudo minicom -D /dev/ttyUSB2 -b 115200
AT
ATE1
AT+CPIN?
AT+CSQ
AT+COPS?
AT+CEREG?
AT+CGDCONT?
\end{atcode}

\item \textbf{Bring up the WAN.} Prefer RNDIS or QMI over the point-to-point protocol.
\begin{shellcode}
# RNDIS path: usb0 typically appears when the HAT USB cable is connected
ip link
sudo dhcpcd usb0 || sudo dhclient usb0
ip route
# If usb0 did not appear, load qmi_wwan / option and use ModemManager:
sudo apt-get install -y modemmanager libqmi-utils
mmcli -L
# Do not run PPP in parallel with QMI or RNDIS.
\end{shellcode}

\item \textbf{Confirm and lock the default route}, so the modem wins over wlan0.
\begin{shellcode}
ip route | grep default
# optional: give wlan0 a higher metric so LTE wins
sudo ip route del default dev wlan0
\end{shellcode}

\item \textbf{Run the event publisher}, saved as \texttt{/home/pi/p01\_shock.py}.
\begin{pycode}
import time, math, json, subprocess, smbus
bus = smbus.SMBus(1); A = 0x53
bus.write_byte_data(A, 0x2D, 0x08)   # measure
bus.write_byte_data(A, 0x31, 0x0B)   # full-res +/-16g

def g(lo, hi):
    v = (hi << 8) | lo
    if v & 0x8000:
        v -= 65536
    return v / 256.0

cool = 0
while True:
    b = bus.read_i2c_block_data(A, 0x32, 6)
    ax, ay, az = g(b[0], b[1]), g(b[2], b[3]), g(b[4], b[5])
    mag = math.sqrt(ax*ax + ay*ay + az*az)
    if abs(mag - 1.0) > 0.35 and time.time() > cool:
        payload = json.dumps({"mag": round(mag, 3), "ax": round(ax, 3)})
        subprocess.run(["mosquitto_pub", "-h", "test.mosquitto.org",
                        "-t", "lab/p01/shock", "-m", payload], check=False)
        cool = time.time() + 2
    time.sleep(0.05)
\end{pycode}

\item \textbf{Drive the LEDs from state, not from hope.} Green follows \texttt{AT+CEREG?}, yellow follows the presence of a default route on \texttt{usb0}, red follows the last publish. Poll \texttt{AT+CSQ} every 60 s and log it.
\end{steps}

\subsection*{Acceptance test}

\begin{itemize}
\item A fist on the table produces exactly one MQTT JSON message and lights the red LED.
\item \texttt{ping -I usb0 8.8.8.8} succeeds, or the operator's own DNS address does.
\item \texttt{ip route | grep default} names the modem interface, not wlan0.
\item Outdoors, \texttt{AT+CGPS=1} then \texttt{AT+CGPSINFO} returns non-blank fields. That is a side check on the hardware, not this lab's product.
\end{itemize}

\subsection*{Practices}

User plane on USB, AT on a ttyUSB. Token-bucket the publisher, which is what the two-second cool-off is. Log \texttt{AT+CSQ} every 60 s so a weak-signal night is visible the next morning. Do not seat the MCC 118 or the 3.5 inch LCD at the same time: one 40-pin HAT per host.

\subsection*{Pitfalls}

\begin{itemize}
\item \textbf{No antenna fitted.} The module will not attach, and transmitting into an open port stresses the power amplifier. Fit both antennas before power.
\item \textbf{A 2.5 A supply.} The Pi 4 alone is close to that budget. A Cat-4 transmit burst on top of it produces undervoltage warnings in \texttt{dmesg} and random resets that look like software faults.
\item \textbf{Running the point-to-point protocol next to QMI.} Two user planes, one radio. The symptom is a route that appears and then carries nothing.
\item \textbf{Picking the wrong AT device.} Several \texttt{ttyUSB} nodes appear. The wrong one is silent rather than an error, so try USB2, then USB3, then USB1.
\item \textbf{Publishing every sample.} Without the cool-off, one impact produces a burst and the dead-band logic buys nothing.
\item \textbf{5 V on the accelerometer.} The SEN0032 goes to pin 1, never pin 2.
\end{itemize}

\subsection*{Sources}

\begin{itemize}
\item Waveshare SIM7600E-H 4G HAT wiki, \url{https://www.waveshare.com/wiki/SIM7600E-H_4G_HAT}
\item SIMCom SIM7600 series AT command manual (module vendor documentation)
\item DFRobot SEN0032 ADXL345 product wiki, \url{https://wiki.dfrobot.com/}
\item Analog Devices ADXL345 data sheet, register map 0x2D, 0x31, 0x32
\item ModemManager and libqmi documentation, \url{https://www.freedesktop.org/wiki/Software/ModemManager/}
\item Raspberry Pi hardware documentation, 40-pin header, \url{https://www.raspberrypi.com/documentation/computers/raspberry-pi.html}
\end{itemize}

\project{02}{NB-IoT PSM/eDRX CoAP field node}{Raspberry Pi 3 + SIM7020E}{NB-IoT and the power-saving duty cycle}

\begin{unique}
Owns NB-IoT and the PSM/eDRX duty cycle. No default IPv4 WAN, no GNSS, no MQTT. The opposite of P01.
\end{unique}

\begin{blueprint}
Draft P4 seated the SIM7020E on the NanoPi. That header is 24-pin. NB-IoT lives on the Pi 3 in this book.
\end{blueprint}

\subsection*{Intent}

The SIM7020E sits on a Pi 3 as a power-saving NB-IoT endpoint. Every 10 minutes the node wakes, sends an approximately 20-byte CoAP/UDP datagram, and returns to idle. Acceptance is the duty cycle, not a continuous ping. A green LED on BCM27 reports one fact only, that \texttt{AT+CEREG?} says the module is registered on a cell.

Nothing else belongs here. There is no default route on the Pi, because the UDP socket lives inside the module and is opened with \texttt{AT+CSOC}. There is no GNSS, because that is P03. There is no MQTT, because NB-IoT is not a phone network and the public broker is usually unreachable from an NB APN. Energy figures are measured in P10 with the PPK2, not estimated here.

\begin{kit}
Pi 3, SIM7020E HAT with an NB antenna and an NB-IoT SIM, green LED on BCM27.
\end{kit}

\begin{note}{Without a SIM}
A data SIM is assumed only for P01 to P03 and P13. Without one this lab stops at AT bring-up, and that is still a pass. Steps 1 and 2 up to \texttt{AT+CPIN?} exercise the UART, the overlay and the HAT seating, which is most of what goes wrong.
\end{note}

\subsection*{System architecture}

\diagram{p02_arch}{One serial line and one GPIO. The Linux host holds no IP address for this path: the UDP socket is created inside the module by \texttt{AT+CSOC}, so the kernel network stack is not involved at all.}

The whole path is a character device. The Pi opens \texttt{/dev/ttyAMA0} at 115200, writes AT commands and reads responses, and that is the entire interface between Linux and the radio. The module holds the socket, the APN context and the power-saving state. This is the opposite arrangement from P01, where the modem presents a network interface and the kernel routes packets through it.

The consequence is that the Pi cannot ping anything over this link, and should not try. The evidence that the node works is a datagram arriving at the CoAP endpoint once per wake, plus \texttt{AT+CEREG?} and \texttt{AT+CSQ} showing a cell between wakes. A continuous ping is a fail, because it also prevents the module from ever entering power-saving mode.

\begin{table}[H]\centering\small
\begin{tabular}{p{34mm}p{112mm}}
\toprule
Command & What it proves \\
\midrule
\texttt{AT+CPIN?} & The SIM is seated and unlocked \\
\texttt{AT+CBAND=20} & The band is set. EU deployments are commonly B8 or B20, so this is operator-specific \\
\texttt{AT+COPS=0} & Automatic operator selection, the normal case for an NB SIM \\
\texttt{AT+CEREG?} & Registration on a cell. This is what the green LED follows \\
\texttt{AT+CGATT?} & The packet service is attached, which is a step beyond registration \\
\texttt{AT+CSQ} & Signal quality, the number to write in the lab book with the antenna in place \\
\texttt{AT+CGDCONT?} & The APN context the module will actually use \\
\texttt{AT+CPSMS?} & The power-saving values as granted, which may differ from the request \\
\bottomrule
\end{tabular}
\caption{The bring-up commands and what each one settles. Run them in this order: a failure at any line makes the following lines meaningless.}
\end{table}

\subsection*{Wiring and schematic}

The electrical work is almost nothing. The HAT seats on the header and takes its power and its serial line from there. One LK-LED10 module is the only part on the breadboard, and it carries its own resistor, so nothing is added beside it. The configuration work, on the other hand, is the reason this section exists: on a Pi 3 the PL011 is wired to the Bluetooth controller by default, and the port that carries the AT dialogue only becomes \texttt{ttyAMA0} after the overlay moves Bluetooth out of the way.

\diagram{p02_schematic}{The HAT seats on the 40-pin header and reaches the Pi over the PL011 UART. The UART jumper is in the Waveshare B position, so the Pi controls the modem. The LED branch is the only flying wiring in the lab.}

\begin{table}[H]\centering\small
\begin{tabular}{p{40mm}p{24mm}p{22mm}p{62mm}}
\toprule
Signal & Pi header pin & BCM GPIO & Note \\
\midrule
HAT 5 V & 2 and 4 & & Supplied through the header when the HAT is seated \\
UART TXD0 & 8 & GPIO14 & Pi transmit, to the module receive pin \\
UART RXD0 & 10 & GPIO15 & Pi receive, from the module transmit pin \\
HAT GND & 6 & GND & \\
Green module, S1 & 13 & GPIO27 & CEREG registered. Lit when the pin is driven high \\
Green module, G & 9 & GND & The return. The module's resistor is inside it \\
NB antenna & & & On the NB port of the HAT, not a 4G antenna \\
\bottomrule
\end{tabular}
\caption{Wiring. The HAT occupies the whole 40-pin header, so the LED goes on the pass-through pins the HAT exposes. Confirm the pass-through on the silkscreen before soldering anything permanent.}
\end{table}

\begin{asciiart}
HAT seated on Pi 3. UART jumper = Waveshare "B" (Pi controls modem).
Disable Bluetooth to free PL011 ttyAMA0:
  /boot/firmware/config.txt   dtoverlay=disable-bt
  raspi-config -> Serial: login=OFF, hardware=ON
LED GRN BCM27 = CEREG registered.
Antenna on the NB port. No 4G antenna swap, the matching is different.
\end{asciiart}

\begin{rulebox}{Before the supply goes in}
The SIM7020E owns the 40-pin header: fit one HAT, and tear it off before the next lab. Pi I/O is 3.3 V, so the LED branch goes to a header pin and never to a 5 V rail. Fit the NB antenna before power, and fit the NB antenna, not the 4G one.
\end{rulebox}

\subsection*{Bench layout}

\diagram{p02_bench}{Bench layout. One HAT, one antenna, one LED. The CoAP endpoint is a host you control on the operator's network, because a public echo server is frequently unreachable from an NB-IoT APN.}

Seat the HAT with the power off, fit the NB antenna, insert the SIM, then apply power. The NB antenna and the 4G antenna are not interchangeable: the matching is different, and swapping them produces a module that reports a poor \texttt{AT+CSQ} and never completes the attach.

The rest of the bench is deliberately empty. One LED, one breadboard, one supply. There is no second HAT, because the SIM7020E owns the 40-pin header, and there is no USB gadget to watch, because this radio reaches Linux through the header UART rather than through a cable. If the bench looks too quiet for a day's work, that is the shape of an NB-IoT node: most of its life is spent doing nothing, on purpose.

\subsection*{Software design (UML)}

\diagram{p02_uml}{The duty cycle as a state machine. The long timer is the requested tracking-area update, the short one is the requested active time. Between them the module is idle and the Pi has nothing to do.}

The lab is the state machine, not the payload. Two timers are requested with \texttt{AT+CPSMS}: a long tracking-area update interval that sets how often the module must announce itself, and a short active time that sets how long it stays reachable after each exchange. The encoded bitmaps are requests, not guarantees. The network grants what it grants, which is why \texttt{AT+CPSMS?} is read back and why the fallback in step 5 exists.

The software on the Pi is correspondingly small. There is one process, and its loop has four phases: read the registration state, open the socket, send, sleep. It holds no buffer worth protecting and no connection worth keeping, so a restart costs one missed datagram at most. Resist the urge to add a retry queue before the duty cycle itself is measured, because an aggressive retry is exactly the behaviour that keeps the radio awake and makes the lab prove the opposite of its intent.

One consequence is worth stating plainly before the steps begin. Everything the Pi knows about this radio arrives as text on one file descriptor, so every diagnostic in the lab is a command you can type by hand in \texttt{minicom} first and automate afterwards. Type it by hand first.

\subsection*{Data flow (ASCII)}

\begin{asciiart}
  Pi 3                                        SIM7020E HAT
  +------------------------------+            +---------------------------+
  | p02_coap.py                  |            |  NB-IoT radio, band 20    |
  |   every 10 minutes:          |            |    |                      |
  |     AT+CSOC=1,2,1            |            |    v                      |
  |     AT+CSOCON=0,5683,"..."   |  ttyAMA0   | UDP socket INSIDE the     |
  |     AT+CSOSEND=0,8,"50494E47"|<==========>| module (no Linux route)   |--> CoAP
  |   between sends: nothing     |  115200    |                           |    5683
  |                              |            | PSM: idle, radio off      |
  | gpio: GRN BCM27 <- CEREG     |            |                           |
  +------------------------------+            +---------------------------+
        ^                                             |
        +--- AT+CEREG? and AT+CSQ on each wake -------+
  duty cycle:  |<-- ~20 B --> . . . . . . . . . . . . . . . . |<-- ~20 B -->
               wake           idle, about 10 minutes           wake
\end{asciiart}

The payload budget is about 20 bytes. That is not a limitation of the radio so much as a discipline: a field node that sends a compact binary record every 10 minutes stays inside the smallest data plan on the market, and one that sends a formatted JSON document with a timestamp string does not. Decide the record layout before the first send, and write it in the lab book next to the endpoint address.

\subsection*{Steps}

\begin{steps}
\item \textbf{Assemble before power.} Seat the HAT on the Pi 3 with the supply unplugged, check that the UART jumper is in the Waveshare B position, fit the NB antenna, and insert the NB-IoT SIM. A SIM inserted under power is not detected until the next boot, and a module powered with no antenna fitted will not attach.

\item \textbf{Free the PL011 and open the port.} Bluetooth owns \texttt{ttyAMA0} on a Pi 3 until the overlay moves it, so apply \texttt{dtoverlay=disable-bt} in \texttt{/boot/firmware/config.txt}, set the serial login shell off and the serial hardware on under \texttt{raspi-config}, and reboot.
\begin{shellcode}
sudo systemctl disable --now hciuart
sudo apt-get install -y minicom python3-serial
sudo minicom -D /dev/ttyAMA0 -b 115200
\end{shellcode}

\item \textbf{Identity and attach.} The band is operator-specific. In the EU it is commonly B8 or B20.
\begin{atcode}
AT
AT+CPIN?
AT+CBAND=20
AT+COPS=0
AT+CEREG?
AT+CGATT?
AT+CSQ
AT+CGDCONT?
\end{atcode}

\item \textbf{Enable power-saving mode.} The encoded tracking-area-update and active-time bitmaps are operator-checked. The values below request a long update interval and a short active window. Confirm them with the SIM sheet before leaving this running unattended.
\begin{atcode}
AT+CPSMS=1,,"00100011","00000101"
AT+CEDRXS=1,5,"0010"
AT+CPSMS?
\end{atcode}

\item \textbf{Open a UDP socket toward a CoAP endpoint you control.} Use \texttt{coap.me} only as a lab echo, and expect it to be unreachable: many NB APNs cannot reach the public Internet. The payload below is four bytes of hexadecimal, eight characters on the wire.

The three arguments of \texttt{AT+CSOC} select an IPv4 datagram socket, and \texttt{AT+CSOCON} gives that socket its destination port and address. Confirm the encoding against the SIMCom AT manual for your firmware revision before you rely on it in a script, because the argument order is a module convention and not a standard.
\begin{atcode}
AT+CSOC=1,2,1
AT+CSOCON=0,5683,"172.16.0.1"
AT+CSOSEND=0,8,"50494E47"
\end{atcode}

\item \textbf{Wrap the send in a schedule.} A 10-minute cron entry or a Python loop is enough. If the operator does not honour power-saving mode from this HAT, power-cycle the radio between sends with \texttt{AT+CFUN=0} and \texttt{AT+CFUN=1} instead of trusting the timers.
\begin{pycode}
import time, serial

ser = serial.Serial("/dev/ttyAMA0", 115200, timeout=2)

def at(cmd, wait=1.0):
    ser.write((cmd + "\r\n").encode())
    time.sleep(wait)
    return ser.read(ser.in_waiting or 1).decode(errors="ignore")

while True:
    print(at("AT+CEREG?"), at("AT+CSQ"))
    at("AT+CSOC=1,2,1")
    at("AT+CSOCON=0,5683,\"172.16.0.1\"")
    print(at("AT+CSOSEND=0,8,\"50494E47\""))
    # if PSM is not granted, uncomment the two lines below
    # at("AT+CFUN=0", 2.0)
    # at("AT+CFUN=1", 8.0)
    time.sleep(600)
\end{pycode}

\item \textbf{Drive the LED from the registration state.} Green on BCM27 follows \texttt{AT+CEREG?} and nothing else. It is not a link light and not a transmit light, because on this radio neither of those is meaningful between wakes.

\item \textbf{Count the duty cycle at the endpoint, not on the Pi.} The product of this lab is a record on the receiving side. Leave the node running for an hour and count the datagrams that arrived.
\begin{shellcode}
# on the CoAP endpoint host
sudo tcpdump -n -i any udp port 5683 -ttt
# six datagrams in an hour is the target, not six hundred
\end{shellcode}
\end{steps}

\subsection*{Acceptance test}

\begin{itemize}
\item \texttt{AT+CSQ} and \texttt{AT+CEREG?} show a cell.
\item One datagram per 10 minutes appears on the server. Count them over an hour: six, not five hundred.
\item A continuous ping is a fail. It is also impossible on this path, because the Pi holds no route.
\item \texttt{AT+CPSMS?} reads back the requested values, or shows what the network granted instead.
\item Without a SIM the lab stops at AT bring-up, and that is still a pass.
\item The green LED is on while \texttt{AT+CEREG?} reports a cell and off when it does not, with no other condition attached to it.
\item Between wakes the Pi holds no route, no socket and no open connection. \texttt{ip route} shows wlan0 or eth0 only.
\end{itemize}

\subsection*{Practices}

NB-IoT is not a phone network. Do not expect \texttt{test.mosquitto.org} to answer, and do not treat an unreachable public host as a hardware fault. Use the APN printed on the SIM sleeve, not the one from the Cat-4 SIM in P01.

\begin{description}
\item[Endpoint first.] Have the CoAP endpoint listening and logging before the first \texttt{AT+CSOSEND}. A datagram that nothing recorded did not happen, and this radio gives you no second chance to look at it.
\item[One instrument per question.] Measure watts in P10, with the PPK2 in the power rail. A figure copied from a datasheet is not a measurement, and this bench has no current shunt in it.
\item[One HAT per host.] The MCC 118, the Explorer700 and the 3.5 inch LCD stay on the shelf while this lab is seated. Power off before unseating anything.
\end{description}

\subsection*{Pitfalls}

\begin{itemize}
\item \textbf{Bluetooth still owns the PL011.} Without \texttt{dtoverlay=disable-bt} the AT port is the mini UART, whose baud rate follows the core clock. The symptom is silence or corrupted characters at 115200, which reads like a dead module.
\item \textbf{The serial login shell is still enabled.} The getty consumes the module's responses. The symptom is an AT port that echoes but never returns \texttt{OK}.
\item \textbf{A 4G antenna on the NB port.} The matching is different. The symptom is a low \texttt{AT+CSQ} and an attach that never completes.
\item \textbf{An overlay that was edited but not applied.} \texttt{dtoverlay=disable-bt} takes effect at boot. The symptom is a port that behaves exactly as it did before the edit, which invites a second and unnecessary edit.
\item \textbf{The wrong APN.} \texttt{AT+CGDCONT?} shows what the module will use. The symptom is a registration that succeeds while \texttt{AT+CSOCON} times out.
\item \textbf{Expecting a public broker.} An NB APN commonly has no route off the operator's network. Point the datagram at an endpoint you control.
\item \textbf{Proving the link with a ping loop.} It defeats power-saving mode, so the module never sleeps and the lab measures nothing.
\item \textbf{Trusting the requested timers.} \texttt{AT+CPSMS} is a request. Read it back, and use the \texttt{AT+CFUN} fallback when the network declines.
\item \textbf{Taking the AT port away from the HAT.} If the UART jumper is not in the B position the Pi is not the controller, and the port answers nothing that the Pi sends.
\item \textbf{Seating a second HAT.} The SIM7020E owns the header. The MCC 118, the Explorer700 and the 3.5 inch LCD each want the same pins, and stacking two of them is how a lab afternoon is lost to a fault that is not a fault.
\item \textbf{Leaving minicom open while the script runs.} Two writers on one serial port interleave their commands and read each other's responses. The symptom is an AT dialogue that looks intermittent rather than broken, which is the hardest kind to diagnose. Close the terminal session first.
\item \textbf{A cold module and an impatient timeout.} The attach on a first power-up takes longer than a warm one. Give \texttt{AT+CEREG?} several attempts before deciding the antenna or the band is at fault.
\item \textbf{Reading the LED as a link light.} Green means registered. It does not mean the last datagram arrived, and on this radio nothing on the Pi can tell you that. The endpoint's log can.
\end{itemize}

\subsection*{Sources}

\begin{itemize}
\item Waveshare NB-IoT HAT wiki for the SIM7020E, jumper positions and antenna ports, \url{https://www.waveshare.com/wiki/}
\item SIMCom SIM7020 series AT command manual, \texttt{AT+CBAND}, \texttt{AT+CPSMS}, \texttt{AT+CEDRXS}, \texttt{AT+CSOC}, \texttt{AT+CSOCON}, \texttt{AT+CSOSEND} (module vendor documentation)
\item The Constrained Application Protocol, RFC 7252, \url{https://www.rfc-editor.org/rfc/rfc7252}
\item Raspberry Pi documentation, UART configuration and the Bluetooth overlay, \url{https://www.raspberrypi.com/documentation/computers/configuration.html}
\item Raspberry Pi hardware documentation, 40-pin header, \url{https://www.raspberrypi.com/documentation/computers/raspberry-pi.html}
\item The operator's own NB-IoT APN and band notes, which arrive with the SIM and outrank every other document listed here
\item P01 in this volume for the Cat-4 contrast, and P10 for the energy measurement this lab deliberately does not attempt
\item P13 in this volume, which reuses this attachment unchanged as the backup signalling path
\end{itemize}

\project{03}{Cat-M MQTT node with GNSS stamp}{Raspberry Pi 3B+ + SIM7070G}{Cat-M and the GNSS-timestamp-in-payload rule}

\begin{unique}
Owns Cat-M and the GNSS-timestamp-in-payload rule. Dual-mode silicon, single-mode lab.
\end{unique}

\begin{blueprint}
Draft P19 drove the SIM7070 from an ESP32 as the application processor. That is valid with flying UART leads, but then it stops being an embedded-Linux lab. We keep the Pi 3B+ as the host and use the module's built-in MQTT and GNSS AT set. Power GNSS down during transmission.
\end{blueprint}

\subsection*{Intent}

Lock the SIM7070G to Cat-M with \texttt{AT+CMNB=1} and keep it there. The module is dual-mode silicon: it will also do NB-IoT and it will fall back to GPRS. This lab is single-mode on purpose, because a node that quietly changes bearer changes its latency, its power profile and its cost without telling anyone, and the log then describes a network that was never configured.

Every publish carries a GNSS time, not \texttt{date} from the Pi. That is the rule this chapter owns. The Pi clock on a Pi 3B+ is whatever the last network time synchronisation left behind, and in the field there may not have been one. The GNSS engine on the same module has a better one, so the payload starts with the stamp from \texttt{+CGNSINF} and the Pi clock is not consulted at all.

\begin{kit}
Pi 3B+, SIM7070G HAT with LTE and GNSS antennas and a Cat-M SIM, yellow LK-LED10 module on BCM16.
\end{kit}

\begin{note}{Without a SIM}
A data SIM is assumed only for P01 to P03 and P13. Without one this lab stops at AT bring-up, and that is still a pass. \texttt{AT+CGNSPWR=1} and \texttt{AT+CGNSINF} still work with no SIM at all, so the GNSS half of the lab can be completed on its own.
\end{note}

\subsection*{System architecture}

\diagram{p03_arch}{Two subsystems inside one module, reached over one AT channel. The MQTT session and the GNSS engine never run at the same moment, and that ordering is the lab.}

The SIM7070G carries its own MQTT client. The session is configured with \texttt{AT+SMCONF}, opened with \texttt{AT+SMCONN} and used with \texttt{AT+SMPUB}, all over the AT character device. Nothing in Python holds a socket, which is the opposite arrangement from a host-side client library. The advantage is that the session, and any transport security you later add to it, stays on the modem where the radio state also lives. The cost is that every publish is a short AT dialogue with return codes to parse, so the host script has to be written as a protocol driver rather than as a call to a client library.

The GNSS engine is the second subsystem, and it is the one that has to be managed. It is powered on with \texttt{AT+CGNSPWR=1}, polled with \texttt{AT+CGNSINF}, and powered off again with \texttt{AT+CGNSPWR=0} before the module transmits. The host cycle is therefore: acquire a fix, power the receiver down, then open the MQTT session and publish the fix that was just cached. USB is preferred for the AT and NMEA traffic, so the character devices appear when the HAT USB cable is connected.

\begin{table}[H]\centering\small
\begin{tabular}{p{34mm}p{112mm}}
\toprule
Command & What it settles \\
\midrule
\texttt{AT+CPIN?} & The SIM is seated and unlocked \\
\texttt{AT+CNMP=38} & LTE only, so the module does not choose a 2G bearer \\
\texttt{AT+CMNB=1} & Cat-M, not NB-IoT. This is the line the lab is named after \\
\texttt{AT+CPSI?} & The bearer actually in use. It must mention CAT-M or eMTC \\
\texttt{AT+CGNSPWR} & Powers the GNSS receiver on, and off again before every transmission \\
\texttt{AT+CGNSINF} & The fix: run status, fix status, UTC time, latitude, longitude \\
\texttt{AT+SMCONF} & Broker address, port and client identifier, held on the module \\
\texttt{AT+SMCONN} & Opens the session from the module, not from Python \\
\texttt{AT+SMPUB} & Publishes one payload to one topic \\
\bottomrule
\end{tabular}
\caption{The AT set this lab uses and what each command settles. Two of them, \texttt{AT+CMNB} and \texttt{AT+CGNSPWR}, carry the whole discipline of the chapter.}
\end{table}

\subsection*{Wiring and schematic}

The seating rules are the same as P01: one 40-pin HAT, jumpers set before power, antennas fitted before power. The only flying wire in the lab is the yellow LK-LED10 module on BCM16.

\diagram{p03_schematic}{The HAT on the header, the USB cable that carries the AT and NMEA channels, and the single status LED. The GNSS antenna needs sky view, which on most benches means a window and a longer lead.}

\begin{table}[H]\centering\small
\begin{tabular}{p{40mm}p{24mm}p{22mm}p{62mm}}
\toprule
Signal & Pi header pin & BCM GPIO & Note \\
\midrule
HAT 5 V & 2 and 4 & & Supplied through the header when the HAT is seated \\
HAT GND & 6 & GND & \\
Yellow module, S1 & 36 & GPIO16 & Cat-M confirmed and the session open. Lit when driven high \\
Yellow module, G & 39 & GND & The return. The module's resistor is inside it \\
PWRKEY & 7 & GPIO4 & The HAT's own. Never drive it from anything else \\
Pin 15 & & GPIO22 & The draft's choice. Free on this HAT, see the note \\
VCCIO jumper & & & 3V3 \\
PWR jumper & & & Auto-on, as in P01 \\
LTE antenna & & & Main port on the HAT \\
GNSS antenna & & & Separate port, needs sky view \\
HAT USB cable & & & To a Pi USB-A port: AT and NMEA \\
\bottomrule
\end{tabular}
\caption{Wiring. Confirm the jumper block legend on the silkscreen before power, since the HAT revision decides which position is which.}
\end{table}

\begin{note}{What this HAT claims, and the three ways its vendor spells one pin}
The draft put the indicator on GPIO22. On the SIM7600E-H of P01 the draft's indicator pins turned out to be modem control lines, GPIO27, GPIO22 and GPIO23 for ring indicator, data terminal ready and clear to send. The SIM7020E of P02 was at first recorded as claiming GPIO27 for its power key; the vendor's wiki puts that key on GPIO4 like this board's, and that record is withdrawn in P02. This board was read next, and it and the SIM7020E are the two that leave the draft's pin alone.

The module's own hardware design document settles only half the question. It confirms the SIM7070G has that family of lines to route, PWRKEY on module pin 1, UART1\_DTR on 3, UART1\_RI on 4, UART1\_DCD on 5, UART1\_CTS on 7 and UART1\_RTS on 8, plus STATUS on 66 and NETLIGHT on 52. Which Raspberry Pi header pin each reaches is the carrier board's decision, so the HAT's own pages are what answer it.

They do, and they claim remarkably little: the UART on GPIO14 and GPIO15, the module's power key, and one optional line.

\begin{table}[H]\centering\small
\begin{tabular}{p{34mm}p{26mm}p{20mm}p{56mm}}
\toprule
HAT connection & Header pin & BCM GPIO & Note \\
\midrule
PWRKEY & 7 & GPIO4 & Default. The PWR jumper's other position takes it to the supply instead \\
DTR & 37 & GPIO26 & \textbf{Only} with the DTR jumper in position A. The default is disconnected \\
IO level select & none & none & Default 3.3 V. The other position makes the HAT's logic 5 V \\
\bottomrule
\end{tabular}
\caption{What the SIM7070G HAT puts on the header besides the UART, from the vendor's product page and wiki.}
\end{table}

\textbf{So GPIO22 is free on this board and the draft's pin would have worked.} It stays moved anyway, for a reason worth stating: one indicator assignment that is safe on all four cellular labs is easier to hold in the head than four, and BCM16 is free on every HAT checked so far. What is not free is GPIO26, which an intermediate version of this correction had used and which is exactly what the DTR jumper connects when it is in position A. That pin was abandoned for an unrelated reason, the MCC 118 claiming it in P05, and the right answer arrived by luck rather than by reading. The reading is what settles it now.

\textbf{One pin, three spellings, one vendor.} The product page says PWRKEY goes to "P4". The wiki says it goes to "P7 (wiringPi number)". Those are the same conductor: wiringPi 7 is BCM GPIO4 is physical header pin 7. Neither page says which scheme it is using, and a reader who assumes physical numbering on the first page wires the power key to 5 V. Every pin in this book is given as header pin and BCM number together for exactly this reason.

\textbf{Never move the IO level jumper.} It defaults to 3.3 V and its other position makes the HAT's logic 5 V, on a header whose every pin is a 3.3 V input. That is the single connection P15 exists to prevent, available on this board as a jumper.
\end{note}

\begin{asciiart}
Same seating rules as P01. USB preferred for AT + NMEA.
VCCIO = 3V3, auto PWR.
GNSS antenna needs sky view.
No second HAT.
LED YEL BCM16 pin 36 = Cat-M confirmed, module MQTT session open.
\end{asciiart}

\begin{rulebox}{Before the supply goes in}
One 40-pin HAT per host: the SIM7070G cannot serve P02, and P02's SIM7020E cannot serve this lab. Fit both antennas before power. Pi I/O is 3.3 V, so the LED branch goes to a header pin and never to a 5 V rail.
\end{rulebox}

\subsection*{Bench layout}

\diagram{p03_bench}{Bench layout. The GNSS antenna is the part that decides whether the lab finishes today: indoors, under a concrete ceiling, a cold start may never complete.}

Put the GNSS antenna where it can see sky, and accept that a first fix can take minutes. This is the reason the software caches the last good fix rather than blocking on the receiver. Fit the LTE antenna too, because a module that transmits into an open port is stressing its own power amplifier for no result.

\subsection*{Software design (UML)}

\diagram{p03_uml}{The publish cycle. Note the ordering: the GNSS receiver is powered down before the module opens its MQTT session, and the fix in the payload is the one cached a moment earlier.}

The sequence is the product of this lab. Read it from the top: lock the bearer, confirm it, power the GNSS engine, poll until \texttt{+CGNSINF} carries usable fields, power the engine down, then open the session and publish. The host script never holds both subsystems active, and it never blocks the publish on a cold start. If no fix has been acquired yet, the cached value from the previous cycle is used and marked as such, and if there is no cached value the cycle is skipped rather than filled with the Pi clock.

\subsection*{Data flow (ASCII)}

\begin{asciiart}
  Pi 3B+                                      SIM7070G HAT
  +---------------------------------+         +-----------------------------+
  | p03_fix.py                      |         | GNSS engine                 |
  |   AT+CNMP=38 ; AT+CMNB=1        |         |   AT+CGNSPWR=1  (on)        |
  |   AT+CPSI?      -> CAT-M ?      |  USB    |   AT+CGNSINF    -> fix      |
  |   AT+CGNSPWR=1                  |<=======>|   AT+CGNSPWR=0  (off)       |
  |   parse +CGNSINF -> cache fix   |  ttyUSB2|                             |
  |   AT+CGNSPWR=0     <-- ALWAYS   |         | Cat-M radio + MQTT client   |
  |   AT+SMCONF / AT+SMCONN         |         |   AT+SMCONN                 |
  |   AT+SMPUB "lab/p03/fix",0,1    |         |   AT+SMPUB ---------------------> broker
  |   payload: <GNSS stamp>,lat,lon |         |                             |    1883
  |   sleep 60 s                    |         +-----------------------------+
  | gpio: YEL BCM16 = Cat-M + up    |
  +---------------------------------+
  ordering:  GNSS on -> fix -> GNSS OFF -> transmit.  Never GNSS on during TX.
\end{asciiart}

The cycle is one minute long and it has exactly one slow part, the fix. Everything else is a short AT exchange. When a cycle produces no new fix, the cached one is published again, which is honest as long as the stamp inside the payload is the stamp the receiver reported, because a repeated stamp is visibly a repeat. That property is lost the moment a host clock is substituted, which is the whole argument for this lab's rule.

\subsection*{Steps}

\begin{steps}
\item \textbf{Assemble before power.} Seat the HAT on the Pi 3B+ with the supply unplugged, set VCCIO to 3V3 and leave PWR on auto-on, fit the LTE and GNSS antennas, insert the Cat-M SIM, and connect the HAT USB cable to a Pi USB-A port. No second HAT.

\item \textbf{AT bring-up on the USB AT port.} Lock the bearer before anything else, and confirm it with \texttt{AT+CPSI?} rather than assuming.
\begin{atcode}
sudo minicom -D /dev/ttyUSB2 -b 115200
AT+CPIN?
AT+CNMP=38          # LTE only
AT+CMNB=1           # Cat-M, NOT NB-IoT
AT+CPSI?            # must mention CAT-M / eMTC
AT+CSQ
\end{atcode}

\item \textbf{GNSS with radio courtesy.} Power the receiver, wait outdoors, read the fix, and power the receiver down again before any transmission.
\begin{atcode}
AT+CGNSPWR=1
# wait, outdoors
AT+CGNSINF
# before MQTT TX:
AT+CGNSPWR=0
\end{atcode}

\item \textbf{Module MQTT.} This keeps the session, and any transport security you add later, on the modem rather than in Python.
\begin{atcode}
AT+SMCONF="URL","test.mosquitto.org",1883
AT+SMCONF="CLIENTID","p03-catm"
AT+SMCONN
AT+SMPUB="lab/p03/fix",0,1
# type payload:  2026-09-14T10:00:00Z,lat,lon
\end{atcode}

\item \textbf{Automate.} A 30-line Python script talks AT over \texttt{serial.Serial}, parses \texttt{+CGNSINF}, powers GNSS off, publishes, and sleeps 60 s.
\begin{pycode}
import time, serial

ser = serial.Serial("/dev/ttyUSB2", 115200, timeout=2)
last_fix = None

def at(cmd, wait=1.0):
    ser.write((cmd + "\r\n").encode())
    time.sleep(wait)
    return ser.read(ser.in_waiting or 1).decode(errors="ignore")

at("AT+CNMP=38"); at("AT+CMNB=1")
print(at("AT+CPSI?"))          # must mention CAT-M / eMTC

while True:
    at("AT+CGNSPWR=1")
    inf = at("AT+CGNSINF", 2.0)
    for line in inf.splitlines():
        if line.startswith("+CGNSINF:"):
            f = line.split(":", 1)[1].split(",")
            if len(f) > 5 and f[1].strip() == "1" and f[2].strip():
                last_fix = (f[2].strip(), f[3].strip(), f[4].strip())
    at("AT+CGNSPWR=0")         # power down BEFORE the transmission
    if last_fix:
        at("AT+SMCONF=\"URL\",\"test.mosquitto.org\",1883")
        at("AT+SMCONF=\"CLIENTID\",\"p03-catm\"")
        at("AT+SMCONN", 4.0)
        at("AT+SMPUB=\"lab/p03/fix\",0,1")
        ser.write((",".join(last_fix) + "\x1a").encode())
    time.sleep(60)
\end{pycode}

\item \textbf{Keep the evidence.} Append the \texttt{+CPSI} line and the published payload to a log on every cycle. The pair is what proves the lab afterwards: the bearer that was in use, and the stamp that went out with it.
\begin{shellcode}
# one line per cycle, appended by the script
# 2026-09-14T10:00:00Z  CPSI=LTE CAT-M1,...  lab/p03/fix  published
tail -f /home/pi/p03.log
\end{shellcode}

\item \textbf{Drive the LED from two conditions, not one.} Yellow on BCM16 is on when \texttt{AT+CPSI?} last contained CAT-M and the module reported the session open. If the bearer changed, the LED goes out, which is the fastest way to notice a silent fallback from across the bench.

The module has a status LED of its own and it answers a different question, which is why this lab fits a second one rather than reading the first. NETLIGHT on module pin 52 drives a network indicator with four states defined in the hardware design document: 64 ms on and 800 ms off while no network is registered, 64 ms on and 3000 ms off once registration succeeds, 64 ms on and 300 ms off while data is moving, and off for power down or power saving. Useful, and silent on the one thing this chapter exists for: all three lit patterns look identical whether the bearer is Cat-M, NB-IoT or a GPRS fallback. The host LED reports the bearer and the session; the module's reports that a network exists. Watch both and the pair tells you which half failed.
\end{steps}

\subsection*{Acceptance test}

\begin{itemize}
\item \texttt{AT+CPSI?} contains CAT-M.
\item The broker payload starts with a GNSS stamp, not the Pi clock. Compare the two deliberately once: set the Pi clock wrong and confirm the payload does not move.
\item A module that silently falls to GPRS fails the lab. Fix \texttt{AT+CMNB} before continuing.
\item \texttt{AT+CGNSPWR} is 0 during every transmission, and the transcript shows it.
\item Without a SIM the lab stops at AT bring-up, and that is still a pass.
\item One publish per minute arrives at the topic, and the yellow LED is out whenever the last \texttt{+CPSI} line did not say CAT-M.
\item The GNSS half of the lab passes on its own: \texttt{AT+CGNSINF} returns non-blank fields with the antenna at a window, with or without a SIM.
\end{itemize}

\subsection*{Practices}

Cache the last good fix. Do not block MQTT on a cold GNSS start. The same HAT cannot serve P02, so plan the sitting: this lab and the NB-IoT lab are different modules on different hosts, and the source's build order puts P01, then this lab, then P02, then P13.

\begin{description}
\item[Confirm, do not assume, the bearer.] \texttt{AT+CMNB=1} is a request to the module, and \texttt{AT+CPSI?} is the evidence. Log the \texttt{+CPSI} line with every publish cycle so a fallback is visible in the record and not just on the LED.
\item[Stamp at the source.] The payload carries the time reported by the GNSS receiver. A stamp added by the broker, or by the Pi, describes when the message was handled, which is a different quantity.
\item[One radio question per lab.] Cat-4 routing is P01, NB-IoT duty cycling is P02, and failover across two of them is P13. Do not blend them here.
\item[Give the cold start its own session.] Prove \texttt{AT+CGNSINF} by hand at a window before wiring it into the loop. A receiver that has never had a fix and a script that has never had a fix are two separate problems, and finding out which one you have costs an afternoon if they are debugged together.
\end{description}

\subsection*{Pitfalls}

\begin{itemize}
\item \textbf{A silent fall to GPRS.} The module attaches, the publish succeeds and nothing looks wrong. The symptom is in \texttt{AT+CPSI?}, which stops mentioning CAT-M. Check it on every cycle, not once at bring-up.
\item \textbf{GNSS still powered during transmission.} The source's rule exists because the two subsystems share the module. Power the receiver down first, every time, and put the command in the script rather than in the operator's memory.
\item \textbf{Blocking on a cold start.} A first fix can take minutes and may never arrive indoors. A publisher that waits for it sends nothing at all, which is worse than publishing a cached fix marked as cached.
\item \textbf{The Pi clock in the payload.} It is the failure this lab is built to prevent. If the script ever calls \texttt{date} or \texttt{time.time()} for the stamp, the lab is not passed.
\item \textbf{The GNSS antenna indoors.} Under a concrete ceiling there is no fix to acquire. Move the antenna to a window before suspecting the module.
\item \textbf{Reading a fix status of zero as a coordinate.} \texttt{+CGNSINF} returns its fields whether or not the fix is valid. Check the fix-status field before caching anything from that line.
\item \textbf{Picking the wrong AT device.} Several \texttt{ttyUSB} nodes appear on the HAT USB cable. The wrong one is silent rather than an error, so try the next one before suspecting the module.
\item \textbf{Two HATs in one sitting.} The SIM7070G and the SIM7020E of P02 both want the 40-pin header. Finish one lab, power off, unseat, then start the other.
\item \textbf{Leaving minicom open while the script runs.} Two writers on one serial port interleave their commands and read each other's responses.
\item \textbf{The payload terminator.} \texttt{AT+SMPUB} asks for the payload after the command line, and the module decides how that input ends. Confirm the terminator in the AT manual for your firmware revision: a publish that never completes usually means the module is still waiting for it.
\item \textbf{A session left open across a bearer change.} \texttt{AT+SMCONN} succeeded once, the radio then changed state, and the next \texttt{AT+SMPUB} returns an error that reads like a broker problem. Re-read \texttt{AT+CPSI?} before blaming the broker host.
\item \textbf{Expecting the public broker to be reachable.} It usually is on a Cat-M APN, unlike the NB APN in P02, but it is a lab convenience and not a product. Point at your own broker, for example the one built in P12, as soon as the path is proven.
\end{itemize}

\subsection*{Sources}

\begin{itemize}
\item Waveshare SIM7070G Cat-M/NB-IoT/GPRS HAT wiki, jumper positions and antenna ports, \url{https://www.waveshare.com/wiki/}
\item SIMCom SIM7070 series AT command manual, \texttt{AT+CNMP}, \texttt{AT+CMNB}, \texttt{AT+CPSI}, \texttt{AT+CGNSPWR}, \texttt{AT+CGNSINF}, \texttt{AT+SMCONF}, \texttt{AT+SMCONN}, \texttt{AT+SMPUB} (module vendor documentation)
\item SIMCom application note on the module MQTT AT set, for the return codes the driver must parse
\item Eclipse Mosquitto documentation, \url{https://mosquitto.org/documentation/}
\item Raspberry Pi hardware documentation, 40-pin header, \url{https://www.raspberrypi.com/documentation/computers/raspberry-pi.html}
\item P01 in this volume for the seating rules this lab reuses, and P12 for the broker that should replace the public one
\item P02 in this volume, which owns NB-IoT: the second mode of this module is deliberately not exercised here
\item The operator's own Cat-M coverage and APN notes, which arrive with the SIM
\end{itemize}

\project{13}{Dual-radio signalling failover}{Pi 4 + SIM7600 and Pi 3 + SIM7020, with the P12 broker}{multi-homing with a source tag}

\begin{unique}
Owns multi-homing. Two Pis, two modems, one MQTT topic with a source tag. Not link bonding, and not a stacked-HAT fantasy.
\end{unique}

\subsection*{Intent}

The Pi 4 with the SIM7600 is the primary path, reusing the P01 WAN exactly as it was built there. The Pi 3 with the SIM7020 is the backup signalling path, reusing P02 exactly as it was built there. A supervisor on the P12 broker watches one topic and marks the system PRIMARY or FAILOVER. Nothing is bonded, nothing is load balanced, and no kernel is asked to hold two default routes at once.

The word signalling is doing real work in that sentence. The backup path carries a health sample every two minutes, roughly twenty bytes at a time. It cannot carry the traffic of the primary path and it is not meant to. What it can do is tell an operator that the node is alive when the Cat-4 route has gone, which is the property this chapter is about.

\begin{kit}
Pi 4 with the SIM7600 attachment from P01, Pi 3 with the SIM7020 attachment from P02, and the P12 broker if it is built, otherwise \texttt{test.mosquitto.org}.
\end{kit}

\begin{note}{This lab has no wiring of its own}
Every wire in this chapter was made in P01 and P02. The schematic below is a topology drawing, not a new attachment: for pin numbers, jumper positions and antenna ports, use the wiring tables in those two chapters. If either attachment is not yet passing its own acceptance test, finish that lab first.
\end{note}

\subsection*{System architecture}

\diagram{p13_arch}{Two complete hosts, two radio paths, one topic. The supervisor is the only component that knows about both, and it knows about them only as tagged samples that arrive or fail to arrive.}

Each host is independent all the way down. The Pi 4 holds a default route on \texttt{usb0} or \texttt{wwan0} and publishes with a normal MQTT client over it. The Pi 3 holds no route at all: its sample leaves through a UDP socket inside the SIM7020E module, as built in P02, and reaches the same topic by whatever path P02 actually reached. Neither host knows that the other exists.

The supervisor holds the entire state of the system, and that state has one input: the arrival time of the most recent sample tagged \texttt{cat4}. This is the smallest arrangement that does the job, and its smallness is the point. A supervisor that also probes the radios, or reads modem status, or tries to bring an interface up, has made itself a second failure surface on top of the one it was meant to observe.

\begin{table}[H]\centering\small
\begin{tabular}{p{38mm}p{52mm}p{52mm}}
\toprule
 & Primary path & Backup path \\
\midrule
Host & Pi 4 & Pi 3 \\
Radio & SIM7600E-H, LTE Cat-4 & SIM7020E, NB-IoT \\
Built in & P01 & P02 \\
Route on the host & Default route on usb0 or wwan0 & None. The socket is inside the module \\
Transport used here & MQTT client over the WAN & UDP datagram from the module \\
Sample interval & 15 s & 120 s \\
Source tag & \texttt{cat4} & \texttt{nb} \\
What its loss means & The supervisor declares FAILOVER after 45 s & The node is unreachable by both paths \\
\bottomrule
\end{tabular}
\caption{The two paths, side by side. They are not equivalent and are not meant to be: one carries a route, the other carries a signal.}
\end{table}

\subsection*{Wiring and schematic}

The wiring is P01's and P02's. Nothing on either host changes for this lab, including the jumper positions, the antennas and the LEDs. What follows is the topology.

\diagram{p13_schematic}{The two-host topology. Two radios, two operator paths, one topic, one supervisor. The two hosts share no cable and no kernel, and the only thing they have in common is the string in the source tag.}

\begin{table}[H]\centering\small
\begin{tabular}{p{32mm}p{34mm}p{30mm}p{52mm}}
\toprule
Host & Attachment & Built in & Path used here \\
\midrule
Pi 4 & SIM7600E-H HAT, LTE and GNSS antennas, ADXL345, LEDs on BCM 16/20/21 & P01, unchanged & Default route on usb0 or wwan0, a normal MQTT client \\
Pi 3 & SIM7020E HAT, NB antenna, green LED on BCM16 & P02, unchanged & UDP socket inside the module, no route on the host \\
Pi 4 with the DSI panel & Official touchscreen and keyboard, Mosquitto & P12, unchanged & Broker and supervisor, on the LAN \\
\bottomrule
\end{tabular}
\caption{The three attachments this lab borrows. It adds no wiring of its own, and it changes nothing on any of them.}
\end{table}

\begin{asciiart}
No new wiring in this lab.

  Pi 4  + SIM7600E-H   ......  exactly as built in P01 (HAT, 2 antennas, USB cable)
  Pi 3  + SIM7020E     ......  exactly as built in P02 (HAT, NB antenna, UART jumper B)
  Pi 4  + DSI panel    ......  exactly as built in P12 (broker, keyboard, glass)

  The only physical action in this chapter is on an antenna connector:
    unscrew the LTE main antenna  -> the state must flip to FAILOVER
    refit the LTE main antenna    -> the state must return to PRIMARY
  Refit it before the next transmission. No antenna means no attach, and
  transmitting into an open port stresses the power amplifier.
\end{asciiart}

\begin{rulebox}{Before the antenna comes off}
Two hosts, two HATs, one per host. Do not stack the Explorer700 or a second modem on either Pi. Power off before unseating anything, and refit the LTE antenna as soon as the failover test is recorded.
\end{rulebox}

\subsection*{Bench layout}

\diagram{p13_bench}{Bench layout. Three boards, three supplies and one pair of hands on one antenna connector. The glass shows the supervisor state, so the operator can watch the transition while performing it.}

Put the two radio hosts far enough apart that their antennas are not touching, and put the operator glass where you can see it from the antenna. The transition takes up to 45 seconds, so the test is a slow one: unscrew, wait, watch the state change, refit, wait, watch it return. Write both times in the lab book.

Three supplies, three power switches and one topic. Label the two radio hosts on their cases before you start, because from across the bench a Pi 3 with a black HAT and a Pi 4 with a black HAT look identical, and this is the one lab where pulling the wrong antenna proves nothing.

\subsection*{Software design (UML)}

\diagram{p13_uml}{The supervisor. Two states, one timer and one input. The 45 second window is three missed samples from a publisher that sends every 15 seconds, which is long enough to survive one lost message and short enough to notice a lost route.}

The supervisor is a subscriber with a clock. Every message on \texttt{lab/p13/health} carries a source tag. When the tag is \texttt{cat4} the arrival time is recorded. When the tag is \texttt{nb} the sample is kept as the latest backup reading. On every tick, if the most recent \texttt{cat4} sample is older than 45 seconds the state becomes FAILOVER and the \texttt{nb} sample is the live one, and when a \texttt{cat4} sample arrives again the state returns to PRIMARY.

\subsection*{Data flow (ASCII)}

\begin{asciiart}
  Pi 4 + SIM7600 (P01)                          Pi 3 + SIM7020 (P02)
  +-------------------------+                   +-------------------------+
  | default route on usb0   |                   | no route on the host    |
  | mosquitto_pub every 15 s|                   | AT+CSOSEND every 120 s  |
  |   {"src":"cat4","ok":1} |                   |   {"src":"nb","ok":1}   |
  +-----------+-------------+                   +-----------+-------------+
              |  LTE Cat-4                                  |  NB-IoT
              v                                             v
      [ operator network ]                          [ operator network ]
              |                                             |
              +---------------------+-----------------------+
                                    v
                        topic  lab/p13/health
                                    |
                                    v
                    +-------------------------------+
                    | supervisor on the P12 broker  |
                    |   last cat4 sample < 45 s ?   |
                    |     yes -> PRIMARY            |
                    |     no  -> FAILOVER, nb live  |
                    +-------------------------------+
\end{asciiart}

Read the diagram from the bottom. The supervisor sees one stream of tagged records and nothing else. It does not see two interfaces, it cannot see a route, and it has no opinion about radios. That is what makes the test at the end of this chapter meaningful: when the antenna comes off, the only thing that changes in the supervisor's world is that records with one tag stop arriving.

\subsection*{Steps}

\begin{steps}
\item \textbf{Do not undress P01 and P02 on the same afternoon} if you still need their individual acceptance tests. This lab assumes both already attach. Confirm that first: \texttt{AT+CEREG?} on each host, and a default route on \texttt{usb0} on the Pi 4.

\item \textbf{Publish the primary health sample from the Pi 4}, every 15 s, through the P01 WAN.
\begin{pycode}
import json, subprocess, time

BROKER = "localhost"        # the P12 broker, or test.mosquitto.org
while True:
    payload = json.dumps({"src": "cat4", "ok": 1})
    subprocess.run(["mosquitto_pub", "-h", BROKER,
                    "-t", "lab/p13/health", "-m", payload], check=False)
    time.sleep(15)
\end{pycode}

\item \textbf{Publish the backup health sample from the Pi 3}, every 120 s, over CoAP/UDP or whatever P02 actually reached. The payload is the same record with a different tag, and it leaves through the module socket, not through a host route.
\begin{atcode}
AT+CSOC=1,2,1
AT+CSOCON=0,5683,"172.16.0.1"
# the record {"src":"nb","ok":1} as hex, length counted as in P02
AT+CSOSEND=0,38,"7B22737263223A226E62222C226F6B223A317D"
\end{atcode}

\item \textbf{Run the supervisor} on the P12 broker host, or on a laptop on the same LAN.
\begin{pycode}
import json, time
import paho.mqtt.client as mqtt

last_cat4 = 0.0
last_nb = None
state = "FAILOVER"          # nothing has been heard yet

def on_message(client, userdata, msg):
    global last_cat4, last_nb, state
    rec = json.loads(msg.payload.decode(errors="ignore"))
    if rec.get("src") == "cat4":
        last_cat4 = time.time()
        if state != "PRIMARY":
            state = "PRIMARY"
            print(time.strftime("%H:%M:%S"), "state ->", state)
    elif rec.get("src") == "nb":
        last_nb = rec

c = mqtt.Client()
c.on_message = on_message
c.connect("localhost", 1883, 60)
c.subscribe("lab/p13/health")
c.loop_start()

while True:
    if state == "PRIMARY" and time.time() - last_cat4 > 45:
        state = "FAILOVER"
        print(time.strftime("%H:%M:%S"), "state ->", state, "live:", last_nb)
    time.sleep(1)
\end{pycode}

\item \textbf{Unscrew the LTE main antenna.} The state must flip to FAILOVER within 45 seconds of the last primary sample. Refit the antenna. The state must return to PRIMARY on the next primary sample. Record both transitions with their timestamps.

\item \textbf{Record both sources to a file.} The log is the acceptance artifact, and it has to show both tags and the two transitions with their timestamps.
\begin{shellcode}
mosquitto_sub -h localhost -t 'lab/p13/health' -v | ts >> /var/log/lab-p13.txt
# then, after the antenna test:
grep -c '"src":"cat4"' /var/log/lab-p13.txt
grep -c '"src":"nb"'   /var/log/lab-p13.txt
ip route | grep default        # on each host: one modem route, not two
\end{shellcode}

\item \textbf{Show the state on the glass.} A \texttt{watch} terminal on the P12 panel, or the dashboard page from that chapter, is enough. The operator should be able to read the state from across the bench without a keyboard.

\item \textbf{Refit the antenna and confirm the return.} The test is not finished at the flip. Screw the LTE main antenna back on, wait for the next primary sample, and confirm that the state returns to PRIMARY on its own, with no restart of the supervisor and no intervention on either host.
\end{steps}

\subsection*{Acceptance test}

\begin{itemize}
\item The antenna-pull test above: unscrew the LTE main antenna, the state flips to FAILOVER; refit it, the state returns to PRIMARY.
\item Logs show both sources. A \texttt{cat4} record roughly every 15 s and an \texttt{nb} record roughly every 120 s, each with its tag intact.
\item Never a bonded \texttt{metric 0} default route on one kernel. Check \texttt{ip route} on both hosts: the Pi 4 has one default route through the modem, and the Pi 3 has none through its modem at all.
\item The transition is visible on the P12 glass without opening a terminal.
\item The FAILOVER transition happens no later than 45 seconds after the last \texttt{cat4} record, and the return to PRIMARY happens on the first \texttt{cat4} record after the antenna is refitted.
\item The backup publisher keeps its own rhythm through both transitions. Its interval does not change because the state changed, because the two are not connected.
\item The LTE antenna is back on the connector at the end of the sitting, and the lab book says so.
\end{itemize}

\subsection*{Practices}

This is signalling failover. Do not advertise a 50 Mbit NB-IoT backup, it does not exist. The backup path carries a health record and nothing else, and any design that assumes otherwise will be found out on the first day it is needed.

\begin{description}
\item[Tag at the source, not at the broker.] The string \texttt{cat4} or \texttt{nb} is written by the host that owns the radio. A tag added by the supervisor describes what the supervisor believed, which is the thing under test.
\item[One timer, one state.] The 45 second window is the only tunable in the supervisor. Resist adding a second timer for the backup path until the first one has been observed over a working day.
\item[Keep the two publishers dumb.] Neither host should try to detect the other's state or change its own behaviour when the supervisor changes state. All the intelligence in this design sits in one place, and that is why it can be read in twenty lines.
\item[Refit the antenna.] The teardown rule from the source applies here more than anywhere else, because this lab asks you to remove an antenna on purpose. Refit it before the next transmission.
\end{description}

\subsection*{Pitfalls}

\begin{itemize}
\item \textbf{Treating this as link bonding.} Two radios, two kernels, two independent paths. There is no interface that carries both, and no kernel is asked to choose between them.
\item \textbf{A supervisor that probes.} If the supervisor pings the Pi 4, it is no longer measuring the uplink the publisher uses. Arrival of a tagged sample is the only evidence this design accepts.
\item \textbf{A timeout shorter than the publish interval.} With a 15 second publisher, a 20 second window reports FAILOVER every time one message is lost. The 45 second window tolerates two.
\item \textbf{A backup interval longer than the timeout, read as a fault.} The NB path sends every 120 seconds by design. A FAILOVER state may therefore be minutes old before the backup sample refreshes, and that is the honest behaviour of this path, not a defect.
\item \textbf{One SIM in both hosts, swapped between tests.} The two paths have to be live at the same time for the test to mean anything. Two SIMs, two radios, two hosts.
\item \textbf{Both scripts publishing from the same host.} It proves nothing. The test needs two hosts, because the failure being modelled is the loss of one host's radio.
\item \textbf{Undressing P01 or P02 first.} If either attachment is broken down, the lab starts with two unproven radios and one unproven supervisor, and the first failure will be ambiguous.
\item \textbf{Leaving the antenna off.} Transmitting into an open port stresses the power amplifier. Refit it as soon as the transition is recorded.
\item \textbf{Publishing both tags with one client identifier.} If both hosts reach the broker with the same client identifier, the broker will keep dropping one of the sessions and the log will look like a radio problem.
\item \textbf{A supervisor clock that is not the broker's.} If the supervisor runs on a laptop that sleeps, the 45 second window measures the laptop's suspend behaviour rather than the radio. Run it on the P12 host.
\item \textbf{Declaring success on the first flip.} The interesting half of the test is the return to PRIMARY. A supervisor with a latch, or with an uninitialised timer, will flip once and stay flipped.
\end{itemize}

\subsection*{Sources}

\begin{itemize}
\item P01 in this volume, for the Cat-4 attachment, the WAN and the LED conventions this lab reuses unchanged
\item P02 in this volume, for the NB-IoT attachment and the module socket this lab reuses unchanged
\item P12 in this volume, for the broker, the panel and the operator glass the supervisor runs on
\item Eclipse Mosquitto documentation, broker configuration and \texttt{mosquitto\_pub}, \url{https://mosquitto.org/documentation/}
\item Eclipse Paho Python client documentation, \url{https://eclipse.dev/paho/}
\item SIMCom SIM7600 and SIM7020 AT command manuals, for \texttt{AT+CEREG?} and \texttt{AT+CSOSEND} (module vendor documentation)
\item The lab sequence at the end of this volume, which places this chapter after P01, P03 and P02 for exactly this reason
\item The operator's own coverage notes for both SIMs, since a backup path that shares one cell with the primary is not a backup
\end{itemize}


\part{Instruments: vibration, analog, ranging, energy}
\project{04}{STWIN.box vibration and ultrasound USB gateway}{Raspberry Pi 4 + STEVAL-STWINBX1 + official DSI touchscreen}{6 kHz vibration and the ultrasonic microphones}

\begin{unique}
Owns the only 6 kHz vibration sensor, the IIS3DWB, and the only ultrasound and industrial MEMS microphones. USB-CDC into Linux. The DSI glass is the waterfall. No HAT.
\end{unique}

\begin{blueprint}
Draft P1 called USB-C a "VCP" and invoked Intel MKL. USB-C is the connector; CDC-ACM exists only after HSDatalog or USB-CDC firmware is present. Use a NumPy FFT. Draft P20 reused the STWIN as a BLE factory tier, which would double-own this silicon. Anomaly-only uplink is a policy in this lab, not a second project.
\end{blueprint}

\subsection*{Intent}

The STWIN.box is a USB industrial probe. The Pi 4 isolates one CPU core for the FFT thread with \texttt{isolcpus} and \texttt{taskset}, and classifies NORMAL, WARNING or FAULT from spectral peaks. The DSI panel carries the waterfall, and the 40-pin header stays empty, which is what makes this lab compatible with the panel in the first place.

The demonstration is an orthogonality test, and it is worth stating before the wiring. Touch the case and the vibration energy rises. Speak near the digital microphone and the microphone RMS rises while the vibration band does not move. Two physical quantities, two independent paths through the same cable, and one host that can tell them apart. A system that cannot separate them is not measuring what its labels claim.

\begin{kit}
Pi 4, official DSI touchscreen, STWIN.box, USB-C cable. The optional STWIN LiPo only if it already sits in the kit box.
\end{kit}

\begin{note}{The firmware decides whether this lab exists}
CDC-ACM appears only after HSDatalog or a USB-CDC firmware is present on the probe. Until then the USB-C connector carries power and a DFU interface and nothing else. Flash first, then look for \texttt{/dev/ttyACM0}, and read the frame format from ST's tools rather than guessing at it.
\end{note}

\subsection*{System architecture}

\diagram{p04_arch}{The probe and the host. The microcontroller moves samples and stamps them. Linux does every arithmetic operation that matters, on a core that the scheduler has been told to leave alone.}

The boundary is the design. On the STWIN side, the STM32U585 reads the IIS3DWB over SPI and the microphones over their own interfaces, moves the samples with DMA, applies a timestamp and a sequence number, and packetises the result onto the USB CDC endpoint. It does not transform the data. On the Linux side, a reader pulls frames from \texttt{/dev/ttyACM0} into a ring buffer, a windowed FFT runs on the isolated core, and a policy decides NORMAL, WARNING or FAULT from the resulting peaks.

Note what is missing from that list. There is no analysis on the probe, no storage on the probe, and no decision on the probe. The device is a well-behaved source of stamped samples, and everything downstream of the cable can be rewritten, restarted and rerun against a recorded file without touching the firmware at all.

That split is not a stylistic preference. The Cortex-M33 in the probe has better access to the sensors and worse access to everything else: no file system worth the name, no display, no package manager and no way to change the analysis without a reflash. Linux owns the system. The microcontroller is a peripheral that happens to contain a processor, and the moment an FFT moves into it, the parts of the lab you most want to change become the parts that are hardest to change. Do not FFT on the M33.

\begin{table}[H]\centering\small
\begin{tabular}{p{32mm}p{46mm}p{68mm}}
\toprule
Part & What it measures & Why it is here \\
\midrule
IIS3DWB & Vibration to 6 kHz & The only 6 kHz vibration sensor in the kit. The ADXL345 of P01 and P19 cannot do this band \\
ISM330DHCX & Inertial, with a machine learning core & On the probe, not on the IKS5A1. The draft placed it on the wrong board \\
IMP34DT05 & Digital microphone & The acoustic half of the orthogonality test \\
IMP23ABSU & Analog microphone, ultrasonic range & The only ultrasonic path in the kit \\
STM32U585 & Nothing. It moves samples & DMA, timestamp, sequence number, packetise. No arithmetic on the signal \\
\bottomrule
\end{tabular}
\caption{What is on the probe. One physics per lab: these are not the IKS4A1 sensors of P07 and not the ADXL345 of P19.}
\end{table}

\subsection*{Wiring and schematic}

\diagram{p04_schematic}{The probe's internal signal paths and the two cables that leave it. Nothing here touches the 40-pin header, which is what lets the DSI panel and this probe coexist on one Pi.}

\begin{table}[H]\centering\small
\begin{tabular}{p{40mm}p{44mm}p{64mm}}
\toprule
Connection & Connector & Note \\
\midrule
STWIN.box to Pi 4 & USB-C on the probe to USB-A on the Pi & Data and 5 V bus power on one cable. Keep it short and strain-relieved \\
Touchscreen to Pi 4 & DSI ribbon to the DISPLAY connector & Not the 40-pin, and not the CAMERA connector \\
40-pin header & None & Left empty for the whole lab \\
IIS3DWB to STM32U585 & SPI, inside the probe & The 6 kHz vibration path \\
IMP34DT05 to STM32U585 & Digital microphone interface, inside the probe & The digital microphone path \\
IMP23ABSU to STM32U585 & Analog microphone input, inside the probe & The ultrasonic analog path \\
\bottomrule
\end{tabular}
\caption{Connections. This lab has no header pins to number: everything arrives over two cables. Confirm the connector legend on the silkscreen before seating the ribbon.}
\end{table}

\begin{asciiart}
STWIN USB-C  --->  Pi 4 USB-A     (data + 5V bus power)
DSI ribbon   --->  Pi 4 DISPLAY   (not the 40-pin)
40-pin left empty. Strain-relieve the USB cable.
Mount STWIN on the machine under test, not on a ringing desk if you
can avoid it. The desk is a valid first stimulus.
\end{asciiart}

\begin{rulebox}{Before the probe goes on the machine}
CDC-ACM appears only after the firmware declares it. USB-C is the connector, not a protocol, and \texttt{/dev/ttyACM0} is the evidence that the right image is running. Seat the DSI ribbon with the Pi powered off.
\end{rulebox}

\subsection*{Bench layout}

\diagram{p04_bench}{Bench layout. A short USB cable, a probe mounted on the thing under test, and a panel that shows the spectrum. The 40-pin header is visible and empty in this drawing on purpose.}

Mount the STWIN.box on the machine under test, not on a ringing desk if you can avoid it. The desk is a valid first stimulus, and for the first capture it is the right one, because a hand on a desk produces a repeatable excitation that everyone at the bench can perform. Once the pipeline works, move the probe onto something that actually turns.

Keep the panel out of the mechanical path. It is the readout, not part of the experiment, and a panel resting against the machine under test will put its own resonance into every capture.

Strain-relieve the USB cable at both ends. The probe is mounted on something that moves, and the cable is the one part of this lab that is asked to tolerate that movement while carrying high-rate data and the probe's own supply. Tape it to the bench, leave a service loop, and keep the run short.

\subsection*{Software design (UML)}

\diagram{p04_uml}{The boundary the source insists on. Everything left of the dashed line is firmware that moves and labels samples. Everything right of it is Linux, on a core the scheduler does not touch.}

Read the diagram as a contract. The firmware guarantees three properties and no more: samples arrive in order, each frame carries a timestamp, and each frame carries a sequence number that increases by one. Those three are exactly what the host needs to detect a gap, and a gap is the only transport fault that matters here. The host side then owns the window, the transform, the baseline and the decision, which are the four things you will want to change during the lab.

The sequence number is the cheapest instrument in the lab. It costs four bytes per frame and it is the only way to tell a quiet machine from a stalled reader, because both look like an absence of energy in the spectrum. Check it on every frame, and count the gaps rather than logging them, so that a long run produces one number instead of a scroll.

The isolated core is part of the same argument. With \texttt{isolcpus=3 nohz\_full=3} on the kernel command line and \texttt{taskset -c 3} on the analysis process, the completely fair scheduler stops migrating the FFT thread and stops interrupting it with unrelated work. Without it the transform still runs, but its timing wanders with whatever else the desktop is doing, and a wandering pipeline produces dropped frames that look like a hardware fault.

\subsection*{Data flow (ASCII)}

\begin{asciiart}
IIS3DWB --SPI--> STM32U585 --USB-C--> Pi 4 /dev/ttyACM*
IMP34DT05 / IMP23ABSU --I2S/analog--> STM32U585 --same CDC-->
Pi: serial reader -> ring buffer -> window+FFT -> peak bins
     -> NORMAL | WARNING | FAULT  (DSI waterfall)
Boundary: STM32 does DMA + timestamp + sequence + packetize.
          Linux does DSP + policy. Do not FFT on the M33.
\end{asciiart}

Two properties of that block are worth naming. The first is that the same cable carries the vibration stream and both microphone streams, so a transport fault appears in all of them at once, which makes it easy to recognise. The second is that the DSI panel sits on its own connector, so the display costs the analysis nothing on the header and very little on the bus.

\subsection*{Steps}

\begin{steps}
\item \textbf{Flash the probe.} On a PC or on the Pi, load ST HSDatalog or the Zephyr \texttt{steval\_stwinbx1} sensors sample, which enumerates as USB CDC. STM32CubeProgrammer works in USB DFU mode on the same USB-C connector.
\begin{shellcode}
lsusb | grep -i st
# hold BOOT / DFU as in the UM of the STWIN.box, then
STM32_Programmer_CLI -c port=USB1 -w firmware.bin 0x08000000
\end{shellcode}

\item \textbf{Isolate a core} so the scheduler does not migrate the FFT thread.
\begin{shellcode}
# append to /boot/firmware/cmdline.txt:
#   isolcpus=3 nohz_full=3
sudo reboot
ls /dev/ttyACM*
sudo apt-get install -y python3-serial python3-numpy python3-matplotlib
\end{shellcode}

\item \textbf{Ingest.} If the firmware is the Zephyr text dashboard, parse lines. If it is HSDatalog binary, use ST's host decoder first. Do not invent a frame format.
\begin{shellcode}
stty -F /dev/ttyACM0 921600 raw -echo
taskset -c 3 python3 p04_fft.py
\end{shellcode}

\item \textbf{Write the analysis.} \texttt{p04\_fft.py} in outline: read 2048 int16 samples at the configured output data rate, apply a Hann window, call \texttt{numpy.fft.rfft}, take the magnitude, and compare a high-frequency band against a baseline captured on an idle machine. Print NORMAL, WARNING or FAULT.
\begin{pycode}
import numpy as np, serial

N = 2048
win = np.hanning(N)
baseline = None                 # captured once, on an idle machine

ser = serial.Serial("/dev/ttyACM0", 921600, timeout=2)

def band_energy(block):
    spec = np.abs(np.fft.rfft(block * win))
    return spec

while True:
    block = read_2048_samples(ser)      # parser for the firmware you flashed
    spec = band_energy(block)
    hf = spec[len(spec)//2:].sum()      # the high-frequency band
    if baseline is None:
        baseline = hf
        continue
    ratio = hf / baseline
    # WARN and FAULT are measured on your own idle machine.
    # The cookbook prescribes no numbers here: record what you used.
    state = "NORMAL"
    if ratio > WARN:
        state = "WARNING"
    if ratio > FAULT:
        state = "FAULT"
    print(state, round(ratio, 2))
\end{pycode}

The two thresholds are the one part of this pipeline with no prescribed value. Measure them: capture the idle baseline, capture the stimulus, and set the two ratios where they separate the two populations on your bench. Write both numbers in the lab book beside the baseline they belong to, because a threshold without its baseline is not a measurement.

\item \textbf{Optional policy: publish only FAULT.} That is the anomaly-only uplink from draft P20, implemented here as a flag rather than as a second project. The uplink itself belongs to whichever radio lab is seated on another host.

\item \textbf{Check the transport before trusting the spectrum.} Capture two seconds to a file and confirm that the sequence number is monotonic and that no frame is missing. A spectrum computed over a stream with holes in it is a picture of the holes.
\begin{shellcode}
# capture two seconds, then count the gaps offline
timeout 2 cat /dev/ttyACM0 > /home/pi/p04_capture.bin
ls -l /home/pi/p04_capture.bin
# the decoder reports: frames, first and last sequence number, gaps = 0
\end{shellcode}

\item \textbf{Run the orthogonality test and write down both numbers.} Touch the case, note the vibration band. Speak near the digital microphone, note the RMS and the vibration band again. The second number is the one that decides whether the mounting is honest.
\end{steps}

\subsection*{Acceptance test}

\begin{itemize}
\item Touch the STWIN case: the vibration-band energy rises.
\item Speak near the digital microphone: the RMS rises and the vibration band does not. That orthogonality is the lab.
\item A 2-second capture file has a monotonic sequence number and no dropped frames.
\item \texttt{ls /dev/ttyACM*} shows the CDC device, which is the evidence that the firmware declares the interface.
\item The baseline, the two thresholds and the mounting are all written down together, so the run can be repeated next week.
\item The same capture file, replayed through the analysis offline, produces the same classification. If it does not, the pipeline has state it should not have.
\item The classifier prints NORMAL on an idle machine and moves to WARNING or FAULT under the stimulus, against a baseline you captured rather than a constant you guessed.
\item The FFT process is on the isolated core: \texttt{taskset -cp} on its process identifier reports core 3 and nothing else.
\item The 40-pin header is empty at the end of the lab, as it was at the start.
\item Two consecutive idle captures give the same baseline within a tolerance you decided before looking at the second one.
\end{itemize}

\subsection*{Practices}

Short USB cable. No Wi-Fi scan during a capture. Timestamp with \texttt{CLOCK\_MONOTONIC}, because a wall clock can step backwards during a run and a monotonic clock cannot. USB-powered is the Linux-gateway mode; the LiPo is for standalone logging you are not doing today.

\begin{description}
\item[Capture the baseline deliberately.] The idle spectrum of the machine under test is data, not a formality. Record it, keep it in the file next to the captures it is compared against, and recapture it when the probe is remounted.
\item[Keep the policy separate from the transform.] The FFT produces numbers; the thresholds turn numbers into words. Keeping them in different functions is what lets you re-tune the lab without touching the signal path.
\item[One physics per lab.] The IIS3DWB is not the ADXL345 of P01 and P19, and it is not the IKS4A1 of P07. Six different instruments, never interchangeable.
\item[Let the transport prove itself first.] Sequence numbers and frame counts before spectra. The order matters, because a dropped frame produces a broadband artefact that reads like a real mechanical event.
\item[Name the mounting in the file name.] The probe on a desk and the probe on a motor are two different instruments. A capture whose name does not record which one it was cannot be compared with anything later.
\item[Keep the probe on the probe's job.] The anomaly-only uplink is a flag in this lab's policy function. It is not a reason to add a radio to this host, and it is not a second product.
\end{description}

\subsection*{Pitfalls}

\begin{itemize}
\item \textbf{No CDC device after flashing.} USB-C is the connector. CDC-ACM exists only after the firmware declares it, so an absent \texttt{/dev/ttyACM0} points at the image, not at the cable.
\item \textbf{A capture started before the probe settled.} A probe just remounted, or just plugged in, is still moving. Give it a few seconds before the baseline run.
\item \textbf{Inventing a frame format.} HSDatalog produces a binary layout that ST documents and decodes. Guessing at it produces a spectrum that looks plausible and means nothing.
\item \textbf{Running the FFT on the microcontroller.} It moves the part of the lab you most want to change into the part that is hardest to change, and it gives up the host's memory and libraries for no gain.
\item \textbf{Skipping isolcpus.} The transform still runs, but its timing wanders with the desktop's load, and the dropped frames that follow look like a hardware fault.
\item \textbf{A hub between the probe and the host.} Bus power and a sustained high-rate stream through a shared hub produce drops that arrive in bursts. Use a port on the Pi itself.
\item \textbf{A long or unsecured USB cable.} Bus power and high-rate data on a cable that moves with the machine under test produces enumeration faults in the middle of a capture.
\item \textbf{A Wi-Fi scan during a capture.} It costs CPU and bus time at an unpredictable moment. Hold the scan until the file is closed.
\item \textbf{Reading the microphone RMS as vibration.} The two paths share a cable and nothing else. If a sound moves the vibration band, the mounting is resonating and the mechanical setup needs attention before the software does.
\item \textbf{Seating a HAT.} The 40-pin header stays empty in this lab. The DSI panel does not consume it, which is the only reason the panel is here.
\item \textbf{A baseline captured on a running machine.} The reference has to be the idle state. A baseline taken while the machine turns hides exactly the energy the classifier is meant to notice.
\item \textbf{Comparing today's ratio against last week's baseline after remounting.} Mounting stiffness changes the coupling. Recapture the baseline whenever the probe moves.
\item \textbf{Quoting a threshold without its baseline.} The ratios are specific to one mounting on one machine. Reported alone, they describe nothing.
\item \textbf{A wall clock in the timestamp.} \texttt{CLOCK\_MONOTONIC} cannot step backwards; a wall clock can, and a time synchronisation during a capture will put a negative interval in the middle of the file.
\item \textbf{Treating a LiPo run as the same experiment.} USB-powered is the Linux-gateway mode. Standalone logging on the battery is a different lab, and it is not this one.
\end{itemize}

\subsection*{Sources}

\begin{itemize}
\item ST STEVAL-STWINBX1 product page and user manual, for the DFU procedure and the connector legend, \url{https://www.st.com/en/evaluation-tools/steval-stwinbx1.html}
\item Raspberry Pi documentation, the DSI display connector, \url{https://www.raspberrypi.com/documentation/}
\item ST HSDatalog host tools and the documented file format, \url{https://www.st.com/}
\item Zephyr \texttt{steval\_stwinbx1} board support and the sensors sample, \url{https://docs.zephyrproject.org/}
\item IIS3DWB, ISM330DHCX, IMP34DT05 and IMP23ABSU data sheets (ST)
\item The STWIN.box user manual, for the BOOT and DFU procedure referenced in step 1
\item NumPy FFT documentation, \url{https://numpy.org/doc/stable/reference/routines.fft.html}
\item The kernel parameter list for \texttt{isolcpus} and \texttt{nohz\_full}, \url{https://www.kernel.org/doc/html/latest/admin-guide/kernel-parameters.html}
\item P12 in this volume for the DSI panel, and P18 for the other high-rate USB-CDC recorder in this book
\item P10 in this volume, which owns the energy figures this lab does not attempt
\item P07 and P08 in this volume, for the two Nucleo shields that are deliberately not this instrument
\item STM32CubeProgrammer documentation, for the DFU command line used in step 1, \url{https://www.st.com/en/development-tools/stm32cubeprog.html}
\end{itemize}

\project{05}{MCC 118 100 kS/s voltage recorder}{Raspberry Pi 4 + MCC 118 + official DSI touchscreen + POW-BB}{the only calibrated 10 V analog path}

\begin{unique}
Owns the only calibrated 10 V analog path in the book. Not an IMU lab.
\end{unique}

\begin{blueprint}
Draft P2 and P17 claimed Linux IIO with DMABUF and io\_uring on this HAT. The MCC 118 talks SPI to the official daqhats library. There is no mainline IIO driver and no \texttt{IIO\_BUFFER\_DMABUF\_ATTACH\_IOCTL} for this board. Draft P17 also closed a PID loop onto a green LED, which is a demo, not a process plant. We record voltages correctly first.
\end{blueprint}

\subsection*{Intent}

Eight single-ended 12-bit inputs, 100 kS/s as the board maximum, and factory calibration held in the board's EEPROM. The Pi 4 with the DSI panel is the instrument's screen. The POW-BB 3.3 V rail on CH0 is the known reference, which is what turns a plausible number into a checked one.

The product of this lab is a file, not a plot. A hardware-paced scan that runs for one second at 10 kS/s per channel must produce exactly 10 000 samples per enabled channel, with a monotonic sample index and no holes. Everything else in the chapter exists to make that sentence true: the common analog ground, the warm-up, the single-shot sanity read and the refusal to stack a second SPI board.

\begin{kit}
Pi 4, MCC 118, official DSI touchscreen, POW-BB, jumpers.
\end{kit}

\begin{rulebox}{The analog limits before anything is connected}
Never exceed $\pm$10.1 V on any input. The Pi GPIO is not an analog source: do not wire a header pin into a screw terminal to "generate" a test voltage. The 5 V rail is acceptable on CH1 only because 5 is less than 10.1. Warm the board up for one minute before quoting a number.
\end{rulebox}

\begin{note}{Where the 100 kS/s goes}
100 kS/s is the board cap across the scan, not per channel. Two channels at 10 kS/s each is a comfortable first run and is what the acceptance test uses. Eight channels at the cap is a different problem, because the host still has to move, convert and store every sample, and the file writing is not paced by anything.
\end{note}

\subsection*{System architecture}

\diagram{p05_arch}{The supported path. SPI in the kernel, the daqhats library in user space, and no IIO device node anywhere in the diagram. The screen is on the DSI connector, so it costs the header nothing.}

The board is an SPI peripheral and the supported software path is the vendor's userspace library. There is no industrial input and output device node to open, no buffer to map, and no ring to share with the kernel. A Python process calls into daqhats, the library talks SPI through the kernel's SPI interface, and the board returns blocks of samples that it converted on its own clock. That last point is the one that matters: the pacing belongs to the board, not to the process, which is why a busy Pi still produces an evenly spaced file.

The display side is separate and cheap. The official panel is on the DSI connector, so it takes no header pins and shares no bus with the converter. That is the only reason a screen appears in an analog lab at all: it costs the measurement nothing.

The calibration lives in the board's EEPROM and is applied by the library, which is why \texttt{daqhats\_read\_eeproms} is part of the installation rather than an optional step. A board whose EEPROM has not been read into the host's configuration will still return numbers, and those numbers will be wrong in a way that looks like a wiring fault.

\begin{table}[H]\centering\small
\begin{tabular}{p{44mm}p{102mm}}
\toprule
Call & What it does \\
\midrule
\texttt{daqhats\_read\_eeproms} & Reads the factory calibration out of the board's ID EEPROM. Part of the installation \\
\texttt{daqhats\_list\_boards} & Names the boards the host can see, and at which address \\
\texttt{hat\_list} & The same question from Python, filtered by board type \\
\texttt{mcc118(0)} & Opens board 0, which is what all address jumpers off means \\
\texttt{a\_in\_read} & One single-shot reading from one channel. Sanity, not product \\
\texttt{a\_in\_scan\_start} & Hands the pacing to the board: channel mask, sample count, rate \\
\texttt{a\_in\_scan\_read} & Blocks until the samples arrive or the timeout expires \\
\texttt{a\_in\_scan\_stop} & Ends the conversion \\
\texttt{a\_in\_scan\_cleanup} & Releases the library's scan buffer \\
\bottomrule
\end{tabular}
\caption{The library calls this lab uses. The last four are the hardware-paced scan, and they belong together in that order.}
\end{table}

\subsection*{Wiring and schematic}

\diagram{p05_schematic}{The screw terminals. CH0 carries the POW-BB 3.3 V rail as the known reference, CH1 carries the 5 V rail, and AGND is the common analog ground that makes both of them meaningful.}

\begin{table}[H]\centering\small
\begin{tabular}{p{36mm}p{30mm}p{24mm}p{62mm}}
\toprule
Signal & Connection & Pi header pin & Note \\
\midrule
CH0 & POW-BB 3V3 & Screw terminal & Expect about 3.30 V, the known reference \\
CH1 & POW-BB 5V & Screw terminal & Only because 5 is less than 10.1 \\
AGND & POW-BB GND & Screw terminal & Common analog ground \\
DGND & Pi GND & Screw terminal & Digital ground \\
SPI0 CE0 & Through the header & 24 & GPIO8, the board's chip select \\
SPI0 SCLK & Through the header & 23 & GPIO11 \\
SPI0 MOSI & Through the header & 19 & GPIO10 \\
SPI0 MISO & Through the header & 21 & GPIO9 \\
Board address & Address jumpers & 32, 33, 37 & GPIO12, GPIO13, GPIO26. All jumpers off is board 0 \\
ID EEPROM & HAT ID pins & & Read once with \texttt{daqhats\_read\_eeproms} \\
\bottomrule
\end{tabular}
\caption{Wiring. The HAT sits alone on the 40-pin header. Confirm the terminal order on the silkscreen before landing a wire, since the block legend is printed on the board.}
\end{table}

\begin{asciiart}
MCC 118 seated ALONE on the 40-pin
  uses SPI0 CE0, ID EEPROM, address GPIO12/13/26
  address jumpers all OFF = board 0

Screw terminals:
  CH0  ---- POW-BB 3V3     expect ~3.30 V
  CH1  ---- POW-BB 5V      only because 5 < 10.1
  AGND ---- POW-BB GND     common analog ground
  DGND ---- Pi GND
Never exceed +/-10.1 V. Pi GPIO is not an analog source.
Warm up 1 minute before quoting a number.
\end{asciiart}

\subsection*{Bench layout}

\diagram{p05_bench}{Bench layout. Four wires from the rail splitter into the terminal block, a panel on the DSI connector, and nothing else on the header. The multimeter is optional but it is the only way to know the reference is what the label says.}

Keep the analog wires short and keep them away from the supply lead. This is a bench with a rail splitter in it, and the two rails that matter are the ones being measured, so anything that couples into those four wires becomes part of the reading.

Land the four wires, tighten the terminals and check them with a gentle pull before power. Then leave the board alone for a minute. A converter that has just been powered is still settling, and a number quoted in the first seconds is a number about temperature rather than about the rail.

\subsection*{Software design (UML)}

\diagram{p05_uml}{The hardware-paced scan. The library call that blocks is \texttt{a\_in\_scan\_read}, and the board is filling its buffer on its own clock the whole time. The two closing calls are not optional.}

A closed-loop demonstration belongs after all of this, if at all. Driving an output from a reading is easy to show and easy to misread as process control, and the draft that did it had not yet produced an honest file. Record correctly first, then decide whether the loop adds anything.

Four calls make the product. \texttt{a\_in\_scan\_start} tells the board which channels to convert and how fast, and from that moment the board paces itself. \texttt{a\_in\_scan\_read} blocks until the requested samples have arrived or the timeout expires, and it returns a block whose data are already scaled and calibrated. \texttt{a\_in\_scan\_stop} ends the conversion, and \texttt{a\_in\_scan\_cleanup} releases the buffer the library allocated. Leaving either of the last two out leaves the board scanning into a buffer nobody is emptying, and the next run then starts against a board that is already busy.

\subsection*{Data flow (ASCII)}

\begin{asciiart}
  POW-BB                 MCC 118 (board 0)            Pi 4
  +-----------+          +---------------------+      +----------------------+
  | 3V3 rail  |--CH0---->|                     |      | daqhats (userspace)  |
  | 5V  rail  |--CH1---->| 8 ch, 12-bit        |      |   a_in_scan_start    |
  | GND       |--AGND--->| 100 kS/s max        |      |   a_in_scan_read <-- blocks
  +-----------+          | factory cal in      | SPI0 |   a_in_scan_stop     |
                         | the ID EEPROM       |<====>|   a_in_scan_cleanup  |
   Pi GND ------DGND---->|                     | CE0  |          |           |
                         +---------------------+      |          v           |
                                                      | CSV: index,ch0,ch1   |
                                                      |          |           |
                                                      |          v           |
                                                      | matplotlib on the DSI|
                                                      +----------------------+
  Limits: never exceed +/-10.1 V.  Pi GPIO is not an analog source.
\end{asciiart}

One more property of that drawing deserves attention. The POW-BB is the only voltage source in the lab, and it is a rail splitter with printed labels rather than a calibrator. It is good enough to be a known reference on a bench because its rails are regulated and because a multimeter can confirm them in ten seconds. It is not a traceable standard, and nothing in this chapter claims that it is.

\subsection*{Steps}

\begin{steps}
\item \textbf{Install the library and read the board.} The EEPROM read is part of the installation, not an extra.
\begin{shellcode}
git clone https://github.com/mccdaq/daqhats.git
cd daqhats && sudo ./install.sh
sudo daqhats_read_eeproms
daqhats_list_boards
\end{shellcode}

\item \textbf{Land the wires and wait.} Tighten the four screw terminals, confirm the two rails with a multimeter and write the numbers in the lab book, then power the board and leave it for a minute before reading anything. The wait is part of the procedure, not a courtesy.

\item \textbf{Single-shot sanity.} One read per channel, to confirm the reference before anything is paced.
\begin{pycode}
from daqhats import mcc118, hat_list, HatIDs
print(hat_list(filter_by_id=HatIDs.MCC_118))
h = mcc118(0)
print([round(h.a_in_read(ch), 4) for ch in range(8)])
\end{pycode}

\item \textbf{Hardware-paced scan.} This is the product, not the single read.
\begin{pycode}
from daqhats import mcc118, OptionFlags
h = mcc118(0)
h.a_in_scan_start(0x03, 10000, 10000, OptionFlags.DEFAULT)
block = h.a_in_scan_read(20000, 5000)
print(len(block.data), block.data[:8])
h.a_in_scan_stop(); h.a_in_scan_cleanup()
\end{pycode}

The channel mask \texttt{0x03} selects CH0 and CH1. At 10 kS/s per channel, a 10 000 sample request is one second of recording per channel, and the second argument of \texttt{a\_in\_scan\_read} is a timeout in milliseconds. The returned block carries the two channels interleaved, which is why the sample count that comes back is twice the per-channel figure.

\item \textbf{Write the file.} One second at 10 kS/s on CH0 and CH1, as CSV with a monotonic sample index, then plot it on the DSI panel with matplotlib.
\begin{pycode}
import csv
from daqhats import mcc118, OptionFlags

h = mcc118(0)
h.a_in_scan_start(0x03, 10000, 10000, OptionFlags.DEFAULT)
block = h.a_in_scan_read(20000, 5000)
h.a_in_scan_stop(); h.a_in_scan_cleanup()

data = block.data                      # CH0, CH1, CH0, CH1, ...
with open("/home/pi/p05_scan.csv", "w", newline="") as f:
    w = csv.writer(f)
    w.writerow(["index", "ch0", "ch1"])
    for i in range(len(data) // 2):
        w.writerow([i, round(data[2*i], 6), round(data[2*i + 1], 6)])
print("rows:", len(data) // 2)         # expect 10000
\end{pycode}

\item \textbf{Check the file before you look at the picture.} Count the rows, check that the index is monotonic, and confirm that the CH0 column sits on the reference. A plot hides a hole; a row count does not.
\begin{shellcode}
wc -l /home/pi/p05_scan.csv          # 10001 lines: 10000 rows and one header
head -n 3 /home/pi/p05_scan.csv
awk -F, 'NR>1 {s+=$2; n++} END {print "mean ch0:", s/n}' /home/pi/p05_scan.csv
\end{shellcode}
\end{steps}

\subsection*{Acceptance test}

\begin{itemize}
\item CH0 reads 3.30 V plus or minus 0.05 V against the POW-BB rail, after the one-minute warm-up.
\item A 10 kS/s by 1 s file has 10 000 samples per enabled channel and no holes.
\item \texttt{daqhats\_list\_boards} names one board at address 0, which is what all jumpers off means.
\item The sample index in the file is monotonic from zero with no repeats.
\item Nothing else is seated on the 40-pin header, and \texttt{daqhats\_list\_boards} still reports the board after a reboot.
\item CH1 sits near 5 V, which is inside the limit and is the only reason that rail is on an input at all.
\item The multimeter reading of the 3.3 V rail and the CH0 column of the file agree, and both are written down.
\item The scan ends with \texttt{a\_in\_scan\_stop} and \texttt{a\_in\_scan\_cleanup}, and a second run starts cleanly without a reboot.
\end{itemize}

\subsection*{Practices}

Common AGND. Do not stack the 3.5 inch LCD or the Explorer700: both want SPI, and this board owns it for the duration. 100 kS/s is the board cap, so budget CPU before enabling all eight channels: the conversion is paced by the board, but the transfer and the file writing are not. A closed-loop demonstration that drives a PWM output at an LED is a stretch goal only after the file is honest.

\begin{description}
\item[Quote the reference with the reading.] The sentence is "CH0 reads 3.30 V against the POW-BB 3.3 V rail", not "the board reads 3.30". One of those can be checked next week and the other cannot.
\item[Warm up before quoting.] One minute, every session. It costs a minute and it removes a whole class of argument about the first decimal place.
\item[Count the rows before plotting.] A plot interpolates over a missing sample and a row count does not. The file is the product, and the picture is a view of it.
\item[Treat the EEPROM as part of the instrument.] The factory calibration is the difference between this board and a generic converter. Read it once at installation and do not assume it survived a reflash of the host.
\end{description}

\subsection*{Pitfalls}

\begin{itemize}
\item \textbf{Treating the single-shot read as the measurement.} \texttt{a\_in\_read} is a sanity check. It is not paced, so a sequence of them is not a recording and their spacing is whatever the interpreter managed.
\item \textbf{Looking for an IIO device.} There is none. The supported path is SPI plus the daqhats userspace library, and time spent looking for \texttt{/sys/bus/iio} is time spent on a device that does not exist here.
\item \textbf{A second SPI board on the header.} The 3.5 inch LCD of P06 and P19 and the Explorer700 of P11 both want the same bus. One 40-pin HAT per host, and this lab's HAT is the MCC 118.
\item \textbf{No common analog ground.} Without AGND tied to the POW-BB ground, the inputs float and the readings drift with whatever else is on the bench. The symptom is a reference that moves when a hand comes near the wires.
\item \textbf{An input left floating.} An unconnected channel reads whatever the input stage drifts to. It is not a zero, and it should not be quoted as one.
\item \textbf{A Pi GPIO used as an analog source.} It is a digital output with a series impedance you do not control, and it is not a voltage reference. The rail splitter is.
\item \textbf{Quoting the first reading.} A converter that has just been powered is still settling. Wait the minute.
\item \textbf{Skipping the EEPROM read.} The board returns numbers either way. Only one set of them is calibrated.
\item \textbf{A rail measured at the splitter but read at the terminal.} If the two disagree, the wire or the terminal is the cause, and that is worth knowing before the converter is suspected.
\item \textbf{Leaving the scan running.} Without \texttt{a\_in\_scan\_stop} and \texttt{a\_in\_scan\_cleanup}, the board keeps converting into a buffer nobody empties, and the next run meets a busy board.
\item \textbf{Reading the block as one channel.} With a two-channel mask the samples are interleaved. Splitting them incorrectly produces a file that looks like noise at half amplitude.
\item \textbf{Two boards at the same address.} All jumpers off is board 0. If a second MCC board ever joins the bench, the address jumpers are how they are told apart, and \texttt{daqhats\_list\_boards} is where the answer appears.
\item \textbf{A loose screw terminal.} It reads as noise on one channel and nothing on the others. Pull each wire gently after tightening.
\item \textbf{Calling the plot the result.} The acceptance test counts samples. A picture with a straight line through a gap passes no test.
\item \textbf{Enabling all eight channels on the first run.} 100 kS/s is the board maximum across the scan, and the host still has to move and store every sample. Start with two.
\end{itemize}

\subsection*{Sources}

\begin{itemize}
\item MCC DAQ HAT library, installation and the Python API, \url{https://github.com/mccdaq/daqhats}
\item MCC 118 documentation, input range, resolution and the 100 kS/s board maximum, \url{https://mccdaq.github.io/daqhats/}
\item The daqhats Python examples, which include a scan written in the same four calls used here
\item Digilent and Measurement Computing MCC 118 product documentation, terminal block legend and address jumpers
\item P06 in this volume, the other SPI board in the book, for why the two cannot share a sitting
\item JOY-iT SBC-POW-BB documentation, for the rail voltages used as the reference
\item NumPy and matplotlib documentation, for the plot drawn on the panel, \url{https://matplotlib.org/}
\item Raspberry Pi hardware documentation, 40-pin header and SPI0, \url{https://www.raspberrypi.com/documentation/computers/raspberry-pi.html}
\item P12 in this volume for the DSI panel, P15 for the rail discipline this lab depends on, and P11 and P19 for the two boards that must stay off the header
\end{itemize}

\project{06}{8x8 ToF occupancy kiosk}{Raspberry Pi 3 + 3.5 inch LCD + Nucleo-H7A3ZI-Q + X-NUCLEO-53L8A1}{multi-zone ranging and the SPI heat map}

\begin{unique}
Owns VL53L8CX multi-zone ranging and the Waveshare SPI LCD as a local heat map.
Not an IMU and not V4L2.
\end{unique}

\begin{blueprint}
Draft P8 loaded stsw-img042 as a V4L2 camera node and drove the LCD with flexfb/fbtft.
On Bookworm the Waveshare 3.5 inch (A) uses the vendor overlay / DRM path, and ST's ToF
Linux stack is not a Pi camera. Stream 64 millimetre values over the Nucleo USB-CDC.
\end{blueprint}

\subsection*{Intent}

The Nucleo-H7A3ZI-Q carries the X-NUCLEO-53L8A1 and produces an 8x8 grid of millimetre
ranges. The Pi 3, wearing the Waveshare 3.5 inch LCD (A), paints occupancy or desk
presence from those 64 numbers. Nothing else happens on either board. The ToF sensor is
an instrument with its own ranging core and its own timing, and the useful product of
this lab is a grid that arrives complete, frame after frame, and lands on glass a person
can read from across the room.

This is a two-host lab, and that is the design rather than an accident. The 40-pin header
on the Pi 3 is taken by the LCD, so the ToF sensor cannot sit there. The Nucleo owns the
sensor and hands Linux a line of text; Linux owns storage, the framebuffer and the
threshold policy. Read the boundary as the book's standing rule made physical: one USB
cable, firmware on one side, Linux on the other. P19 uses the same LCD on the same header
for an ADXL345 tilt meter, so P19 and P06 cannot share a Pi in one sitting. Do that on a
different afternoon.

\begin{kit}
Pi 3, Waveshare 3.5 inch RPi LCD (A), Nucleo-H7A3ZI-Q, X-NUCLEO-53L8A1, USB cable.
\end{kit}

\subsection*{System architecture}

\diagram{p06_arch}{Two hosts, one USB cable. Everything left of the cable is STM32
firmware, everything right of it is Linux. The VL53L8CX is never visible to the Pi as an
I2C device; only the 64 millimetre values cross the boundary.}

The figure has one load-bearing feature: the single USB cable in the middle. On the
firmware side the VL53L8CX sits on the STM32 I2C bus at address 0x29 and is serviced by
ST's driver inside the sketch, which collects a complete 8x8 result and formats one text
line per frame. On the Linux side that line arrives on a CDC ACM character device, and a
Python consumer turns it into 64 integers. No ToF driver, no V4L2 node and no ST Linux
stack is installed on the Pi, because none of those is the supported path on Bookworm.

The second half of the Linux side is the panel. The 3.5 inch LCD (A) is an SPI display
that the vendor overlay binds at boot, so the painting code writes to a framebuffer or a
DRM device rather than to the SPI bus itself. That separation keeps the lab honest: if
the grid stops arriving, the display stays up and shows stale values, which is the
symptom you want rather than a blank screen that could mean either problem.

\subsection*{Wiring and schematic}

\diagram{p06_schematic}{The two stacks. On the left, the 53L8A1 on the Nucleo Arduino UNO
headers with J9 set to I2C. On the right, the LCD seated on the Pi 3 40-pin header, which
it owns alone. The only electrical link between the two is the USB cable at the bottom.}

\begin{table}[H]\centering\small
\begin{tabular}{p{40mm}p{30mm}p{20mm}p{58mm}}
\toprule
Signal or part & Connector & Pin & Note \\
\midrule
X-NUCLEO-53L8A1 & Nucleo Arduino UNO headers & stacked & One shield only, no flying leads \\
J9 SPI\_I2C\_N & 53L8A1 jumper & 2-3 & Selects I2C, the default for this lab \\
VL53L8CX I2C & STM32 I2C bus & 0x29 & Behind the STM32, not on the Pi \\
Cover glass and spacer & sensor aperture & 0.5 mm & As shipped, do not remove \\
Nucleo USER or ST-LINK USB & Pi 3 USB-A & any & CDC ACM, one cable, the only link \\
LCD 3V3 & Pi 40-pin & 1 & Panel supply from the header \\
LCD 5V & Pi 40-pin & 2 & Backlight rail on the HAT \\
LCD GND & Pi 40-pin & 6 & Also pins 9, 20, 25 \\
LCD MOSI & Pi 40-pin & 19 & GPIO10, SPI0 \\
LCD MISO & Pi 40-pin & 21 & GPIO9, SPI0, touch read-back \\
LCD SCLK & Pi 40-pin & 23 & GPIO11, SPI0 \\
LCD CE0 & Pi 40-pin & 24 & GPIO8, panel chip select \\
LCD CE1 & Pi 40-pin & 26 & GPIO7, touch controller chip select \\
LCD control and touch IRQ & Pi 40-pin & vendor overlay & The exact GPIO assignment comes
from the vendor overlay, confirm against the Waveshare wiki and the silkscreen \\
\bottomrule
\end{tabular}
\caption{Wiring. The LCD owns the Pi 40-pin header, so nothing else seats on it, and the
ToF sensor reaches Linux only over USB.}
\end{table}

\begin{asciiart}
53L8A1 stacked on Nucleo Arduino UNO headers.
J9 SPI_I2C_N = I2C (pins 2-3) unless you deliberately choose SPI.
Nucleo USER USB or ST-LINK USB --> Pi USB-A   (CDC ACM)
Waveshare 3.5" LCD seated on Pi 3 40-pin (this lab's HAT).
Two hosts, one USB cable between them.
Point the aperture at 200-1500 mm of free space; no sunlight in the FOV.
Cover glass + 0.5 mm spacer as shipped.
\end{asciiart}

The shield reaches the STM32 through the Arduino UNO headers and nothing else. There are
no flying leads in this lab, no level shifter and no second I2C master. If you deliberately
choose SPI on J9 instead, the firmware example has to change with it, so make that a
decision you write in the lab book rather than a jumper you find in the wrong position
later.

\subsection*{Bench layout}

\diagram{p06_bench}{Bench layout. The Pi 3 wears the LCD as its HAT, the Nucleo wears the
53L8A1 as its shield, and one USB cable joins them. The sensor aperture looks into 200 to
1500 mm of free space with no sunlight in the field of view.}

Put the two stacks side by side with the LCD facing you and the ToF aperture facing the
space you want to measure. Give the sensor a clear volume between 200 and 1500 mm, and
keep direct sunlight out of that volume, because ambient infrared raises the noise floor
and the far zones start to report nonsense. Leave the cover glass and the 0.5 mm spacer
as they shipped: the calibration in the part assumes them.

\begin{note}{One header and one HAT per sitting}
The 3.5 inch LCD (A) is one of the six boards that own the 40-pin header outright. While
it is seated, the MCC 118, the Explorer700 and every cellular HAT stay in their bags, and
P19 waits for another afternoon even though it wants the same panel. Power the Pi off
before unseating anything, and build one lab at a time.
\end{note}

\subsection*{Software design (UML)}

\diagram{p06_uml}{One frame, from the ranging core to the glass. The firmware waits for
data-ready, reads a complete 8x8 result, and prints one line. Linux parses 64 integers,
applies the threshold it measured against an empty field of view, and paints 64
rectangles.}

The consumer is deliberately dull. It reads a line, checks that the line starts with
\texttt{Z,}, splits on commas and converts to integers. A line that fails any of those
checks is dropped rather than repaired, because a partial frame painted on the glass
looks like a real reading and is not. Occupancy is one comparison per zone against a
threshold you measured yourself with the field of view empty, not a constant somebody
wrote down.

\subsection*{Data flow (ASCII)}

\begin{asciiart}
  Nucleo-H7A3ZI-Q  (firmware)                Pi 3  (Linux)
  +------------------------------+           +-------------------------------+
  | X-NUCLEO-53L8A1              |           | /dev/ttyACM0  115200 8N1      |
  |   VL53L8CX 8x8 ToF           |           |   |                           |
  |   I2C 0x29, J9 = I2C         |           |   v                           |
  |        |                     |           | p06_kiosk.py                  |
  |        v  data-ready         |    USB    |   startswith("Z,") ? drop : ok|
  | ST driver: read 64 ranges    |==========>|   64 ints, min / max in mm    |
  |        |                     |  CDC ACM  |   range < threshold = occupied|
  |        v                     |           |   |                           |
  | printf "Z,d00,d01,...,d63"   |           |   v                           |
  +------------------------------+           | 8x8 rectangles -> /dev/fb0    |
                                             +-------------------------------+
                                                         |
                                                         v  SPI0, vendor overlay
                                             Waveshare 3.5" LCD (A) 480x320

  Nothing crosses the USB cable except text frames. The Pi has no ToF driver,
  no V4L2 node and no I2C path to the sensor: i2cdetect on the Pi is empty.
\end{asciiart}

\subsection*{Steps}

\begin{steps}
\item \textbf{Flash the Nucleo.} Use the STM32duino
\texttt{X\_NUCLEO\_53L8A1\_HelloWorld\_I2C} sample or the Cube VL53L8CX example. Whichever
you pick, make it print one line per frame, in millimetres:
\begin{plaincode}
Z,d00,d01,...,d63
\end{plaincode}
That format is the contract between the two hosts. Sixty-five comma-separated fields, the
first of which is the literal \texttt{Z}.

\item \textbf{Stack one shield and one shield only.} The 53L8A1 goes on the Nucleo Arduino
UNO headers. Set J9 SPI\_I2C\_N to I2C, pins 2-3, unless you deliberately choose SPI. Do
not also stack IKS4A1 or IKS5A1: those are P07 and P08, and the Arduino stack takes one
shield.

\item \textbf{Bring the LCD up on Bookworm.} Follow the current Waveshare 3.5 inch (A)
note for Bookworm, which uses the vendor \texttt{dtoverlay}, not the 2016 flexfb recipe.
Then confirm that a framebuffer actually exists before you write to it.
\begin{shellcode}
ls /dev/fb*
ls /dev/dri/
\end{shellcode}
A \texttt{/dev/fb0} or a DRM card must be present. If neither is, the overlay did not load
and no amount of Python will paint anything.

\item \textbf{Find the Nucleo.} Connect the single USB cable from the Nucleo USER USB or
ST-LINK USB port to a Pi USB-A port.
\begin{shellcode}
lsusb
ls /dev/ttyACM* 2>/dev/null
sudo apt-get install -y python3-serial
\end{shellcode}
More than one \texttt{ttyACM} node can appear, because the ST-LINK enumerates as well.
The node that carries \texttt{Z,} lines is the one you want.

\item \textbf{Run the consumer.} Save this as \texttt{/home/pi/p06\_kiosk.py} and confirm
the numbers before you paint anything.
\begin{pycode}
import serial
ser = serial.Serial("/dev/ttyACM0", 115200, timeout=1)
while True:
    line = ser.readline().decode(errors="ignore").strip()
    if line.startswith("Z,"):
        z = [int(x) for x in line.split(",")[1:]]
        print(min(z), max(z), "mm", len(z))
\end{pycode}
The printed length must be 64 on every frame. If it is not, the firmware line is truncated
and the serial settings or the print statement are the cause, not the sensor.

\item \textbf{Measure the threshold, do not guess it.} With the field of view empty, watch
the printed minimum for a minute and write the wall distance in the lab book. Occupied is
a range below a threshold you set under that number with margin. Repeat the measurement if
you move the bench, because the wall distance is a property of where the sensor sits and
not of the part.

\item \textbf{Paint the grid.} Draw an 8x8 rectangle grid on the framebuffer, one rectangle
per zone, filled by range. Occupied means the zone range is below the threshold you just
measured against the empty field of view. Keep the painting in the same loop as the parse,
so a stalled sensor stops the repaint rather than leaving a stale frame that looks live.

\item \textbf{Tear the stack down when the acceptance test passes.} Power the Pi off before
you unseat the LCD, put the 53L8A1 back in its own bag, and leave the Nucleo bare for P18.
The next lab that wants this Pi wants the header empty.
\end{steps}

\subsection*{Acceptance test}

\begin{itemize}
\item A hand at 300 mm occupies 4 to 10 zones.
\item An empty field of view reports the wall at a stable range, plus or minus 20 mm.
\item The LCD shows a coarse heat map, 64 values every frame.
\item \texttt{len(z)} is 64 on every accepted line, and lines that do not start with
\texttt{Z,} are dropped rather than parsed.
\end{itemize}

\subsection*{Practices}

One Arduino shield on the Nucleo. Do not also stack IKS4A1 or IKS5A1. One 40-pin HAT on
the Pi, which in this lab is the LCD, so the MCC 118, the Explorer700 and any cellular HAT
stay on the shelf. The ToF sensor is one of six different instruments in this kit and it is
not interchangeable with the ADXL345 or with the STWIN.box.

Keep the millimetre grid in its own log file if you record it, and keep the threshold you
measured in the lab book beside the wall distance it came from. A threshold without the
measurement behind it is a number somebody will copy into the next bench, where the wall
is somewhere else.

\subsection*{Pitfalls}

\begin{itemize}
\item \textbf{Looking for the ToF sensor with \texttt{i2cdetect} on the Pi.} It is not
there. The VL53L8CX sits on the STM32 I2C bus at 0x29 and reaches Linux only as text over
USB. P07 makes that boundary its teaching point; here it is simply the reason the scan is
empty.
\item \textbf{Installing the ST Linux ToF stack.} stsw-img042 is not a drop-in Bookworm
camera device and the Pi has no V4L2 node for this part. The symptom is a long detour that
ends without a \texttt{/dev/video} node.
\item \textbf{Following the 2016 flexfb or fbtft recipe.} On Bookworm the 3.5 inch (A) uses
the vendor overlay and the DRM path. The symptom is a white or black panel with no error,
because nothing failed loudly.
\item \textbf{Sunlight or a bright window in the field of view.} Ambient infrared raises
the noise floor. The symptom is far zones that jump by hundreds of millimetres between
frames while near zones stay calm.
\item \textbf{Stacking a second Arduino shield.} The IKS4A1 and IKS5A1 want the same stack
and share addresses with each other. The symptom is a firmware scan that finds the wrong
part, or nothing at all.
\item \textbf{Trying to run P19 on the same Pi in the same sitting.} P19 wants the same LCD
on the same header with an ADXL345 on the pass-through. Power off, tear the HAT down, and
build one lab at a time.
\item \textbf{Painting a partial frame.} A line that lost characters still parses into a
short list. Check the length, drop the frame, and let the glass hold the last good grid
rather than a half-drawn one.
\item \textbf{The wrong ACM node.} With the ST-LINK enumerating as well, more than one
\texttt{ttyACM} can appear. The wrong one is silent rather than an error, so check which
node carries \texttt{Z,} lines.
\item \textbf{Removing the cover glass or the spacer.} The part is calibrated with them
fitted. The symptom is a constant offset across every zone that no threshold explains.
\end{itemize}

\subsection*{Sources}

\begin{itemize}
\item ST X-NUCLEO-53L8A1 expansion board page,
\url{https://www.st.com/en/ecosystems/x-nucleo-53l8a1.html}
\item ST VL53L8CX 8x8 multizone ranging sensor data sheet,
\url{https://www.st.com/en/imaging-and-photonics-solutions/vl53l8cx.html}
\item STM32duino X-NUCLEO-53L8A1 library and the HelloWorld I2C example,
\url{https://github.com/stm32duino/X-NUCLEO-53L8A1}
\item ST NUCLEO-H7A3ZI-Q board page,
\url{https://www.st.com/en/evaluation-tools/nucleo-h7a3zi-q.html}
\item Waveshare 3.5 inch RPi LCD (A) wiki, for the current Bookworm overlay note,
\url{https://www.waveshare.com/wiki/3.5inch_RPi_LCD_(A)}
\item Raspberry Pi hardware documentation, 40-pin header and SPI0,
\url{https://www.raspberrypi.com/documentation/computers/raspberry-pi.html}
\item Linux USB CDC ACM driver, for the \texttt{ttyACM} naming,
\url{https://www.kernel.org/doc/html/latest/usb/index.html}
\end{itemize}

\project{10}{PPK2 energy characterisation bench}{Raspberry Pi 4 + Nordic PPK2 + ESP32 NodeMCU as the DUT}{microamp figures}

\begin{unique}
Owns microamp figures. No telemetry product.
\end{unique}

\begin{blueprint}
Draft P9 put the PPK2 in series with a SIM7070 HAT that was still powered from the Pi 5 V
pins. That measures nothing useful. Either lift the HAT 5 V pins and feed VOUT into the
modem power input, or measure a flying-lead DUT (ESP32) that is not also USB-powered. PPK2
source mode is 5 V-class and about 1 A; an LTE TX peak can exceed that, so do not source a
Cat-4 power amplifier from the PPK2.
\end{blueprint}

\subsection*{Intent}

Three firmware states on the ESP32: radio-off idle, BLE advertise, and deep-sleep. The
Pi 4 logs the PPK2 stream and you write down two numbers per state, a median and a p95.
That is the product of this lab. There is no telemetry here, no broker and no uplink: the
deliverable is a small table of current figures that somebody else can act on.

The single idea that makes the measurement valid is that the DUT has exactly one source of
power, and the PPK2 is it. An ESP32 that is also plugged into USB draws most of its current
through the USB rail, and the PPK2 then reports the small remainder that happens to come
through its own output. Traces taken that way look reasonable and mean nothing. Pulling the
DUT USB cable is not a detail of the setup; it is the experiment.

\begin{kit}
Pi 4, PPK2, ESP32 NodeMCU as the DUT, POW-BB, jumpers.
\end{kit}

\subsection*{System architecture}

\diagram{p10_arch}{The measurement path and the control path. The PPK2 sits in the DUT
power rail as a source-measure instrument, and it reports samples to the Pi over its own
USB cable. The DUT's own USB is not part of the picture, because it is not connected.}

Read the figure as two separate cables that meet only inside the PPK2. The power path runs
from a bench 5 V supply into PPK2 VIN, through the instrument, and out of VOUT into the
ESP32 3V3 pin, in source mode, set to 3300 mV. The measurement path runs from the PPK2 over
USB to the Pi 4, which logs samples at 100 kS/s. Ground is common across the PPK2, the DUT
and the Pi, because a current measurement with two ground references is not a measurement.

The instrument has limits and they are part of the design. The PPK2 in source mode is
5 V-class and about 1 A. That is ample for an ESP32 in any of the three states, and it is
not ample for a cellular power amplifier: a Cat-4 transmit peak can exceed it. This is why
the lab measures a flying-lead ESP32 and not the modem from P01, and why the source's
correction of the blueprint draft is repeated on the schematic itself.

\subsection*{Wiring and schematic}

\diagram{p10_schematic}{The whole lab in one drawing. Bench 5 V into VIN, VOUT to the ESP32
3V3 pin at 3300 mV in source mode, one common ground, PPK2 USB to the Pi, and the ESP32 USB
cable drawn where it is not: disconnected.}

\begin{table}[H]\centering\small
\begin{tabular}{p{34mm}p{34mm}p{26mm}p{56mm}}
\toprule
From & To & Setting & Note \\
\midrule
Bench 5 V (or Pi USB) & PPK2 VIN & 5 V & The supply the instrument draws from \\
PPK2 VOUT & ESP32 3V3 pin & source mode, 3300 mV & The DUT's only power source \\
PPK2 GND & DUT GND & common & Also common with the Pi \\
DUT GND & Pi GND & common & One ground reference for the whole bench \\
PPK2 USB & Pi USB-A & CDC & The sample stream and the instrument control \\
ESP32 USB & nothing & disconnected & This is the whole experiment \\
POW-BB & bench rails & labelled & Rails and a common ground, as in P15 \\
Sample rate & PPK2 to Pi & 100 kS/s & Two seconds per firmware state \\
\bottomrule
\end{tabular}
\caption{Wiring. Two cables carry everything: one power path through the instrument, one
USB path to the host. The third cable, the DUT's own USB, stays out of the bench.}
\end{table}

\begin{asciiart}
PPK2 VIN   <--  Pi USB or bench 5V
PPK2 GND   ---  DUT GND --- Pi GND (common)
PPK2 VOUT  -->  ESP32 3V3 pin     source mode, set 3300 mV
PPK2 USB   -->  Pi USB-A
ESP32 USB  DISCONNECTED           this is the whole experiment
\end{asciiart}

Five lines of wiring, and the fifth one is the experiment. The first four put the
instrument in series with the only supply the DUT has, and the fifth removes the supply it
would otherwise prefer. A bench that gets the first four right and the fifth wrong produces
traces that look like measurements, which is worse than producing nothing.

\begin{rulebox}{Two limits that do not move}
The PPK2 in source mode is 5 V-class and about 1 A. Never source a Cat-4 power amplifier
from it: an LTE transmit peak can exceed that current, and the instrument is not a bench
supply for a radio. And never leave the DUT USB cable connected during a capture, because
a second power path turns every number in the run into an unknown fraction of the truth.
\end{rulebox}

\subsection*{Bench layout}

\diagram{p10_bench}{Bench layout. The ESP32 sits on its own with two flying leads and no
USB cable. The PPK2 is between the supply and the DUT, and the only cable leaving the bench
goes to the Pi.}

Lay the bench out so that the disconnected USB cable is visibly elsewhere, not coiled
beside the DUT where somebody will plug it back in between runs. Bring 3V3 and ground to
the ESP32 with two labelled jumpers from the PPK2 output, use the POW-BB discipline from
P15 for any other rail on the bench, and check the ground continuity with a meter before
the first capture rather than after the third.

\begin{note}{The rails come from P15}
This bench uses the POW-BB the way P15 teaches it: labelled rails, a measured 5 V, a
measured 3.3 V and a common ground written in the lab book. Nothing here needs a 5 V pull-up
and nothing here goes onto a GPIO, so the only electrical discipline this lab adds to P15 is
that the DUT has one supply and the instrument is it.
\end{note}

\subsection*{Software design (UML)}

\diagram{p10_uml}{One capture run, across the three firmware states. Each state is flashed,
allowed to settle, captured for two seconds at 100 kS/s, and reduced to a median and a p95.
The run ends with a comparison: if the three traces overlay, the DUT still has a second
power source.}

The activity is deliberately repetitive, because the value of the lab comes from doing the
same thing three times. Flash, settle, capture two seconds, reduce to two numbers, write
them down. The reduction matters as much as the capture: a median says what the state
normally costs, and a p95 says what it costs when the radio is busy, and neither of them is
a screenshot of a trace.

\subsection*{Data flow (ASCII)}

\begin{asciiart}
  bench 5V                 nRF PPK2                      Pi 4
  +---------+          +------------------+          +--------------------+
  |   5 V   |---VIN--->| source mode      |          | nRF Util /         |
  +---------+          |   set 3300 mV    |---USB--->|   Power Profiler   |
                       | 100 kS/s ammeter |   CDC    |   or ppk2-python   |
                       +------------------+          |        |           |
                              |    ^                 |        v           |
                            VOUT   | GND             | 2 s per state      |
                              |    |                 | median and p95     |
                              v    |                 +--------------------+
                       +------------------+
                       | ESP32 NodeMCU    |      USB port:  NOT CONNECTED
                       |   3V3 pin, GND   |      (the whole experiment)
                       | idle | BLE | deep|
                       +------------------+

  One power path, one host path, one ground.  A second power path makes
  every number in the run an unknown fraction of the truth.
\end{asciiart}

\subsection*{Steps}

\begin{steps}
\item \textbf{Install the host tools.} On the Pi, install Nordic's nRF Util with the Power
Profiler command line, or capture CDC samples with the published ppk2-python helpers.
Either path is acceptable; what matters is that the samples land in a file you can reduce
afterwards.

\item \textbf{Wire the rail, then check it.} Bench 5 V into VIN, VOUT to the ESP32 3V3 pin,
grounds common across the PPK2, the DUT and the Pi. Set source mode and 3300 mV before the
DUT is connected, not after.

\item \textbf{Disconnect the DUT USB and keep it disconnected.} The ESP32's own USB cable
stays out of the bench for the whole session. This is the whole experiment.

\item \textbf{Write three tiny sketches.} Radio-off idle, BLE advertise, and deep-sleep. Do
not measure the Arduino blink sketch and call it deep-sleep. Each sketch should do one
thing, so that the state under measurement is the state you named.

\item \textbf{Capture two seconds per state at 100 kS/s.} Let the state settle before the
capture window opens, then take the full two seconds without touching the bench.

\item \textbf{Reduce each capture to two numbers.} Record the median and the p95 for every
state, in a table with the three state names down the side. Quote those numbers afterwards,
not one screenshot of a trace.

\item \textbf{Repeat one state at the end of the run.} Go back to the first state and
capture it again. If the median has moved, something on the bench moved with it: a lead, a
ground, or the state itself. A run that cannot reproduce its own first number is not
finished, whatever the other two states reported.

\item \textbf{Compare the three states before you put the bench away.} Deep-sleep, idle and
BLE advertise should be three visibly different numbers. If they are not, go back to the
DUT USB cable before you go anywhere else.

\item \textbf{Write the honest sentence in the lab book.} Deep-sleep on a dev-board is tens
of microamps to low milliamps because the USB-UART silicon leaks. Write that sentence down
beside the number, so the figure is never quoted later as a bare-module deep-sleep current.
\end{steps}

\subsection*{Acceptance test}

\begin{itemize}
\item Deep-sleep measures tens of microamps to low milliamps on a dev-board, and the lab
book records why: the USB-UART silicon leaks.
\item Active radio measures tens of milliamps.
\item Each state has a median and a p95 written down, not a screenshot.
\item If the three traces overlay, the DUT is still USB-powered. That is a fail, and the
cause is the cable, not the instrument.
\end{itemize}

\subsection*{Practices}

Quote median and p95, not one screenshot. Never source SIM7600 transmit current from the
PPK2. The instrument is 5 V-class and about 1 A in source mode, which covers an ESP32 in
every state this lab measures and does not cover a Cat-4 power amplifier.

Keep this lab free of telemetry. Nothing publishes, nothing subscribes, and no broker is
involved: the numbers are the product. When another lab wants a power figure, it comes here
for it rather than measuring in passing, and it brings its own DUT wiring, because the one
rule that made these numbers valid is that the DUT had exactly one source of power.

Write the bench down as well as the numbers. Which supply fed VIN, what the source voltage
was set to, which sketch each state ran, and how long the state settled before the window
opened. Those four facts are what let somebody repeat the run next month and get the same
table, and they cost four lines in the lab book.

\subsection*{Pitfalls}

\begin{itemize}
\item \textbf{The DUT USB cable still connected.} The symptom is three traces that overlay
at a plausible current. The measurement is of the PPK2's share of a shared rail, which is
not a quantity anyone wants.
\item \textbf{Sourcing a modem from the PPK2.} A Cat-4 transmit peak can exceed about 1 A.
The symptom is a capture with a collapsed rail and a module that resets during transmission.
\item \textbf{Measuring the blink sketch as deep-sleep.} The symptom is a deep-sleep figure
in the milliamps that nobody can reproduce on a bare module.
\item \textbf{Two ground references.} Without a common ground the instrument measures a
loop rather than the DUT. The symptom is a reading that changes when you move a lead.
\item \textbf{Quoting a screenshot.} A trace picture hides the distribution. The symptom is
a power budget built from peaks that were never typical or from averages that were never
measured.
\item \textbf{Capturing before the state settles.} Boot current is not idle current. The
symptom is a median that drifts between repeats of the same state.
\item \textbf{Reporting a dev-board figure as a module figure.} The USB-UART silicon leaks
and the regulator has a quiescent current. Write the sentence in the lab book so the number
is never quoted without it.
\item \textbf{Setting the source voltage with the DUT attached.} Set source mode and
3300 mV first. The symptom is an ESP32 that sees whatever the instrument was last left at.
\item \textbf{Leaving the DUT USB cable coiled beside the board.} It goes back in between
runs, usually without anyone noticing. The symptom is one state in the table that does not
match the others and cannot be reproduced.
\item \textbf{Treating this bench as a telemetry lab.} There is no broker here and no
uplink. The symptom is a capture taken while the DUT is doing something the three sketches
never described, which measures a state nobody named.
\end{itemize}

\subsection*{Sources}

\begin{itemize}
\item Nordic Semiconductor Power Profiler Kit II product page and user guide,
\url{https://www.nordicsemi.com/Products/Development-hardware/Power-Profiler-Kit-2}
\item Nordic nRF Util, with the Power Profiler command line,
\url{https://www.nordicsemi.com/Products/Development-tools/nRF-Util}
\item ppk2-python, the published host helpers for the PPK2 CDC stream,
\url{https://github.com/IRNAS/ppk2-api-python}
\item Espressif ESP32 technical reference and the sleep-mode documentation,
\url{https://docs.espressif.com/projects/esp-idf/en/latest/esp32/api-reference/system/sleep_modes.html}
\item JOY-iT SBC-NodeMCU-ESP32 board documentation,
\url{https://joy-it.net/en/products/SBC-NodeMCU-ESP32}
\item JOY-iT SBC-POW-BB power rail board documentation,
\url{https://joy-it.net/en/products/SBC-POW-BB}
\end{itemize}


\part{Sensor shields on the Nucleo}
\project{07}{Consumer 9-DoF plus climate (IKS4A1)}{Raspberry Pi 4 or 3B+ + Nucleo-H7A3ZI-Q + X-NUCLEO-IKS4A1}{consumer MEMS, humidity and the Qvar electrode}

\begin{unique}
Owns consumer MEMS plus SHT40 humidity and the Qvar electrode. Wearable physics, not
industrial high-g.
\end{unique}

\subsection*{Intent}

The IKS4A1 on the Nucleo streams JSON of the LSM6DSV16X and LSM6DSO16IS inertial parts,
the LIS2MDL magnetometer, the LIS2DUXS12 accelerometer, the LPS22DF barometer, the SHT40
humidity sensor and the STTS22H temperature sensor. Linux records. That is the whole
product: a line of JSON per sample period, written to a file that is honest about what was
measured and when.

The teaching point of this lab is a boundary, not a sensor. Every one of those chips sits
on the STM32 I2C bus, not on the Pi I2C bus, so \texttt{i2cdetect} on the Pi finds nothing
at all. Newcomers read that empty grid as a fault and start pulling jumpers. It is not a
fault, it is the architecture: the microcontroller owns the sensor bus, and Linux owns
storage, time and the file format. The only thing that crosses the USB cable is text the
firmware chose to print. P08 uses the same physical stack with the industrial shield, and
P18 reuses this shield when the M7 runs a filter, but the boundary is identical in all
three.

\begin{kit}
Pi 4 or 3B+, Nucleo-H7A3ZI-Q, X-NUCLEO-IKS4A1, USB cable.
\end{kit}

\subsection*{System architecture}

\diagram{p07_arch}{The boundary that defines this lab. Seven sensor parts sit on the STM32
I2C bus behind the microcontroller. The Pi I2C controller is idle and unwired, which is
why \texttt{i2cdetect -y 1} on the Pi returns an empty grid.}

Read the figure from the left. The shield chips answer at 0x6A for the LSM6 inertial part,
0x1E for the LIS2MDL, 0x44 for the SHT40, 0x5C for the LPS22DF and 0x38 for the STTS22H.
Those addresses live on the microcontroller bus. The scanner that proves them is the one
in the firmware, not the one on the Pi. The Pi I2C controller is drawn deliberately as an
empty box: nothing is wired to pins 3 and 5, and nothing should be.

On the Linux side the work is small and should stay small. A CDC ACM character device
delivers lines, and \texttt{tee} writes them to a file while you watch them go past. Any
processing beyond that, including the tilt computation, is a consumer of the file or of
the same stream, never a replacement for it. Keep the raw capture; a derived quantity you
cannot recompute from the record is a claim rather than a measurement.

\subsection*{Wiring and schematic}

\diagram{p07_schematic}{The IKS4A1 on the Arduino UNO headers, with the UM3239 Mode 1
jumper positions that put every sensor on the microcontroller I2C bus. The Pi side of the
drawing carries one USB connector and nothing else.}

\begin{table}[H]\centering\small
\begin{tabular}{p{40mm}p{28mm}p{22mm}p{58mm}}
\toprule
Signal or part & Connector & Position & Note \\
\midrule
X-NUCLEO-IKS4A1 & Nucleo Arduino UNO headers & stacked & One shield only, IKS5A1 stays bagged \\
UM3239 Mode 1, J4 & IKS4A1 jumper J4 & 1-2 & All sensors on the uC I2C bus \\
UM3239 Mode 1, J4 & IKS4A1 jumper J4 & 11-12 & Second position of the same mode \\
UM3239 Mode 1, J5 & IKS4A1 jumper J5 & 1-2 & All sensors on the uC I2C bus \\
UM3239 Mode 1, J5 & IKS4A1 jumper J5 & 11-12 & Second position of the same mode \\
LSM6DSV16X or LSM6DSO16IS & STM32 I2C bus & 0x6A & Inertial, behind the STM32 \\
LIS2MDL & STM32 I2C bus & 0x1E & Magnetometer \\
SHT40AD1B & STM32 I2C bus & 0x44 & Relative humidity \\
LPS22DF & STM32 I2C bus & 0x5C & Barometer \\
STTS22H & STM32 I2C bus & 0x38 & Temperature \\
LIS2DUXS12 & STM32 I2C bus & see the shield manual & The source does not fix this address, confirm in UM3239 and on the silkscreen \\
Qvar electrode & IKS4A1 jumpers & UM3239 Qvar mode & Left off until the IMU stream is clean \\
Nucleo USB & Pi USB-A & any & CDC ACM, the only link to Linux \\
Pi 40-pin header & Pi & empty & No HAT, no flying leads, nothing on pins 3 and 5 \\
\bottomrule
\end{tabular}
\caption{Wiring. Mode 1 is two jumper positions on each of J4 and J5. Nothing in this lab
connects to the Pi header.}
\end{table}

\begin{asciiart}
IKS4A1 on Nucleo Arduino headers.
UM3239 Mode 1 (all sensors on uC I2C):
  J4: 1-2 and 11-12
  J5: 1-2 and 11-12
Nucleo USB --> Pi.
sudo i2cdetect on the *Pi* will NOT see these chips.
They live behind the STM32. That is the point.
\end{asciiart}

Mode 1 is the Linux-friendly setup because it puts every part on one bus that the firmware
owns and can scan in one pass. The other modes on this shield route parts through the
sensor hub or the ISPU inside the LSM6, which hides devices from that scan. Those modes are
worth a session after the JSON is honest, not before.

\subsection*{Bench layout}

\diagram{p07_bench}{Bench layout. The shield is stacked on the Nucleo, one USB cable
reaches the Pi, and the Pi header stays empty. The two stimuli that prove the climate
sensors are a breath toward the SHT40 and a palm held over the board.}

Put the Nucleo where you can reach the shield with a hand and a breath without moving the
cable. Keep the stack away from a laptop fan and away from a mug: both move relative
humidity and temperature on their own and make the acceptance test ambiguous. The Pi can
sit anywhere the USB cable reaches, because it wears no HAT in this lab and has nothing
plugged into its header.

\subsection*{Software design (UML)}

\diagram{p07_uml}{Component view. The seven sensor components sit behind the STM32 I2C
master. The firmware exposes one interface to the outside world, a CDC ACM port carrying
JSON lines, and the Linux recorder consumes exactly that. There is no interface from the
Pi to the sensor bus, which is why the Pi scanner is empty.}

The component diagram is the argument of this chapter drawn once. Every sensor is a
component with a required I2C interface, and the only component that provides that
interface is the STM32 I2C master. The firmware provides a single external interface: JSON
lines over CDC. The Linux recorder requires that interface and provides a file. Nothing in
the picture connects a Linux component to a sensor component, and no wiring change in this
lab can create such a connection.

\subsection*{Data flow (ASCII)}

\begin{asciiart}
  X-NUCLEO-IKS4A1                       Nucleo-H7A3ZI-Q          Pi 4 or 3B+
  +-----------------------+             +-----------------+      +-----------------+
  | LSM6DSV16X / DSO16IS  |--0x6A--\    |                 |      |                 |
  | LIS2MDL               |--0x1E---\   | STM32 I2C       |      | /dev/ttyACM0    |
  | SHT40AD1B             |--0x44----+->| master + driver |      |   |             |
  | LPS22DF               |--0x5C---/   |     |           | USB  |   v             |
  | STTS22H               |--0x38--/    |     v           |=====>| cat | tee       |
  | LIS2DUXS12            |--see UM3239 | one JSON object |  CDC |   p07.jsonl     |
  | Qvar electrode (off)  |             | per 20 ms       |      |   |             |
  +-----------------------+             +-----------------+      |   v             |
                                                                 | tilt from acc   |
                                                                 | and gyro only   |
                                                                 +-----------------+

  Pi i2c-1:  nothing wired, i2cdetect -y 1 is empty.  That is the architecture,
             not a fault.  The scanner that proves the addresses is in the firmware.
\end{asciiart}

\subsection*{Steps}

\begin{steps}
\item \textbf{Stack the shield and set Mode 1.} The IKS4A1 goes on the Nucleo Arduino
headers. Set the UM3239 Mode 1 positions, which put all sensors on the microcontroller I2C
bus:
\begin{plaincode}
J4: 1-2 and 11-12
J5: 1-2 and 11-12
\end{plaincode}
The IKS5A1 stays in its own bag. The two shields look similar from across the bench, and
mixing them is how the wrong chapter gets built.

\item \textbf{Flash the firmware.} Use an STM32duino X-NUCLEO-IKS4A1 sample, or a sketch
that prints one JSON object per 20 ms:
\begin{plaincode}
{"t":24.1,"rh":41.2,"p":1013.2,"ax":0.02,"ay":0.01,"az":1.00}
\end{plaincode}
One object per line, one line per sample period. Keep the key names stable from the first
capture, because a schema that drifts halfway through a file is worse than a short file.

\item \textbf{Record on the Pi.} Connect the Nucleo USB cable to a Pi USB-A port and write
the stream to a file while you watch it.
\begin{shellcode}
cat /dev/ttyACM0 | tee p07.jsonl
\end{shellcode}

\item \textbf{Prove the boundary.} Run the Pi scanner once, so the empty grid is something
you saw rather than something you were told.
\begin{shellcode}
sudo i2cdetect -y 1
\end{shellcode}
It shows nothing, and it is meant to. These chips live behind the STM32. The addresses
0x6A, 0x1E, 0x44, 0x5C and the 0x38-class temperature part are visible only to the scanner
inside the firmware.

\item \textbf{Compute a running tilt from acc and gyro only.} Magnetometer heading waits on
a hard-iron calibration you actually perform. Do not publish a raw heading and call it
fusion. A tilt that comes from two sensors you trust is worth more than a heading that
comes from three you have not calibrated.

\item \textbf{Test the climate parts with the two stimuli.} Breathe toward the SHT40 and
watch relative humidity move. Hold a palm over the board and watch the STTS22H move by
about 0.5 degrees C. Do them separately: a breath carries both moisture and heat, so doing
both at once tells you less than either alone.

\item \textbf{Leave Qvar for a second session.} The Qvar swipe stays off until the IMU
stream is clean. Then enable the electrode jumpers, in the UM3239 Qvar mode, as its own
session with its own capture file, not as a kitchen sink on top of a stream you are still
validating.
\end{steps}

\subsection*{Acceptance test}

\begin{itemize}
\item The firmware scanner sees 0x6A, 0x1E, 0x44, 0x5C and the 0x38-class addresses.
\item \texttt{sudo i2cdetect -y 1} on the Pi shows an empty grid, and that is a pass.
\item Breathing on the SHT40 moves relative humidity.
\item A palm over the board moves the STTS22H by about 0.5 degrees C.
\item \texttt{p07.jsonl} holds one JSON object per line with the same keys from the first
line to the last.
\end{itemize}

\subsection*{Practices}

Mode 1 is the Linux-friendly setup. Sensor-hub and ISPU modes hide devices behind the
LSM6, which is useful after the JSON is honest and confusing before it. Keep this lab's
numbers in their own log file: P08 measures a different instrument and its schema is not
this one.

Record first and derive later. The file is the artifact, the tilt is a view of it, and a
heading you have not calibrated is neither. When you do run the hard-iron calibration,
record it as its own capture with its own note in the lab book, so the correction can be
recomputed rather than remembered.

\subsection*{Pitfalls}

\begin{itemize}
\item \textbf{Running \texttt{i2cdetect} on the Pi and reading the empty grid as a fault.}
The chips are behind the STM32. The symptom is an hour spent checking jumpers that were
correct, followed by a Mode change that actually does hide devices.
\item \textbf{Choosing a sensor-hub or ISPU mode by accident.} Those modes route parts
through the LSM6, so the firmware scan finds fewer addresses than the shield carries. The
symptom is a JSON object that silently loses keys.
\item \textbf{Publishing a raw magnetometer heading as fusion.} Without a hard-iron
calibration the heading follows the nearest steel object. The symptom is a compass that is
stable and wrong.
\item \textbf{Mixing this schema with P08.} The industrial shield produces different
quantities from different silicon. Two schemas in one file is how a wearable-grade number
ends up labelled industrial.
\item \textbf{Stacking both shields.} The Arduino stack takes one shield. The symptom is an
address clash at 0x6A and a scan that reports the wrong part.
\item \textbf{Enabling Qvar before the IMU stream is clean.} Two new variables at once
means neither is isolated when something looks wrong.
\item \textbf{Breathing on the board to test temperature.} A breath moves humidity and
temperature together. Use the palm for STTS22H and the breath for SHT40, one at a time.
\item \textbf{A schema that drifts.} Adding a key halfway through a capture leaves a file
that no parser reads end to end. Restart the capture instead.
\end{itemize}

\subsection*{Sources}

\begin{itemize}
\item ST X-NUCLEO-IKS4A1 expansion board page,
\url{https://www.st.com/en/ecosystems/x-nucleo-iks4a1.html}
\item ST UM3239, the X-NUCLEO-IKS4A1 user manual, for the J4 and J5 mode tables,
\url{https://www.st.com/resource/en/user_manual/um3239.pdf}
\item STM32duino X-NUCLEO-IKS4A1 library and samples,
\url{https://github.com/stm32duino/X-NUCLEO-IKS4A1}
\item ST LSM6DSV16X and LSM6DSO16IS data sheets,
\url{https://www.st.com/en/mems-and-sensors/inemo-inertial-modules.html}
\item ST LPS22DF pressure sensor and STTS22H temperature sensor data sheets,
\url{https://www.st.com/en/mems-and-sensors/pressure-sensors.html}
\item Sensirion SHT40 humidity and temperature sensor data sheet,
\url{https://sensirion.com/products/catalog/SHT40}
\item ST NUCLEO-H7A3ZI-Q board page,
\url{https://www.st.com/en/evaluation-tools/nucleo-h7a3zi-q.html}
\end{itemize}

\project{08}{Industrial high-g and dual-scale baro (IKS5A1)}{Raspberry Pi 3B+ + Nucleo-H7A3ZI-Q + X-NUCLEO-IKS5A1}{simultaneous low-g and high-g, and the 4060 hPa scale}

\begin{unique}
Owns ISM6HG256X simultaneous low-g and high-g, and ILPS22QS 1260/4060 hPa. No humidity.
A different instrument from P07.
\end{unique}

\begin{blueprint}
Draft P13 put an ISM330DHCX Machine Learning Core on the IKS5A1. That IMU is on the
STWIN.box. IKS5A1 has ISM330IS (ISPU) and ISM6HG256X. Use the chips on the board.
\end{blueprint}

\subsection*{Intent}

Two concurrent streams: a low-g IMU at 104 Hz, and a high-g peak-hold. The barometer runs
in the 4060 hPa scale only when you actually pressurise something, and in its ordinary
scale the rest of the time. That pairing is what the ISM6HG256X offers and what the
IKS4A1 in P07 cannot do: a shock that saturates the low-g channel is still a number on the
high-g channel, in the same event, from the same silicon.

The physical build is the P07 build with the other shield, so this chapter does not repeat
it. Stack the IKS5A1 where the IKS4A1 was, with the IKS4A1 removed and bagged, and read
P07 for the Arduino stack, the single USB cable and the reason the Pi scanner stays empty.
What is new here is the data, and the data has a rule attached: two streams, two keys, two
log files. Collapsing them into one acceleration field is how a wearable-grade number ends
up wearing an industrial label.

\begin{kit}
Pi 3B+, Nucleo-H7A3ZI-Q, X-NUCLEO-IKS5A1 with the IKS4A1 removed, USB cable.
\end{kit}

\subsection*{System architecture}

\diagram{p08_arch}{Five parts on the microcontroller I2C bus feed two streams that never
merge. The low-g channel runs at 104 Hz; the high-g channel reports a peak-hold. Linux
records them into two files with two schemas.}

The sensor side is five chips at five addresses on the STM32 bus, exactly as in P07, and
the Pi again sees none of them. The difference begins at the sampler. The ISM6HG256X
provides a low-g channel and a high-g channel at the same time, so the firmware carries
two producers rather than one, with different rates and different meanings. The low-g
producer emits at 104 Hz. The high-g producer emits a peak-hold, which is an event
quantity: the largest magnitude seen since the last report.

Linux keeps that separation all the way to disk. Two JSON keys, \texttt{lowg} and
\texttt{highg\_peak}, or two MQTT topics, and two files. A reader that opens either one
knows what it holds without asking which part of which event the numbers came from. The
ILPS22QS sits beside both with its own quantity and its own scale, and the scale is part
of the record: a pressure without the range it was measured in is a number with two
possible meanings.

\subsection*{Wiring and schematic}

\diagram{p08_schematic}{The IKS5A1 on the Nucleo Arduino headers, its five default I2C
addresses, and the one USB cable to the Pi 3B+. The stack itself is the P07 stack: see
that chapter for the header detail.}

\begin{table}[H]\centering\small
\begin{tabular}{p{38mm}p{30mm}p{20mm}p{60mm}}
\toprule
Part or signal & Connector & Address & Note \\
\midrule
ISM6HG256X & STM32 I2C bus & 0x6A & Low-g and high-g at the same time \\
ISM330IS & STM32 I2C bus & 0x6B & The ISPU part on this shield \\
ILPS22QS & STM32 I2C bus & 0x5C & 1260 and 4060 hPa scales \\
IIS2MDC & STM32 I2C bus & 0x1E & Magnetometer \\
IIS2DULPX & STM32 I2C bus & 0x19 & Low-power accelerometer \\
X-NUCLEO-IKS5A1 & Nucleo Arduino UNO headers & stacked & The P07 stack with the other shield \\
X-NUCLEO-IKS4A1 & off the board & none & Removed and bagged before this lab \\
Shield mode jumpers & IKS5A1 & see the manual & The source fixes the addresses, not the jumper table: confirm in the shield user manual and on the silkscreen \\
Nucleo USB & Pi 3B+ USB-A & any & CDC ACM, the only link to Linux \\
Pi 40-pin header & Pi 3B+ & empty & No HAT, nothing wired to the sensor bus \\
\bottomrule
\end{tabular}
\caption{Wiring. Five parts, five addresses, one shield. The ISM330DHCX is not among them:
that IMU is on the STWIN.box in P04.}
\end{table}

\begin{asciiart}
The same physical stack as P07, with the other shield.

IKS5A1 on Nucleo Arduino headers   (IKS4A1 removed and bagged)
Default I2C, all behind the STM32:
  ISM6HG256X ---- 0x6A     low-g and high-g together
  ISM330IS ------ 0x6B     ISPU
  ILPS22QS ------ 0x5C     1260 / 4060 hPa
  IIS2MDC ------- 0x1E     magnetometer
  IIS2DULPX ----- 0x19     low-power accelerometer
Nucleo USB --> Pi 3B+      CDC ACM
NOT on this board: ISM330DHCX.  That IMU is on the STWIN.box (P04).
\end{asciiart}

Everything about the stack itself belongs to P07: the Arduino UNO headers, the single USB
cable, and the empty Pi header. Read it there once. What this chapter adds to the wiring is
a list of five addresses and one subtraction: the ISM330DHCX that the blueprint draft put
here is not on this board and cannot be scanned into existence.

\subsection*{Bench layout}

\diagram{p08_bench}{Bench layout. The drop test is 20 mm onto a book. The book is the
point: it gives a repeatable shock that the high-g channel can report and that the bench,
the shield and the connector all survive.}

Drop from 20 mm onto a book, not onto concrete. A concrete drop produces one impressive
number, a damaged connector and no second measurement. The book gives a shock you can
repeat ten times in a row, which is what turns a peak into a distribution. Keep the USB
cable slack during the drop so the cable does not lift the board or pull the plug, and put
the whole stack back on the bench between drops rather than measuring it in flight.

\begin{note}{Two shields that look alike}
Put the IKS4A1 and the IKS5A1 back in separate bags every time. They carry different
silicon, different addresses and different claims, and they are the same shape. The
teardown line in the source says it plainly: they look similar from across the bench. A
label on each bag costs nothing and settles the question at the start of the next session
rather than at the end of it.
\end{note}

\subsection*{Software design (UML)}

\diagram{p08_uml}{Two streams, two schemas, two files. The low-g stream and the high-g
peak-hold share silicon and share a timestamp source, and share nothing else. There is no
component that merges them, because merging them is the error this chapter exists to
prevent.}

The diagram shows the rule as structure rather than as advice. Two producers in the
firmware, two keys on the wire, two writers on the Linux side, two files on disk. An
analysis that wants both opens both and joins on time, which is an explicit operation
somebody has to write down. What it cannot do is find them already merged in a field
called \texttt{acc} and assume the merge was correct.

\subsection*{Data flow (ASCII)}

\begin{asciiart}
  X-NUCLEO-IKS5A1                 Nucleo-H7A3ZI-Q            Pi 3B+
  +---------------------+         +---------------------+    +----------------------+
  | ISM6HG256X   0x6A   |         | low-g producer      |    | reader on ttyACM0    |
  |   low-g channel ----------->  |   104 Hz            |--\ |   |                  |
  |   high-g channel ---------->  | high-g producer     |  | |   +--> p08_lowg.jsonl|
  | ISM330IS     0x6B   |         |   peak-hold, event  |--+-+->  |    key: lowg     |
  | ILPS22QS     0x5C   |         | baro producer       |  | |   |                  |
  | IIS2MDC      0x1E   |         |   scale in the record| | |   +--> p08_highg.jsonl|
  | IIS2DULPX    0x19   |         +---------------------+  | |        key: highg_peak|
  +---------------------+            |   USB CDC ACM       | +----------------------+
                                     +---------------------+

  Never collapse lowg and highg_peak into one "acc" field.
  These numbers live in a different log file from P07.
\end{asciiart}

\subsection*{Steps}

\begin{steps}
\item \textbf{Swap the shield, one at a time.} Remove the IKS4A1 and put it in its own bag
before the IKS5A1 comes out of its bag. The two boards look similar from across the bench,
and a lab that starts with the wrong shield produces P07 numbers under a P08 heading.

\item \textbf{Build the stack as in P07.} Arduino UNO headers, one shield, one USB cable
from the Nucleo to a Pi 3B+ USB-A port, and nothing on the Pi header. The jumper table for
this shield is in its own user manual: confirm the mode positions there and on the
silkscreen before you flash anything.

\item \textbf{Flash the samples.} Use the STM32duino X-NUCLEO-IKS5A1 samples for the
ISM6HG256X 6D and wake-up functions and for ILPS22QS pressure. Confirm in the firmware scan
that the addresses answer:
\begin{plaincode}
0x6A ISM6HG256X   0x6B ISM330IS   0x5C ILPS22QS   0x1E IIS2MDC   0x19 IIS2DULPX
\end{plaincode}

\item \textbf{Publish two topics or two JSON keys.} The names are \texttt{lowg} and
\texttt{highg\_peak}. Never collapse them into one \texttt{acc} field.
\begin{plaincode}
{"lowg":{"ax":0.01,"ay":0.02,"az":1.00}}
{"highg_peak":{"mag":47.2}}
\end{plaincode}
One object per line, one stream per key. If you use MQTT instead, the same rule applies to
the topic names.

\item \textbf{Record into two files.} The low-g stream and the high-g stream go to separate
files, and both go somewhere that is not the P07 capture. A single directory per lab, named
for the lab, is enough discipline to keep the two instruments apart.

\item \textbf{Run the drop-test protocol.} Twenty millimetres onto a book, not onto
concrete. Repeat the drop enough times to see the spread in the peak, and write the drop
height in the lab book beside the numbers, because a peak without a height is not a
measurement of anything.

\item \textbf{Measure the barometer honestly.} Leave the ILPS22QS in its ordinary scale for
room pressure and note the noise you see in a quiet room. Switch to the 4060 hPa scale only
when you actually pressurise something, and record the scale alongside the reading, because
the same number means two different things in the two ranges.
\end{steps}

\subsection*{Acceptance test}

\begin{itemize}
\item On the drop, the high-g peak fires and the low-g channel saturates. Both facts appear
in the record, in their own keys.
\item Quiet-room ILPS22QS noise is a few Pa.
\item The numbers live in a different log file from P07.
\item The firmware scan answers at 0x6A, 0x6B, 0x5C, 0x1E and 0x19, and reports no
ISM330DHCX, because there is none on this board.
\end{itemize}

\subsection*{Practices}

Mixing the P07 and P08 JSON schemas is how fake industrial wearables get invented. Do not.
Keep the shields in separate bags, keep the captures in separate files, and keep the two
chapters' claims apart: consumer MEMS with humidity is P07, simultaneous low-g and high-g
with a dual-scale barometer is this lab.

Record the scale, the drop height and the surface with every number. A peak-hold is an
event quantity, so it is only comparable to another peak-hold taken the same way. When the
high-g channel and the low-g channel disagree about an impact, that is the instrument
working, not a fault to be smoothed away.

\subsection*{Pitfalls}

\begin{itemize}
\item \textbf{Looking for an ISM330DHCX.} It is on the STWIN.box, which is P04. The symptom
is a scan that never finds the address and a sample that never compiles against this
shield.
\item \textbf{Collapsing \texttt{lowg} and \texttt{highg\_peak} into one field.} The
symptom is a file in which a saturated low-g sample and a valid high-g peak are
indistinguishable, and every later analysis inherits the confusion.
\item \textbf{Writing P08 numbers into the P07 capture.} Two instruments, one file. The
symptom is a data set that claims humidity and high-g from the same silicon.
\item \textbf{Leaving the IKS4A1 on the bench next to the IKS5A1.} They look similar. The
symptom is a lab that produces the wrong chapter's addresses and a confident, wrong entry
in the lab book.
\item \textbf{Dropping onto concrete.} One number, then a damaged connector. The symptom is
a peak you cannot reproduce because the hardware no longer enumerates.
\item \textbf{Quoting a pressure without its scale.} The 1260 and 4060 hPa ranges are
different instruments for this purpose. The symptom is a reading that cannot be compared
with yesterday's.
\item \textbf{Expecting \texttt{i2cdetect} on the Pi to answer.} As in P07, these chips are
behind the STM32. An empty grid on the Pi is the architecture working.
\item \textbf{Treating the peak-hold as a rate.} It is the largest magnitude since the last
report, not a sample stream. Averaging peaks produces a number with no physical meaning.
\end{itemize}

\subsection*{Sources}

\begin{itemize}
\item ST X-NUCLEO-IKS5A1 expansion board page,
\url{https://www.st.com/en/ecosystems/x-nucleo-iks5a1.html}
\item STM32duino X-NUCLEO-IKS5A1 library and samples,
\url{https://github.com/stm32duino/X-NUCLEO-IKS5A1}
\item ST ISM6HG256X high-g and low-g inertial module data sheet,
\url{https://www.st.com/en/mems-and-sensors/ism6hg256x.html}
\item ST ISM330IS inertial module with ISPU data sheet,
\url{https://www.st.com/en/mems-and-sensors/ism330is.html}
\item ST ILPS22QS dual full-scale pressure sensor data sheet,
\url{https://www.st.com/en/mems-and-sensors/ilps22qs.html}
\item ST IIS2MDC magnetometer and IIS2DULPX accelerometer data sheets,
\url{https://www.st.com/en/mems-and-sensors/accelerometers.html}
\item ST NUCLEO-H7A3ZI-Q board page,
\url{https://www.st.com/en/evaluation-tools/nucleo-h7a3zi-q.html}
\end{itemize}

\project{18}{Nucleo USB-CDC recorder}{Raspberry Pi 4 + bare Nucleo-H7A3ZI-Q, optional X-NUCLEO-IKS4A1}{high-rate USB gadget traffic}

\begin{unique}
Owns high-rate USB gadget traffic from the bare Nucleo. Linux is the recorder. An EKF may
run on the H7 FPU; it must not replace the Linux host.
\end{unique}

\begin{blueprint}
Draft P15 discards Linux entirely. That is out of scope for an embedded-Linux book unless
a host still captures the UART or CDC stream. We keep the host.
\end{blueprint}

\subsection*{Intent}

The Nucleo-H7A3ZI-Q streams a packed binary or text sample at 1 kHz or more over USB-CDC.
The Pi 4 writes a file with sequence numbers. That is the entire product, and it is a
harder product than it looks: a thousand samples a second through a character device, for
ten seconds, with a counter that lets you prove afterwards how many of them arrived.

The optional half of the lab is a quaternion EKF on the M7 floating-point unit, using the
P07 shield. It is optional because it changes nothing about the architecture. The filter
runs on the microcontroller, the host still records the stream, and the acceptance test is
still about sequence gaps rather than about the elegance of the filter. A bare-metal sketch
is allowed in this book exactly when a Linux host still records what it produces, which is
what this chapter demonstrates.

\begin{kit}
Pi 4, Nucleo-H7A3ZI-Q, USB cable, optional X-NUCLEO-IKS4A1, optional Renkforce USB-TTL as a
second console.
\end{kit}

\subsection*{System architecture}

\diagram{p18_arch}{The recorder. A sequence generator on the M7 produces samples at 1 kHz,
the USB-CDC class carries them, and a Python sink on the Pi counts gaps before it writes.
The shield and the EKF are drawn as the option they are.}

There are only two moving parts, and the interesting one is the counter. The firmware
increments a \texttt{uint32} sequence number and emits it with every sample. The sink on
the Pi keeps the previous value and compares. A gap means samples went missing somewhere
between the two, and the count of gaps over a known interval is the quality number for the
whole path: firmware timing, USB scheduling, the character device and the reader loop, all
in one figure you can quote.

The Linux side must be told not to interpret the stream. A terminal device applies a line
discipline by default, which edits characters on their way through and echoes them back.
The \texttt{stty} call that sets the port to raw at 921600 baud is not decoration: without
it, the capture is a record of what the terminal driver decided to do with the data rather
than a record of what the microcontroller sent.

\subsection*{Wiring and schematic}

\diagram{p18_schematic}{One cable is the whole mandatory wiring. The shield and the USB-TTL
console are optional, drawn dashed, and neither of them changes the recorder path.}

\begin{table}[H]\centering\small
\begin{tabular}{p{38mm}p{28mm}p{24mm}p{58mm}}
\toprule
Signal or part & Connector & Setting & Note \\
\midrule
Nucleo USB & Pi 4 USB-A & CDC ACM & The mandatory link, the recorder path \\
Character device & Pi & /dev/ttyACM0 & More than one ACM node can appear \\
Line settings & Pi & 921600 raw & Set with \texttt{stty} before reading, see the steps \\
Sample rate & firmware & 1 kHz or more & A \texttt{uint32} sequence with every sample \\
X-NUCLEO-IKS4A1 & Nucleo Arduino UNO headers & optional & Only for the EKF session, the P07 shield with the P07 jumper mode \\
Renkforce USB-TTL & Pi USB-A & optional & 3.3 V logic, a second console only, never a second recorder path \\
Nucleo and Pi I/O & both & 3.3 V & Confirm the console pins on the Nucleo silkscreen before connecting the cable \\
Pi 40-pin header & Pi 4 & empty & No HAT in this lab \\
\bottomrule
\end{tabular}
\caption{Wiring. The recorder needs one USB cable. Everything else in the table is an
option that the acceptance test does not depend on.}
\end{table}

\begin{asciiart}
Mandatory:
  Nucleo USER USB  -->  Pi 4 USB-A        CDC ACM, /dev/ttyACM0
  stty -F /dev/ttyACM0 921600 raw
  firmware: uint32 sequence, 1 kHz or more

Optional, and only after the recorder passes:
  X-NUCLEO-IKS4A1 on the Nucleo Arduino headers   (the P07 shield, P07 mode)
  quaternion EKF on the M7 FPU, printing seq,q0,q1,q2,q3
  Renkforce USB-TTL as a second console, 3.3 V logic

Pi 40-pin header: empty.  Linux does not run the filter.
\end{asciiart}

The optional shield is the P07 shield and nothing about it changes here, so set it up as
P07 describes and come back. The Renkforce cable is the P14 jig's tool used as a console;
it does not become a second data path, because two recorders of one stream produce two
files that disagree and no way to say which is right.

\subsection*{Bench layout}

\diagram{p18_bench}{Bench layout. A short USB cable, a Pi with an empty header, and nothing
else competing for the bus during a capture. The optional shield and console are drawn
where they would go.}

Use a short cable and leave the other USB ports alone while a capture runs. A ten-second
capture at 1 kHz is ten thousand samples, and anything that makes the host busy at the
wrong moment shows up as a gap that has nothing to do with the firmware. Strain-relieve the
cable so a knock on the bench does not re-enumerate the device in the middle of the run.

\subsection*{Software design (UML)}

\diagram{p18_uml}{The 1 kHz stream and the gap count. The firmware emits sequence numbers
without waiting for anyone. The sink reads, compares each sequence number with the previous
one, counts what is missing, and writes. The gap count is the acceptance number.}

The sink is a loop with two pieces of state: the previous sequence number and a gap
counter. Every line either continues the sequence or does not, and when it does not, the
difference is added to the counter rather than being repaired. Nothing is interpolated and
nothing is retried, because the point of the exercise is to measure loss, not to hide it.
At the end of ten seconds the loop prints the number of samples it saw and the number it
did not, and those two numbers are the result of the lab.

\subsection*{Data flow (ASCII)}

\begin{asciiart}
  Nucleo-H7A3ZI-Q  (bare, optional shield)       Pi 4  (the recorder)
  +-------------------------------+              +------------------------------+
  | uint32 seq++                  |              | stty 921600 raw              |
  |   1 kHz or more               |              |   |                          |
  |   |                           |    USB       |   v                          |
  |   v                           |=============>| /dev/ttyACM0                 |
  | USB-CDC class                 |   CDC ACM    |   |                          |
  +-------------------------------+              |   v                          |
        ^                                        | p18_sink.py                  |
        |  optional session only                 |   seq - prev == 1 ?          |
  +-------------------------------+              |     no  -> gaps += diff - 1  |
  | X-NUCLEO-IKS4A1 (P07 shield)  |              |   write the line             |
  | quaternion EKF on the M7 FPU  |              |   |                          |
  | prints seq,q0,q1,q2,q3        |              |   v                          |
  +-------------------------------+              | capture file + gap count     |
                                                 +------------------------------+

  Linux does not run the filter.  Linux records, counts and keeps the file.
\end{asciiart}

\subsection*{Steps}

\begin{steps}
\item \textbf{Write the firmware.} A Cube or Arduino USB-CDC sketch that increments a
\texttt{uint32} sequence and prints it. Aim for 1 kHz. Keep the formatting cheap: a
thousand lines a second leaves little room for an expensive print, and the first thing to
check when the rate will not hold is how much work happens between two samples.

\item \textbf{Confirm the device enumerated as CDC.}
\begin{shellcode}
lsusb
ls /dev/ttyACM*
\end{shellcode}
The CDC interface must appear in \texttt{lsusb}. If more than one \texttt{ttyACM} node
exists, the one carrying sequence numbers is the one to record.

\item \textbf{Set the port raw and start the sink.}
\begin{shellcode}
stty -F /dev/ttyACM0 921600 raw
python3 p18_sink.py          # counts gaps in the sequence
\end{shellcode}

\item \textbf{Write the sink so that it counts rather than repairs.} Two pieces of state,
one comparison, no interpolation.
\begin{pycode}
import sys, time, serial
ser = serial.Serial("/dev/ttyACM0", 921600, timeout=1)
prev = None
seen = 0
gaps = 0
t_end = time.monotonic() + 10.0
with open("p18_capture.txt", "w") as out:
    while time.monotonic() < t_end:
        line = ser.readline().decode(errors="ignore").strip()
        if not line:
            continue
        try:
            seq = int(line.split(",")[0])
        except ValueError:
            continue
        if prev is not None and seq != prev + 1:
            gaps += seq - prev - 1
        prev = seq
        seen += 1
        out.write(line + "\n")
print("samples", seen, "gaps", gaps)
\end{pycode}

\item \textbf{Run a ten-second capture and quote the two numbers.} Samples seen and samples
missing. A capture at 1 kHz that reports 0.1 percent or fewer sequence gaps passes. Run it
more than once: a single clean run on a quiet host is a weaker claim than three runs with
the numbers written down.

\item \textbf{Optional EKF session.} Reuse the P07 shield, compute a quaternion on the M7
floating-point unit, and print \texttt{seq,q0,q1,q2,q3}. Linux does not run the filter.
Keep the sequence number in the first field so the same sink still counts gaps, and keep
this capture in its own file: it is a different stream from the bare recorder run.
\end{steps}

\subsection*{Acceptance test}

\begin{itemize}
\item A 10-second capture at 1 kHz has 0.1 percent or fewer sequence gaps.
\item \texttt{lsusb} shows the CDC interface.
\item The sink prints both numbers, samples seen and samples missing, and the capture file
holds the lines it counted.
\item In the optional session, the quaternion is computed on the M7 and Linux still writes
the file.
\end{itemize}

\subsection*{Practices}

Linux owns the system. A bare-metal sketch is allowed in this book only when a Linux host
still records the stream, which is the whole reason this chapter exists. The filter, if you
run one, belongs on the M7; the file, the clock and the gap count belong to the host.

Count gaps rather than trusting the stream. A recorder that cannot say how much it lost is
not a recorder, it is a listener. Keep the capture short and repeat it: ten seconds at
1 kHz is enough to expose a scheduling problem, and three runs of ten seconds say more than
one run of a minute. Leave the other USB ports alone while a capture runs, and keep the
optional console as a console.

\subsection*{Pitfalls}

\begin{itemize}
\item \textbf{Forgetting \texttt{stty ... raw}.} The line discipline edits and echoes the
stream. The symptom is a capture that looks plausible, loses characters at the boundaries
and does not reproduce.
\item \textbf{Recording the wrong \texttt{ttyACM} node.} With the ST-LINK present, more
than one appears. The wrong one is silent rather than an error.
\item \textbf{An expensive print in the firmware loop.} Formatting costs time that the
1 kHz budget does not have. The symptom is a rate that sits well under 1 kHz with no
dropped samples at all, which is a different fault from a gap.
\item \textbf{Interpolating over a gap.} A sink that fills in a missing sample discards the
only measurement the lab produces. Count, do not repair.
\item \textbf{Ignoring the \texttt{uint32} wrap.} After about four billion samples the
counter returns to zero. It will not happen in ten seconds, but a sink that runs for a week
needs the comparison written for it.
\item \textbf{Running the filter on the Pi.} That is the opposite of the optional session
and it proves nothing about the M7. Equally, discarding Linux is the draft-P15 error this
chapter corrects.
\item \textbf{Busy USB during a capture.} A copy, a scan or a second gadget on the same
controller shows up as gaps that the firmware never caused.
\item \textbf{Treating the console cable as a second data path.} Two recorders produce two
files that disagree. The USB-TTL cable is a console, and 3.3 V logic at that.
\end{itemize}

\subsection*{Sources}

\begin{itemize}
\item ST NUCLEO-H7A3ZI-Q board page,
\url{https://www.st.com/en/evaluation-tools/nucleo-h7a3zi-q.html}
\item ST STM32CubeH7 USB device library, CDC class examples,
\url{https://github.com/STMicroelectronics/STM32CubeH7}
\item STM32duino core for STM32, for the Arduino USB-CDC path,
\url{https://github.com/stm32duino/Arduino_Core_STM32}
\item Linux USB CDC ACM driver documentation,
\url{https://www.kernel.org/doc/html/latest/usb/index.html}
\item \texttt{stty} and the terminal line discipline, in the coreutils manual,
\url{https://www.gnu.org/software/coreutils/manual/coreutils.html}
\item pySerial documentation, for the reader loop,
\url{https://pyserial.readthedocs.io/}
\item ST X-NUCLEO-IKS4A1 expansion board page, for the optional session,
\url{https://www.st.com/en/ecosystems/x-nucleo-iks4a1.html}
\end{itemize}


\part{Hosts, consoles and displays}
\project{09}{NanoPi NEO Air headless hub and honest eMMC logging}{NanoPi NEO Air + SEN0032 ADXL345 + Renkforce USB-TTL}{the only non-Raspberry host and the eMMC wear question}

\begin{unique}
Owns the only non-Raspberry Linux host and the only onboard eMMC wear discussion. 24-pin I2C0 plus AP6212 Wi-Fi.
\end{unique}

\begin{blueprint}
Draft P4 and draft P12 seated 40-pin cellular HATs on this board. They do not fit. Draft P6 framed the ADXL345 at 100 Hz as \emph{vibration logging}; that physics belongs to P04. Here the product is the logger and the filesystem, not the spectrum.
\end{blueprint}

\subsection*{Intent}

FriendlyElec Ubuntu runs from the onboard eMMC. The ADXL345 hangs on I2C0, and the board publishes over its onboard Wi-Fi with no monitor, no keyboard and no desktop. The Renkforce USB-TTL cable is a bring-up tool, not part of the product: once the board joins an access point, the cable comes off and the board keeps working.

Two things make this lab different from every Raspberry Pi lab in the book. The header is 24 pins, so the I2C bus number is 0 and no 40-pin HAT can seat. And the storage is soldered down, so an append-heavy log is a wear question rather than a card you can swap. The optional F2FS data partition is the one honest use of the draft-P6 idea: a flash-friendly filesystem on space you repartitioned yourself, never on the rootfs you just booted.

The word \emph{honest} in the title is about that second point. A Raspberry Pi lab can be careless with writes because the medium is a card in a slot, and a worn card costs a few minutes and a re-flash. Here the medium is 8 GB of eMMC soldered to the board, and there is no slot. That changes what a reasonable logger looks like: append rather than rewrite, \texttt{fsync} on a timer rather than on every sample, and a filesystem that was designed for flash translation layers rather than for spinning media.

\begin{kit}
NanoPi NEO Air, SEN0032 ADXL345, Renkforce USB-TTL on UART0, 5 V through the micro-USB socket or through pin 2.
\end{kit}

\begin{rulebox}{Two numbers that are not the Raspberry Pi numbers}
The header has 24 pins, not 40. The sensor bus is 0, not 1. Everything else in this lab is the P01 software with one digit changed.
\end{rulebox}

\subsection*{System architecture}

\diagram{p09_arch}{The headless hub. The console on UART0 is a bring-up path that is removed after Wi-Fi joins; the product path is I2C0 to the publisher to the AP6212 radio. No 40-pin HAT appears anywhere in this figure, because none fits.}

The board carries an Allwinner H3, 512 MB of RAM, 8 GB of eMMC and an AP6212 combo radio. That is the whole system: there is no HAT layer, no USB modem and no display. Everything the lab produces leaves through \texttt{wlan0}, and everything the lab measures arrives through \texttt{i2c-0}.

The number to keep in mind is the bus number. On a Raspberry Pi the sensor bus is \texttt{i2c-1} and the tool is \texttt{i2cdetect -y 1}. On the NEO Air the header bus is I2C0, so it is \texttt{i2cdetect -y 0} and \texttt{smbus.SMBus(0)} in the publisher. The P01 publisher is reused unchanged except for that one digit.

The console deserves its own line in the architecture, because it is the only interface that exists before the network does. A Raspberry Pi can be brought up on HDMI and a keyboard. This board has neither, so \texttt{ttyS0} on the 4-pin debug header is the first and, until \texttt{nmcli} succeeds, the only way in. That is also why the figure shows the console in user space as a getty rather than as part of the product: it is scaffolding, and the acceptance test removes it.

\begin{note}{Which boards are absent from this figure and why}
There is no HAT layer, because the header is 24 pins. There is no USB modem, because the cellular HATs are 40-pin boards and belong to P01, P02 and P03. There is no display stack, because the lab is headless by definition. What is left is a short path from one sensor to one radio, which is the point of using this board at all.
\end{note}

\begin{table}[H]\centering\small
\begin{tabular}{p{44mm}p{26mm}p{74mm}}
\toprule
40-pin board in the bin & Seats here? & Where it belongs \\
\midrule
SIM7600E-H 4G HAT & No & P01, and the failover pair in P13 \\
SIM7020E HAT & No & P02, on a Pi 3 \\
SIM7070G HAT & No & P03, on a Pi 3B+ \\
MCC 118 & No & P05, on a Pi 4 \\
JOY-iT RB-Explorer700 & No & P11, on a Pi 3B+ \\
Waveshare 3.5 inch LCD (A) & No & P06 and P19, on a Pi 3 \\
\bottomrule
\end{tabular}
\caption{The 24-pin rule, written out. No adapter in this kit converts a 40-pin HAT onto this header, and the draft that tried it is corrected at the head of this chapter.}
\end{table}

\subsection*{Wiring and schematic}

\diagram{p09_schematic}{The 24-pin header drawn honestly: four signals to the accelerometer, and a separate 4-pin debug header for the console. The 5 V pin on the debug header is a supply for the cable side, not a signal, and it never reaches SYS\_3V3.}

\begin{table}[H]\centering\small
\begin{tabular}{p{40mm}p{28mm}p{22mm}p{58mm}}
\toprule
Signal & NEO Air header & Pin & Note \\
\midrule
ADXL345 VCC & 24-pin & 1 SYS\_3V3 & Never 5 V \\
ADXL345 SDA & 24-pin & 3 I2C0\_SDA & Bus 0, not bus 1 \\
ADXL345 SCL & 24-pin & 5 I2C0\_SCL & Bus 0, not bus 1 \\
ADXL345 GND & 24-pin & 6 GND & \\
ADXL345 SDO & & & Left floating, selects 0x53 \\
Board 5 V supply & 24-pin & 2 VDD\_5V & Or the micro-USB socket, 5 V at 2 A \\
Console GND & 4-pin debug & GND & Common with the cable \\
Console 5 V & 4-pin debug & 5 V & Cable side only, confirm on the silkscreen \\
Console TX & 4-pin debug & TX & To the USB-TTL receive line \\
Console RX & 4-pin debug & RX & From the USB-TTL transmit line \\
\bottomrule
\end{tabular}
\caption{Wiring. The debug UART0 header is physically separate from the 24-pin header, and the USB-TTL cable must be jumpered for 3.3 V logic.}
\end{table}

\begin{asciiart}
NEO Air 24-pin            ADXL345
1  SYS_3V3 -------------- VCC
3  I2C0_SDA ------------- SDA
5  I2C0_SCL ------------- SCL
6  GND ------------------ GND

UART0 4-pin: GND, 5V, TX, RX --> USB-TTL at 3.3 V logic, 115200 8N1
Do not put 5 V onto SYS_3V3.
\end{asciiart}

\subsection*{Bench layout}

\diagram{p09_bench}{Bench layout. The board is small, the cable bundle is not: strain-relieve the USB-TTL cable so a tug on the console does not lift the sensor jumpers off the 24-pin header.}

Power the board with 5 V at 2 A, either through the micro-USB socket or through pin 2. The H3 plus the AP6212 draw more during a Wi-Fi association than a 500 mA port is willing to lend, and an undervoltage event on this board shows up as a failed association rather than a clear message. Keep the ADXL345 jumpers short and keep the console cable on the opposite side of the board.

Two cable bundles leave the board and they should leave in opposite directions. The sensor bundle is four short jumpers on the 24-pin row. The console is one four-way cable on the debug header, and it is the one you will be pulling on during bring-up. If both leave from the same edge, the pull that unplugs the console also lifts a jumper, and the symptom is an accelerometer that vanished for no visible reason.

\subsection*{Software design (UML)}

\diagram{p09_uml}{First login on UART0, the Wi-Fi join, and headless publishing after the console cable is unplugged. The last frame is the lab: no operator, no cable, and the publishes keep arriving.}

The sequence has three phases and only the third one is the product. Phase one is the serial console: the USB-TTL cable is the only way into a board that has no display and no network yet, so the first login happens there at 115200 8N1. Phase two joins the access point with \texttt{nmcli} and confirms the accelerometer on bus 0. Phase three removes the cable. If the board goes quiet when the console leaves, something in phase two was done in a shell session rather than persisted, and that is the failure this lab is designed to expose.

The publisher itself is the P01 loop and carries the same single piece of state, a cool-off timestamp. It is a token bucket in its simplest form: at most one event every two seconds, however long the surface keeps ringing. On this board there is a second reason to keep it, beyond bandwidth. Every publish that also gets appended to the log is a write to soldered flash, so the cool-off is a wear control as well as a rate control.

\subsubsection*{Where the phases stop being reversible}

Phase one can be repeated as often as you like: reconnect the cable, log in again, nothing is lost. Phase two is the one that has to persist, because \texttt{nmcli device wifi connect} writes a NetworkManager connection profile and that profile is what survives the reboot. Phase three is a test, not a step: if it fails, go back to phase two and check that the profile is on disk rather than in the shell you are standing in.

\subsection*{Data flow (ASCII)}

\begin{asciiart}
  NanoPi NEO Air (Allwinner H3, 512 MB, 8 GB eMMC)
  +----------------------------------------------+
  |  i2c-0 -> 0x53 ADXL345                        |
  |     |                                         |
  |     v                                         |
  |  publisher (P01 loop, SMBus(0))               |
  |     |  magnitude, dead-band, 2 s cool-off     |
  |     +--> mosquitto_pub --> wlan0 (AP6212) -----------> broker on the LAN
  |     |                                         |
  |     +--> optional: append to /data (F2FS)     |
  |            fsync on a timer, not per sample   |
  |                                               |
  |  ttyS0 console on the 4-pin UART0 header      |
  +------------|---------------------------------+
               |  115200 8N1, 3.3 V logic
               v
        Renkforce USB-TTL ---> operator PC   (bring-up only, unplugged for acceptance)
\end{asciiart}

\subsection*{Steps}

\begin{steps}
\item \textbf{Flash FriendlyElec Ubuntu to the eMMC} following the vendor wiki for this board. The first login happens on UART0, because the board has no display output and has not joined a network yet. Open the console at 115200 8N1 with the USB-TTL cable jumpered for 3.3 V logic.

\item \textbf{Install the tools, find the sensor on bus 0, and join the access point.}
\begin{shellcode}
sudo apt-get install -y i2c-tools python3-smbus mosquitto-clients
ls /dev/i2c-*
sudo i2cdetect -y 0         # expect 53
nmcli device wifi connect "YOURAP" password "..."
\end{shellcode}
The bus argument is 0. A Raspberry Pi command line copied from P01 with \texttt{-y 1} will report an empty bus or no such device, and the sensor is not at fault.

\item \textbf{Reuse the P01 publisher with bus number 0.} The dead-band, the magnitude and the two-second cool-off are unchanged. The only edit is the constructor.
\begin{pycode}
import time, math, json, subprocess, smbus
bus = smbus.SMBus(0)                 # bus 0 on the NEO Air, not bus 1
A = 0x53
bus.write_byte_data(A, 0x2D, 0x08)   # measure
bus.write_byte_data(A, 0x31, 0x0B)   # full-res +/-16g

def g(lo, hi):
    v = (hi << 8) | lo
    if v & 0x8000:
        v -= 65536
    return v / 256.0

cool = 0
while True:
    b = bus.read_i2c_block_data(A, 0x32, 6)
    ax, ay, az = g(b[0], b[1]), g(b[2], b[3]), g(b[4], b[5])
    mag = math.sqrt(ax*ax + ay*ay + az*az)
    if abs(mag - 1.0) > 0.35 and time.time() > cool:
        payload = json.dumps({"mag": round(mag, 3), "ax": round(ax, 3)})
        subprocess.run(["mosquitto_pub", "-h", "test.mosquitto.org",
                        "-t", "lab/p09/shock", "-m", payload], check=False)
        cool = time.time() + 2
    time.sleep(0.05)
\end{pycode}

\item \textbf{Optional F2FS data partition}, only if you are willing to repartition unused eMMC space, and not the rootfs you just booted.
\begin{shellcode}
lsblk
# mkfs.f2fs /dev/mmcblk2pX     # X = unused partition YOU created
# mount -o discard /data
# append-only log files, fsync on a timer, not per sample
\end{shellcode}
Read \texttt{lsblk} before you type anything with \texttt{mkfs} in it. The commands above are commented for that reason: the target is a partition you created, and if you cannot name it from the \texttt{lsblk} output then you do not have one yet.

\item \textbf{Confirm the Wi-Fi profile survives a reboot} before you trust it. A connection made in a shell session and never written out is the most common reason this lab appears to fail at the last step.
\begin{shellcode}
nmcli connection show
sudo reboot
# after the board comes back, on the console:
nmcli device status
ip -4 addr show wlan0
\end{shellcode}

\item \textbf{Unplug the console and leave the board alone.} That is the acceptance condition below, and it is also the point of the lab. Confirm from another host on the same network that the publishes are still arriving.
\begin{shellcode}
mosquitto_sub -h test.mosquitto.org -t 'lab/p09/#' -v
free -h
\end{shellcode}
\end{steps}

\subsection*{Acceptance test}

\begin{itemize}
\item With the USB-TTL cable unplugged, the board stays up on Wi-Fi and keeps publishing.
\item \texttt{free -h} shows about 512 MB of RAM.
\item No desktop is running.
\item \texttt{sudo i2cdetect -y 0} shows 53 on bus 0.
\end{itemize}

\subsection*{Practices}

Supply 5 V at 2 A. The AP6212 shares Wi-Fi and Bluetooth, so a Bluetooth scan and a publish burst contend for the same radio and the symptom is jitter in the publish interval, not an error message. Leave the ESP8266 for P17, which is where the air-gap radio pattern lives. Log with \texttt{fsync} on a timer rather than per sample: this is soldered-down flash, and the wear you spend is not recoverable by swapping a card.

Keep the console cable in the kit box rather than on the board. It is a bring-up tool in this lab and a factory tool in P14, and in both places it is jumpered for 3.3 V logic. A cable left permanently on the debug header invites the exact mistake this lab is meant to rule out, which is a board that only works while someone is watching it.

\subsubsection*{What this lab hands to the rest of the book}

P17 reuses this board with the onboard radio switched off and an ESP8266 speaking AT over UART1, which is a different pattern with the same host. P20 sweeps this board as the fourth host in the fleet, so the hostname you set here is the hostname that appears in \texttt{/var/log/lab-fleet.txt}. Nothing else in the book depends on the NEO Air, and nothing that needs a 40-pin header ever will.

\subsection*{Pitfalls}

\begin{itemize}
\item \textbf{Reaching for a 40-pin HAT.} The header is 24 pins. The SIM7020, SIM7070, SIM7600, MCC 118, Explorer700 and the 3.5 inch LCD do not seat on this board, and no adapter in the bin makes them seat.
\item \textbf{Using bus 1.} The symptom is an empty \texttt{i2cdetect} table or a missing device node. The header bus on this board is I2C0.
\item \textbf{5 V onto SYS\_3V3.} Pin 1 is a 3.3 V rail. The 5 V pin on the debug header belongs to the cable side, and the 5 V pin on the 24-pin header is pin 2.
\item \textbf{A 500 mA supply.} The association fails or the board resets during the Wi-Fi join, which reads like a driver fault and is not one.
\item \textbf{Configuring Wi-Fi only inside the console session.} The board then goes quiet the moment the cable comes off, which is exactly the acceptance test failing.
\item \textbf{Running \texttt{mkfs.f2fs} against the rootfs device.} The board no longer boots and the eMMC image has to be flashed again from the vendor wiki.
\item \textbf{Calling a 100 Hz accelerometer stream vibration analysis.} The 6 kHz instrument is the IIS3DWB on the STWIN.box in P04. This lab logs events.
\item \textbf{A USB-TTL cable still jumpered for 5 V from an earlier session.} The board's UART pins are 3.3 V. Check the jumper before the cable touches the debug header, and check it again after P14, where the same cable is used as a factory tool.
\item \textbf{Writing the log with an \texttt{fsync} on every sample.} It works, it is slow, and it spends flash endurance on a board whose storage cannot be swapped. Batch the writes and sync on a timer.
\end{itemize}

\subsection*{Sources}

\begin{itemize}
\item FriendlyElec NanoPi NEO Air wiki, header pinout and eMMC flashing, \url{https://wiki.friendlyelec.com/wiki/index.php/NanoPi_NEO_Air}
\item Allwinner H3 datasheet, I2C0 and UART0 pin multiplexing
\item DFRobot SEN0032 ADXL345 product wiki, \url{https://wiki.dfrobot.com/}
\item Analog Devices ADXL345 data sheet, register map 0x2D, 0x31, 0x32
\item F2FS documentation in the Linux kernel tree, \url{https://www.kernel.org/doc/html/latest/filesystems/f2fs.html}
\item NetworkManager \texttt{nmcli} manual page, for a connection that persists across reboots
\item Broadcom AP6212 combo module notes in the FriendlyElec image, for the shared Wi-Fi and Bluetooth front end
\item \texttt{i2c-tools} manual pages, \texttt{i2cdetect(8)} and the meaning of the bus argument
\item Eclipse Mosquitto client documentation, \texttt{mosquitto\_pub} and \texttt{mosquitto\_sub}, \url{https://mosquitto.org/documentation/}
\item Renkforce USB to TTL cable documentation, for the 3.3 V logic jumper
\end{itemize}

\project{11}{Explorer700 field console}{Raspberry Pi 3B+ + JOY-iT RB-Explorer700}{DS3231 as hwclock, the OLED, IR and the overlay lab}

\begin{unique}
Owns DS3231 as \texttt{hwclock}, the SSD1306 OLED, IR, the joystick, the BMP280 and the PCF8591. The Device Tree overlay lab lives here because the chips do.
\end{unique}

\begin{blueprint}
Draft P7 is this lab. Draft P18 then stacked the same HAT with a SIM7600. Forbidden: that is two 40-pin HATs on one header. Time-series store-and-forward without the second HAT is this: log locally, copy off on Ethernet later.
\end{blueprint}

\subsection*{Intent}

The Pi 3B+ boots, loads an overlay that binds the DS3231 and the BMP280, sets the system time from the real-time clock, and puts status on the OLED. The joystick and the IR receiver are local operator inputs, and the buzzer is the acknowledgement. Nothing in this lab needs a network, which is the point: a field console is a device that is still correct when the network is not there.

This is the Device Tree overlay chapter of the book, and it is here rather than anywhere else because this is where the chips are. Six addressable devices arrive on one board, already wired, with no flying leads to get wrong. That leaves the whole exercise on the kernel side: which node the overlay creates, which driver binds to it, and which user-space interface appears as a result. The acceptance test is a power pull, because an overlay that only works while a shell session is open has not actually bound anything.

\begin{kit}
Pi 3B+, JOY-iT RB-Explorer700. No other 40-pin board.
\end{kit}

\begin{rulebox}{One HAT on this header}
The Explorer700 occupies the 40-pin header for the whole sitting. Do not stack a SIM7600, a SIM7070, a SIM7020, an MCC 118 or the 3.5 inch LCD with it. Power off before unseating it.
\end{rulebox}

\subsection*{System architecture}

\diagram{p11_arch}{Every onboard chip, the kernel binding it gets, and the user-space interface that binding produces. The row is the whole lesson: a chip without a binding is a chip you talk to by address, and a chip with a binding is a device node.}

Read the figure as seven independent rows rather than as one stack. The DS3231 is the row that matters most, because it is the only one where the overlay changes what user space does: with \texttt{dtoverlay=i2c-rtc,ds3231} the chip stops being an address on a bus and becomes \texttt{/dev/rtc0}, which is what \texttt{hwclock} talks to. The BMP280 sits in the middle: it can be bound with a custom overlay, or read from user space through IIO. The remaining chips are addresses and GPIO lines, and that is a legitimate place for them to stay.

The bus assignments are short. I2C1 carries the DS3231 at 0x68, the BMP280 at 0x76 or 0x77, the PCF8591 at 0x48 and the PCF8574 at 0x20. The SSD1306 OLED is on SPI. The IR receiver, the joystick and the buzzer are GPIO lines on the header. The kit cookbook does not print the GPIO numbers for those three, so confirm them on the silkscreen and in the JOY-iT manual before you write them into a script.

Notice what the middle column does not contain. There is no daemon, no polling loop and no service unit anywhere in the binding path. That is the difference this chapter is selling: a device-tree overlay moves work from something that has to be started into something that is simply true after boot. A status script that fails to start leaves you with a blank OLED, which is obvious. A clock script that fails to start leaves you with a plausible wrong time, which is not.

\begin{note}{What an overlay is and what it is not}
A line in \texttt{config.txt} is not a driver. It asks the firmware to graft a node into the device tree before the kernel starts. The node carries a compatible string, the kernel matches that string against the drivers it has, and a successful match produces a device node in user space. If any link in that chain is missing, the chip is still on the bus and still readable by address, and nothing in \texttt{/dev} changes.
\end{note}

\begin{table}[H]\centering\small
\begin{tabular}{p{34mm}p{20mm}p{30mm}p{62mm}}
\toprule
Device & Bus & Address or line & What it is for in this lab \\
\midrule
DS3231 & I2C1 & 0x68 & The system clock across a power pull \\
BMP280 & I2C1 & 0x76 or 0x77 & Pressure and temperature on the status page \\
PCF8591 & I2C1 & 0x48 & Analog inputs on the HAT \\
PCF8574 & I2C1 & 0x20 & Port expander on the HAT \\
SSD1306 OLED & SPI & chip select & Host name and RTC time \\
IR receiver & GPIO & confirm on the silkscreen & A keycode the console reacts to \\
Joystick and buzzer & GPIO & confirm on the silkscreen & Menu input and acknowledgement \\
\bottomrule
\end{tabular}
\caption{The onboard inventory. Addresses come from the default I2C map; the GPIO lines are not printed in the source, so read them off the board.}
\end{table}

\subsection*{Wiring and schematic}

\diagram{p11_schematic}{The HAT is the wiring. One I2C bus carries four devices, SPI0 carries the OLED, and three GPIO lines carry the operator controls. No flying leads are required for any onboard chip.}

\begin{table}[H]\centering\small
\begin{tabular}{p{40mm}p{26mm}p{22mm}p{58mm}}
\toprule
Signal & Pi header pin & BCM GPIO & Note \\
\midrule
HAT 3V3 & 1 & 3V3 & Supplied through the header \\
I2C1 SDA & 3 & GPIO2 & DS3231, BMP280, PCF8591, PCF8574 \\
I2C1 SCL & 5 & GPIO3 & Same four devices \\
HAT GND & 6, 9, 14, 20, 25 & GND & Several ground pins on the header \\
SPI0 MOSI & 19 & GPIO10 & To the SSD1306 \\
SPI0 SCLK & 23 & GPIO11 & To the SSD1306 \\
SPI0 chip select & 24 or 26 & GPIO8 or GPIO7 & Confirm on the silkscreen \\
OLED DC and RESET & & & Not printed in the source, confirm on the silkscreen \\
IR receiver & & & GPIO line, confirm on the silkscreen \\
Joystick & & & GPIO lines, confirm on the silkscreen \\
Buzzer & & & GPIO line, confirm on the silkscreen \\
\bottomrule
\end{tabular}
\caption{Header use. The four I2C devices and the OLED are already routed on the HAT; the entries left blank are ones the source does not print.}
\end{table}

\begin{asciiart}
Pi 3B+ 40-pin                    JOY-iT RB-Explorer700 (seated, no flying leads)
pin 1  3V3  ------------------>  HAT supply
pin 3  GPIO2 SDA -------------+-> DS3231   0x68
                              +-> BMP280   0x76 or 0x77
pin 5  GPIO3 SCL -------------+-> PCF8591  0x48
                              +-> PCF8574  0x20
pin 6  GND  ------------------>  HAT ground (also 9, 14, 20, 25)
pin 19 GPIO10 MOSI -----------+
pin 23 GPIO11 SCLK -----------+-> SSD1306 OLED, SPI
pin 24 or 26 chip select -----+   DC and RESET: confirm on the silkscreen
GPIO lines -------------------->  IR receiver, joystick, buzzer
                                  line numbers: confirm on the silkscreen
Do not stack a second 40-pin HAT on this header.
\end{asciiart}

\subsection*{Bench layout}

\diagram{p11_bench}{Bench layout for the power-pull test. The network cable is unplugged and stays unplugged, and the supply is pulled at the wall rather than by a shutdown, because a graceful shutdown is not the test.}

The bench for this lab is almost empty, which is intentional. One Pi, one HAT, one supply, and a network cable that you remove before the acceptance test. Do not leave a second HAT on the bench where it can be seated by reflex: the Explorer700 owns the header until the lab is finished and torn down.

Check the coin cell before the sitting starts. The DS3231 keeps time from its own backup cell while the board is unpowered, and a cell that has been sitting in a drawer since the kit was bought may not have enough left to carry an hour. The symptom is specific and misleading: the clock is correct while the Pi runs, and wrong the moment it comes back from a power pull, which looks exactly like an overlay that did not load.

\subsection*{Software design (UML)}

\diagram{p11_uml}{Component view of the overlay path. The overlay line in \texttt{config.txt} creates a device-tree node, the node causes a driver to bind, and the driver publishes a user-space interface. The BMP280 shows the alternative: a custom overlay, or no binding at all and an IIO or I2C read from user space.}

The component diagram is the argument of the whole chapter. A line in \texttt{/boot/firmware/config.txt} is not a driver and it is not a device. It is a request to add a node to the device tree, and the node is what makes the kernel look for a driver that matches the compatible string. When the match succeeds you get \texttt{/dev/rtc0}, and \texttt{hwclock} works with no knowledge of I2C at all. When it does not succeed you are left with 0x68 on a bus, which you can still read by hand, and which \texttt{hwclock} will not touch.

That distinction is why the acceptance test is a power pull rather than a command. Reading the correct time out of the DS3231 with a Python script proves that the chip works. Having the correct time after the supply was removed, with no network available, proves that the binding exists and that the boot sequence uses it.

The BMP280 is drawn as two lanes because the source gives a choice and does not resolve it. One lane binds the chip with a custom overlay and gets IIO channels in sysfs. The other lane creates no node at all and reads the chip by address from user space. Both are legitimate here, and the difference matters only if something later has to depend on the pressure reading at boot, which nothing in this book does.

\subsubsection*{Why the status page reads the chip rather than the clock}

Put the time on the OLED from \texttt{/dev/rtc0}, not from the system clock. The system clock is a value the kernel maintains in memory and it can be right for reasons that have nothing to do with this HAT. Reading the device node instead means the page is showing you the chip, so a stopped clock or a dead cell is visible on the glass rather than inferred afterwards.

\subsection*{Data flow (ASCII)}

\begin{asciiart}
  boot                                     steady state
  ----                                     ------------
  config.txt                               joystick -> menu selection
    dtparam=i2c_arm=on                     IR       -> keycode
    dtoverlay=i2c-rtc,ds3231               buzzer   <- acknowledgement
       |                                       |
       v                                       v
  device-tree node rtc@68 on i2c1          +-------------------------+
       |                                   |  status page, python    |
       v                                   |    host name            |
  rtc-ds1307 driver binds                  |    date from /dev/rtc0  |
       |                                   |    BMP280 reading       |
       v                                   +------------+------------+
  /dev/rtc0                                             |
       |                                                v
       +--> systemd reads the RTC at boot        SSD1306 OLED over SPI
       +--> hwclock -w writes it back
       +--> hwclock -r reads it back

  no network is involved in any of this
\end{asciiart}

\subsection*{Steps}

\begin{steps}
\item \textbf{Seat the HAT on a powered-off Pi 3B+}, then install the tools and confirm that all four I2C devices answer.
\begin{shellcode}
sudo apt-get install -y i2c-tools python3-smbus python3-luma.oled
sudo i2cdetect -y 1
# expect 68, 76 or 77, 48, 20
\end{shellcode}
If 0x68 is missing, stop here. Every later step in this lab depends on the real-time clock being on the bus.

\item \textbf{Bind the clock with an overlay rather than with user-space bit-banging.} This is the fragment for \texttt{/boot/firmware/config.txt}.
\begin{dtscode}
dtparam=i2c_arm=on
dtoverlay=i2c-rtc,ds3231
// BMP280 can be bound with a custom overlay or read through iio userspace
\end{dtscode}
Reboot after the edit. After the reboot, \texttt{/dev/rtc0} should exist, and 0x68 usually disappears from the \texttt{i2cdetect} table because the driver now holds the address. That is a binding working, not a device going missing.

\item \textbf{Set and persist the time}, then confirm that the setting survives without a network.
\begin{shellcode}
timedatectl
sudo hwclock -w
sudo hwclock -r
# disconnect Ethernet/Wi-Fi, reboot, confirm the clock survived
\end{shellcode}
Write the clock while the Pi still has correct time from the network, and only then remove the network. A DS3231 written from a wrong system clock is a battery-backed wrong answer.

\item \textbf{Bring up the OLED} with the luma.oled SSD1306 SPI example from the JOY-iT manual. Put two lines on it: the host name and the time read from the real-time clock. That is the status page the acceptance test looks at.

\item \textbf{Map the operator controls.} Use the vendor examples under RB-Explorer700. Map the joystick to a menu, the IR receiver to a keycode, and the buzzer to an acknowledgement. The GPIO line numbers are not printed in the source, so read them from the silkscreen and the JOY-iT manual before writing them into a script.

\item \textbf{Run the power-pull test.} With Ethernet and Wi-Fi disconnected, remove the supply, wait, and power the board back up. Then compare \texttt{date} against a clock you trust.
\begin{shellcode}
timedatectl show -p NTPSynchronized
date
sudo hwclock -r
sudo i2cdetect -y 1
\end{shellcode}
\end{steps}

\subsection*{Acceptance test}

\begin{itemize}
\item After a power pull with no network, \texttt{date} is still correct within 2 s.
\item The OLED shows the host name and the time from the real-time clock.
\item \texttt{i2cdetect} still shows 0x68, or shows it held by the driver after the overlay has bound.
\item No network interface was used to recover the time.
\end{itemize}

\subsection*{Practices}

This HAT is the real-time clock source for any later store-and-forward idea in the book. Do not move it onto a cellular Pi: that would put two 40-pin HATs on one header, which is the error the blueprint box at the head of this chapter corrects. If you want store-and-forward without a second HAT, the honest version is the one the source gives: log locally, and copy the files off over Ethernet later.

Prefer an overlay to a script for anything that has to be true at boot. A script that sets the clock is one more thing that can fail to start, and it fails in a way that looks like a hardware fault. Keep the OLED page short, and read the time on it from \texttt{/dev/rtc0} rather than from the system clock, so the page shows the chip rather than the assumption.

Tear the HAT down when the lab is finished, and note the date you wrote the clock. The next lab on this Pi wants an empty header, and the sequence in the source is explicit about it: P11 comes early, and the HAT comes off before P12 and the display labs.

\subsubsection*{What this lab hands to the rest of the book}

This is the only real-time clock in the kit, so any later idea about store-and-forward, buffered uplinks or timestamped local logs traces back to this chapter. The honest pattern, and the one the source gives, is to log locally with a correct clock and copy the files off over Ethernet later. The dishonest pattern is to move the Explorer700 onto a cellular Pi so that one host has both a clock and a radio, and that is two 40-pin HATs on one header.

\subsection*{Pitfalls}

\begin{itemize}
\item \textbf{Stacking a second 40-pin HAT.} The Explorer700 and a cellular HAT cannot share a header. Split the work across hosts or drop one HAT.
\item \textbf{Reading 0x68 by hand and calling it done.} A Python read proves the chip is alive. Only \texttt{/dev/rtc0} and a power pull prove the binding exists.
\item \textbf{Writing the real-time clock from a wrong system clock.} \texttt{hwclock -w} copies whatever the system believes. Set the system clock from the network first, then write, then disconnect.
\item \textbf{Being surprised when 0x68 leaves the \texttt{i2cdetect} table.} After the overlay binds, the address is held by the driver. That is the expected outcome, not a missing device.
\item \textbf{A graceful shutdown instead of a power pull.} A clean shutdown can write the time on the way down, which hides exactly the failure the test is looking for.
\item \textbf{Leaving the network connected during the test.} Time then arrives from the network and the real-time clock is never exercised.
\item \textbf{Guessing the GPIO numbers for the IR receiver, joystick and buzzer.} They are not printed in the source. Read them from the silkscreen and the JOY-iT manual.
\item \textbf{A depleted coin cell in the DS3231 holder.} The symptom is a clock that is correct until the supply is removed and wrong afterwards, which reads like an overlay problem and is not one.
\item \textbf{Putting the system clock on the OLED instead of the chip.} The page then agrees with itself whatever the hardware is doing, which is the one thing a status page must not do.
\item \textbf{Editing the wrong config file.} On Bookworm the file is \texttt{/boot/firmware/config.txt}. An edit to an older path is accepted by the editor and ignored by the firmware, and the symptom is an overlay that never appears to load.
\item \textbf{Forgetting to reboot after the overlay edit.} Overlays are applied by the firmware before the kernel starts, so a running system does not pick one up.
\end{itemize}

\subsection*{Sources}

\begin{itemize}
\item JOY-iT RB-Explorer700 product page and manual, for the board layout, the joystick, the IR receiver and the OLED example, \url{https://joy-it.net/}
\item Raspberry Pi device tree and overlay documentation, \url{https://www.raspberrypi.com/documentation/computers/configuration.html}
\item The overlay README shipped with Raspberry Pi OS, \texttt{/boot/firmware/overlays/README}, entry \texttt{i2c-rtc}
\item Maxim DS3231 data sheet, address 0x68, the ageing register and the coin-cell backup
\item Bosch BMP280 data sheet, address 0x76 or 0x77 depending on the SDO strap
\item NXP PCF8591 and PCF8574 data sheets, addresses 0x48 and 0x20 to 0x27
\item Solomon Systech SSD1306 data sheet, and the \texttt{luma.oled} documentation, \url{https://luma-oled.readthedocs.io/}
\item \texttt{hwclock(8)} and \texttt{timedatectl(1)} manual pages
\item Linux kernel RTC class documentation, for \texttt{/dev/rtc0} and the boot-time read
\item Linux IIO subsystem documentation, for the sysfs channels a bound pressure driver publishes
\item \texttt{i2c-tools} manual pages, and the reason a bound address leaves the \texttt{i2cdetect} table
\end{itemize}

\project{12}{Official DSI operator glass and MQTT broker}{Raspberry Pi 4 + official DSI touchscreen + official keyboard}{the panel, the keyboard and the broker}

\begin{unique}
Owns the 7 inch DSI panel, the keyboard, and the broker. No 40-pin HAT, so any USB gadget from P04, P16 and P18 can sit beside it.
\end{unique}

\begin{blueprint}
Draft P14 is this lab plus ESP32 actuation. Keep the actuation in P16 so this chapter stays a console and broker lab.
\end{blueprint}

\subsection*{Intent}

The Pi 4 with the official touchscreen runs labwc and Wayland, which is the Bookworm default, and it runs Mosquitto. A full-screen page shows fleet topics. The keyboard is the operator input. That is the whole product: a piece of glass that shows what the rest of the bench is saying, and a broker that the rest of the bench can publish into.

There is a reason the console and the broker are one lab rather than two. An operator console with nothing to show is a browser pointed at a blank page, and a broker with no display is a process you check by running a command. Putting them on the same host gives each one a purpose: the broker has an audience, and the glass has something to render that did not come from itself.

The second half of the lab is a coexistence statement. The official touchscreen connects to the DISPLAY connector, not to the 40-pin header, so this host keeps its header free. That is what makes it the natural home for the broker: it can wear the glass and still accept the USB gadgets from the other labs, the STWIN.box from P04, the ESP32 from P16 and the bare Nucleo from P18. The one thing it does not do today is wear a HAT.

\begin{kit}
Pi 4, official Raspberry Pi touchscreen on the DISPLAY connector, official keyboard.
\end{kit}

\begin{rulebox}{The header stays empty today}
The DSI panel does not consume the 40-pin header. This host is allowed to wear a HAT in a later sitting. Today it does not.
\end{rulebox}

\subsection*{System architecture}

\diagram{p12_arch}{The DSI panel hangs off the DISPLAY connector and the keyboard off USB, so the 40-pin header is untouched and the USB bus is free for the gadgets of P04, P16 and P18. The broker sits in user space with a LAN listener and a local subscriber.}

Read the figure from the connector names rather than from the boxes. DISPLAY is a dedicated connector with its own ribbon, so the panel consumes no GPIO lines and no I2C addresses on the header. USB is a bus with several ports, so the keyboard and the gadget of the day share it without argument. The 40-pin header is drawn empty on purpose: it is the one resource this lab deliberately does not spend.

Mosquitto is the other half. It listens on 1883 for the LAN, it is enabled as a systemd unit so it comes back after a reboot, and it is locked with a password file rather than left open. A full-screen page renders the topics on the glass, either a local HTML file that subscribes over websockets on listener 9001, or three \texttt{watch} terminals tiled on the panel, which is the simpler version and works.

\begin{note}{Enabled is not the same as started}
\texttt{systemctl start mosquitto} runs the broker now. \texttt{systemctl enable mosquitto} arranges for it to run after the next boot. The acceptance test measures the second one, and \texttt{enable --now} does both in a single command. A broker that was only started works perfectly until the reboot and then looks like a hardware problem.
\end{note}

\begin{table}[H]\centering\small
\begin{tabular}{p{40mm}p{34mm}p{78mm}}
\toprule
Resource on the Pi 4 & Used by this lab & Consequence \\
\midrule
DISPLAY connector & The official touchscreen & The panel has its own ribbon and its own connector \\
USB ports & Keyboard, and any gadget & STWIN.box (P04), ESP32 (P16), Nucleo (P18) all fit beside the glass \\
40-pin header & Nothing & Free for a later sitting, empty today \\
Ethernet or Wi-Fi & The broker listener & Other hosts publish into 1883 \\
\bottomrule
\end{tabular}
\caption{What the lab spends and what it leaves alone. The header row is the reason this host can be the broker for the rest of the bench.}
\end{table}

\subsection*{Wiring and schematic}

\diagram{p12_schematic}{There is no flying-lead wiring in this lab. The drawing shows the two connectors that are used, the 40-pin header that is not, and the USB ports that stay available.}

\begin{table}[H]\centering\small
\begin{tabular}{p{42mm}p{30mm}p{80mm}}
\toprule
Connection & Connector & Note \\
\midrule
Official touchscreen & DISPLAY & 15-way ribbon, seated with the Pi powered off \\
Official keyboard & USB-A & Operator input \\
USB gadget of the day & USB-A & Optional: STWIN.box, ESP32 or Nucleo \\
Network & Ethernet or Wi-Fi & Carries the broker listener on 1883 \\
40-pin header & & Left empty for the whole sitting \\
\bottomrule
\end{tabular}
\caption{Connections. The source gives no flying leads for this lab, because there are none.}
\end{table}

\begin{asciiart}
Pi 4                          Official 7 inch touchscreen
DISPLAY connector -------->   15-way DSI ribbon, seated with the Pi powered OFF
USB-A ------------------->    official keyboard
USB-A ------------------->    optional gadget: STWIN.box (P04), ESP32 (P16), Nucleo (P18)
Ethernet or Wi-Fi ------->    LAN, broker listener 1883

40-pin header:  EMPTY.  The DSI panel does not consume it.
                This host may wear a HAT in a later sitting. Today it does not.
\end{asciiart}

\subsection*{Bench layout}

\diagram{p12_bench}{Bench layout. The glass stands on its own frame, the keyboard sits in front of it, and the 40-pin header is visible and empty. The publishing host is somewhere else on the same LAN.}

Seat the DSI ribbon with the Pi powered off, both at the panel end and at the DISPLAY connector. Bookworm detects the official panel, so there is no overlay to add and nothing to configure for the display itself. Leave the header exposed rather than covered: this host is the one place on the bench where an empty header is the feature.

The publishing host in the drawing is deliberately vague. It can be the Pi 3B+ from P11, the NanoPi from P09, or a laptop with \texttt{mosquitto\_pub} installed. What matters for the acceptance test is only that it is a different machine on the same network, because a publish from \texttt{localhost} measures nothing about the path the lab is claiming.

\subsection*{Software design (UML)}

\diagram{p12_uml}{One publish crossing the LAN and reaching the glass. The acceptance number is on the right: the whole path has to complete in under one second, and every hop in it is local.}

The sequence is short because the path is short. A publisher on another host opens a connection to 1883 on this Pi, sends a message on a \texttt{lab/} topic, and the broker fans it out to whoever is subscribed. On this host the subscriber is either the local HTML page over websockets or a \texttt{watch} terminal, and either way the compositor draws the change on the panel. There is no cloud hop in the diagram, which is why one second is a generous budget rather than a tight one.

If the number is missed, the interesting question is which hop consumed it. A slow LAN shows up as a delay before the broker logs the message. A slow page shows up as a delay after it. Subscribing locally with \texttt{mosquitto\_sub} while the page is also running separates the two without any instrumentation.

\subsubsection*{The half of the acceptance test with no messages in it}

The second clause of the acceptance test does not appear in the sequence at all, because it is not a message exchange. Reboot the Pi 4, wait, and check that the broker is running again without anyone having started it. That is the difference between \texttt{systemctl start} and \texttt{systemctl enable}, and it is the only part of this lab that cannot be faked by leaving a terminal open.

\subsection*{Data flow (ASCII)}

\begin{asciiart}
  another host on the LAN                 Pi 4 with the glass
  +---------------------+                 +-------------------------------------+
  | mosquitto_pub       |                 |  mosquitto  (systemd unit, enabled) |
  |   -h pi4            |---- 1883 ------>|    listener 1883                    |
  |   -t lab/...        |                 |    allow_anonymous false            |
  +---------------------+                 |    password_file /etc/mosquitto/... |
                                          +------------------+------------------+
                                                             |
                                        +--------------------+-------------------+
                                        |                                        |
                                        v                                        v
                            mosquitto_sub -t 'lab/#' -v          websockets listener 9001
                            (a watch terminal on the glass)      (a local dash.html page)
                                        |                                        |
                                        +------------------+---------------------+
                                                           v
                                              labwc / Wayland compositor
                                                           v
                                              official 7 inch DSI panel
                                              budget for the whole path: under 1 s

  40-pin header: unused.  USB: keyboard + optional gadget from P04, P16 or P18.
\end{asciiart}

\subsection*{Steps}

\begin{steps}
\item \textbf{Seat the DSI ribbon with the Pi powered off.} Power on and confirm the desktop appears on the glass. Bookworm detects the official panel, so nothing has to be added to \texttt{config.txt} for the display.

\item \textbf{Install the broker and confirm it is running and subscribable.}
\begin{shellcode}
sudo apt-get install -y mosquitto mosquitto-clients
sudo systemctl enable --now mosquitto
mosquitto_sub -h localhost -t 'lab/#' -v
\end{shellcode}
The \texttt{enable} half of that command is the half the acceptance test checks: the unit has to come back on its own after a reboot.

\item \textbf{Put the topics on the glass.} Either point a full-screen Chromium at a local file that subscribes over websockets, with the Mosquitto websockets listener on 9001, or keep it simpler and tile three \texttt{watch} terminals on the panel.
\begin{shellcode}
# the page version
chromium-browser --kiosk file:///home/pi/p12/dash.html
# the simpler version, three terminals tiled on the glass
watch -n 1 "mosquitto_sub -h localhost -t 'lab/p01/#' -v -C 5"
\end{shellcode}
Three terminals are a legitimate answer here. The product of this lab is an operator console, not a web application.

\item \textbf{Lock the broker to the LAN} rather than leaving it anonymous.
\begin{plaincode}
# /etc/mosquitto/conf.d/lab.conf
listener 1883
allow_anonymous false
password_file /etc/mosquitto/passwd
\end{plaincode}
Create the password file before restarting the unit, or every client on the bench stops being able to connect at once.
\begin{shellcode}
sudo mosquitto_passwd -c /etc/mosquitto/passwd lab
sudo systemctl restart mosquitto
systemctl status mosquitto
\end{shellcode}

\item \textbf{Publish from another host on the LAN} and watch the glass.
\begin{shellcode}
# on the other host
mosquitto_pub -h pi4 -t lab/p12/hello -m "from the bench" -u lab -P '...'
\end{shellcode}

\item \textbf{Reboot and check that nothing needed a human.} The broker unit comes back, the panel comes back, and the page or the terminals come back to whatever you arranged to start them.
\end{steps}

\subsection*{Acceptance test}

\begin{itemize}
\item A publish from another host on the LAN appears on the glass in under 1 s.
\item A reboot survives the broker unit: \texttt{systemctl is-enabled mosquitto} reports it enabled and the unit is active without anyone starting it.
\item The 40-pin header is empty at the end of the sitting, as it was at the start.
\item An anonymous connection is refused once \texttt{lab.conf} is in place.
\end{itemize}

\subsection*{Practices}

This host is allowed to wear a HAT in a later sitting. Today it does not. Keep that rule visible while the lab runs, because this Pi 4 is also the natural host for the MCC 118 in P05 and for the SIM7600 in P01, and the temptation to seat one of them while the glass is already working is exactly how a console lab turns into a debugging session about two things at once.

Keep the broker boring. One listener on 1883, one password file, one enabled unit. The websockets listener on 9001 is optional and exists only to serve the local page. Nothing in this book needs the broker to be reachable from outside the LAN, and the lab does not open it.

Write down the user name and the topic prefix somewhere the whole bench can see. Every other lab that publishes needs both, and a password file that only one person can remember turns a shared broker into a single-operator broker.

\subsubsection*{What this lab hands to the rest of the book}

P16 subscribes to this broker and drives its LEDs from it, which is why the actuation stays there and not here. P20 sweeps the fleet from this host or from any Pi with SSH keys, so the keyboard and the glass are a convenient place to watch that sweep run. P04 and P18 both want a USB port and a display, and both get them here without touching the header.

\subsection*{Pitfalls}

\begin{itemize}
\item \textbf{Seating the DSI ribbon under power.} Do it with the Pi powered off, at both ends of the cable.
\item \textbf{Starting the broker without enabling it.} It works all afternoon and is gone after the reboot, which is the acceptance test failing on the one point it actually measures.
\item \textbf{Writing \texttt{allow\_anonymous false} without creating the password file.} Every client on the bench is refused at once, and the cause is the missing file rather than the clients.
\item \textbf{Adding the ESP32 actuation here.} The LEDs and the subscribe-and-drive pattern belong to P16. This chapter stays a console and broker lab.
\item \textbf{Seating a HAT because the header is free.} It is free on purpose. Tear down and seat it in a separate sitting, with the header rule from the front matter in view.
\item \textbf{Measuring the one-second budget with the page as the only subscriber.} Subscribe with \texttt{mosquitto\_sub} at the same time, so a slow broker and a slow page can be told apart.
\item \textbf{Assuming the panel needs an overlay.} Bookworm detects the official display. An overlay copied from a third-party panel note is a way to lose a working screen.
\item \textbf{Publishing from \texttt{localhost} and calling the test passed.} The acceptance test says another host on the LAN. A loopback publish exercises none of the path the number is about.
\item \textbf{Leaving the broker anonymous because it is only a lab.} The bench is a network, other labs publish into it, and an unauthenticated broker on a shared network is a habit worth not forming.
\item \textbf{Running the websockets page without the 9001 listener configured.} The page loads, shows nothing, and looks like a broker fault. Three \texttt{watch} terminals have none of that failure surface.
\end{itemize}

\subsection*{Sources}

\begin{itemize}
\item Raspberry Pi official touchscreen documentation and the DISPLAY connector, \url{https://www.raspberrypi.com/documentation/accessories/display.html}
\item Raspberry Pi OS Bookworm release notes, labwc and the Wayland default
\item Eclipse Mosquitto documentation, \texttt{mosquitto.conf}, listeners and websockets, \url{https://mosquitto.org/documentation/}
\item \texttt{mosquitto\_passwd(1)}, \texttt{mosquitto\_pub(1)} and \texttt{mosquitto\_sub(1)} manual pages
\item \texttt{systemctl(1)} manual page, for the difference between \texttt{start} and \texttt{enable}
\item Raspberry Pi hardware documentation, 40-pin header and the DISPLAY and CAMERA connectors, \url{https://www.raspberrypi.com/documentation/computers/raspberry-pi.html}
\item labwc documentation, the compositor Bookworm uses by default, \url{https://labwc.github.io/}
\item MQTT version 3.1.1 specification, for CONNECT, CONNACK and PUBLISH
\item Chromium command-line switches, for the kiosk mode used by the page version
\end{itemize}

\project{19}{SPI LCD tilt meter}{Raspberry Pi 3 + Waveshare 3.5 inch LCD (A) + SEN0032 ADXL345}{the LCD as an offline field instrument}

\begin{unique}
Owns the 3.5 inch LCD as a field instrument with ADXL345 tilt. No broker, no radio. The LCD owns the 40-pin header, so this cannot run with P06 on the same Pi at the same time, a different sitting.
\end{unique}

\subsection*{Intent}

A Pi 3, a Waveshare 3.5 inch (A) panel and an ADXL345. Pitch and roll on the glass, drawn as two bars, and nothing else. The instrument is battery-unaware: it is an offline meter, not a gateway, and there is no part of it that cares whether a network exists.

That absence is the lab. Every other display in this book is attached to something: P12 shows broker topics, P06 shows ranges arriving from a Nucleo, P04 shows a waterfall from a USB probe. This one shows a quantity it measured itself, on a panel it drives itself, with no client and no uplink. If an MQTT client is running during the acceptance test, the lab has drifted into being a different lab.

\begin{kit}
Pi 3, Waveshare 3.5 inch RPi LCD (A), SEN0032 ADXL345, jumpers for the sensor. This panel has no pass-through, so read the wiring note before planning the sitting.
\end{kit}

\begin{rulebox}{The LCD owns the 40-pin header}
This lab and P06 both want the 3.5 inch LCD on a Pi 3. They cannot share one Pi in the same sitting. Build one, accept it, tear it down, then build the other.
\end{rulebox}

\subsection*{System architecture}

\diagram{p19_arch}{An instrument with no network side. SPI carries the panel, I2C carries the accelerometer, and the loop between them never leaves the board.}

There are exactly two buses in this lab and both of them stay on the host. SPI drives the panel through the vendor overlay, which on Bookworm means the current Waveshare note for the 3.5 inch (A) rather than the flexfb recipes that circulated in 2016. I2C1 carries the ADXL345 at 0x53, and reaching it is the difficulty of this lab rather than a detail of it, because the panel covers the pins the bus lives on. Between them sits one Python loop that reads six registers, computes two angles and paints two bars.

Nothing else is present. There is no broker, no radio, no USB gadget and no second host. Confirm that a framebuffer exists before writing any drawing code, exactly as P06 does: either \texttt{/dev/fb0} or a DRM card, depending on what the vendor overlay gives you on your image.

\begin{note}{This lab and P06 use the same panel}
P06 puts the same Waveshare 3.5 inch (A) on a Pi 3 and paints an 8x8 range grid that arrives over USB from a Nucleo. This lab paints two bars from a sensor on the host's own I2C bus. They share the panel, the host and the header, so they cannot share a sitting. The overlay work is the same in both, which is why this chapter points at P06 for it rather than repeating it.
\end{note}

\subsection*{Wiring and schematic}

\diagram{p19_schematic}{The panel covers the pins the sensor needs, and the two ways past that. SDO left floating selects 0x53, the same address the sensor uses in P01 and P09.}

\begin{table}[H]\centering\small
\begin{tabular}{p{42mm}p{26mm}p{22mm}p{58mm}}
\toprule
Signal & Pi header pin & BCM GPIO & Note \\
\midrule
ADXL345 VCC & 1 & 3V3 & Never 5 V \\
ADXL345 SDA & 3 & GPIO2 & I2C1 data \\
ADXL345 SCL & 5 & GPIO3 & I2C1 clock \\
ADXL345 GND & 6 & GND & \\
ADXL345 SDO & & & Left floating, selects 0x53 \\
LCD panel & 40-pin header & & The HAT owns the header for the sitting \\
\bottomrule
\end{tabular}
\caption{Wiring as the source gives it, which this panel cannot be wired to. The four rows above assume a pass-through, and the Waveshare 3.5 inch (A) has none. Read the note before touching a jumper.}
\end{table}

\begin{note}{This panel has no pass-through, and the lab is blocked until that is worked around}
The source says to wire the sensor to pins 1, 3, 5 and 6 "if the LCD HAT pass-through exposes them (it usually does)". On this panel it does not. The board mates with the first 26 pins through a female header on its underside and carries no header on top, so that whole corner sits under the panel and no jumper reaches it.

The arithmetic is worse than needing different pins. Hardware I2C1 exists only on pins 3 and 5, so there is no second I2C bus to move to. Pin 1 and pin 17 are the only 3.3 V pins on the header and both are inside the covered range, so with the panel seated there is no supply pin left either. What remains above pin 26 is grounds on 30, 34 and 39 and GPIOs on 29, 31, 32, 33, 35, 36, 37, 38 and 40, with pins 27 and 28 left alone as the identity EEPROM lines.

One detail from the panel's own pinout is what makes a fix possible rather than impossible: it declares pins 3 and 5 not connected. The I2C bus is electrically free and only physically buried, so nothing is contending for it.

\textbf{There was a second way out, and it is closed.} The attractive route was to leave the panel seated, take 3.3 V from the POW-BB rail that P15 already labels, ground to pin 30, 34 or 39, and run a bit-banged bus through the \texttt{i2c-gpio} overlay on two of the free GPIOs above pin 26. It was attractive because it costs nothing and because it would have made P15's rails load-bearing in a later lab rather than a one-off.

It does not work, and the carrier's own schematic is what says so. Its Raspberry Pi connector symbol is drawn as thirteen rows of two, numbered 1 and 2 down to 25 and 26, and it stops there; the netlist carries twenty-six pads and nothing above. Pins 27 to 40 are not contacts on this board and are not listed even as not connected. The carrier outline is 85.06 by 56.21 mm against the Pi's 85 by 56 mm, so the board runs the full length of the header and past it, and pin 1 of the socket sits at the SD card and power end. Those fourteen pins are electrically untouched and mechanically buried the moment the panel is seated, so there is nowhere to attach a jumper to them.

\textbf{So the lab needs one part: a 2x20 stacking header, or a tall header whose pins 27 to 40 clear the 26-pin socket.} It is not on this bench. That is a better position than the lab held before, because it is no longer blocked on something unknown; it is blocked on a purchase it can name.

Settled on Thursday 8 October 2026, and worth recording how. An earlier version of this note asked for a ruler along the carrier's long edge, and an inference from the vendor's wiki sentence and the glass panel's dimension drawing already pointed at this same answer. That inference was deliberately not written down, because it rested on a marketing phrase and a drawing of a different part. The schematic settled it without any measuring at all, which is the lesson worth keeping: the cheapest measurement is often a document nobody had opened.
\end{note}

\begin{asciiart}
ADXL345                 Pi 40-pin (NOT reachable: the panel covers pins 1-26)
VCC ------------------- pin 1  3V3      NEVER 5V
GND ------------------- pin 6  GND
SDA ------------------- pin 3  GPIO2
SCL ------------------- pin 5  GPIO3
SDO floating => 0x53

Waveshare 3.5" LCD (A) seated on the Pi 3 40-pin header: this lab's HAT.
Use the current Bookworm vendor overlay, not the 2016 flexfb recipe.
This cannot share a Pi with P06 in the same sitting.
\end{asciiart}

\subsection*{Bench layout}

\diagram{p19_bench}{Bench layout. The board is meant to be picked up and tilted, so the sensor has to be fixed to the same body as the panel and the jumpers have to be short enough to survive being moved.}

The instrument gets handled during its own acceptance test, which is unusual for this book. Fix the ADXL345 to the same piece of board or case as the Pi, and keep the four jumpers short. A sensor that swings on long leads while you rotate the assembly reports its own swing as tilt, and the bars will tell you so.

Note the sensor's orientation on the assembly before the first rotation, because pitch and roll are defined by which axis points where. There is no calibration step in this lab, so the mapping from a physical rotation to a bar is whatever the mounting made it. If the wrong bar moves, rotate the sensor rather than editing the formulas.

\subsection*{Software design (UML)}

\diagram{p19_uml}{The activity view of one pass: read the six registers, convert to g, compute pitch and roll, draw two bars, repeat. There is no branch in this diagram that reaches a network.}

The loop is deliberately flat. Read six bytes starting at 0x32, convert each axis pair to g the way P01 does, compute the two angles, and draw. Pitch is $\arctan$ of $a_x$ against $\sqrt{a_y^{2} + a_z^{2}}$, and roll is $\arctan$ of $a_y$ against $\sqrt{a_x^{2} + a_z^{2}}$, both taken with the two-argument form so the sign survives. Everything after that is drawing: two bars, full scale at the ends of the range you chose.

One thing is worth stating because it is easy to want: there is no filter here and no calibration step. The source gives the two formulas and the drawing, and that is the whole product. Anything more belongs to the labs that own the physics of orientation, not to a panel demonstration.

The other thing the diagram shows by omission is that there is no failure path worth drawing. An instrument with a network has states for connecting, retrying and giving up. This one has a read that either returns six bytes or raises, and a draw that either paints or does not. That is the trade the chapter is making: less reach, far less to go wrong.

\subsection*{Data flow (ASCII)}

\begin{asciiart}
  Pi 3, offline
  +------------------------------------------------------------------+
  |  i2c-1 -> 0x53 ADXL345                                            |
  |     |  read 6 bytes at 0x32                                       |
  |     v                                                             |
  |  ax, ay, az  (g, full resolution)                                 |
  |     |                                                             |
  |     +--> pitch = atan2(ax, sqrt(ay*ay + az*az))                   |
  |     +--> roll  = atan2(ay, sqrt(ax*ax + az*az))                   |
  |               |                                                   |
  |               v                                                   |
  |  two bars on the framebuffer  ->  SPI  ->  3.5" LCD (A), 480x320  |
  +------------------------------------------------------------------+

  no broker, no radio, no USB gadget, no second host
  no MQTT client is running during the acceptance test
\end{asciiart}

\subsection*{Steps}

\begin{steps}
\item \textbf{Seat the LCD} on a powered-off Pi 3 and confirm the vendor overlay as in P06. On Bookworm that means the current Waveshare 3.5 inch (A) note, not the flexfb recipe from 2016. Confirm that a framebuffer exists before writing any drawing code.
\begin{shellcode}
ls /dev/fb* /dev/dri/card* 2>/dev/null
\end{shellcode}

\item \textbf{Wire the ADXL345} by whichever of the two routes in the wiring note applies to the panel in front of you, because the printed pins 1, 3, 5 and 6 are underneath it. Then confirm the address.
\begin{shellcode}
sudo apt-get install -y i2c-tools python3-smbus
sudo i2cdetect -y 1          # must show 53
\end{shellcode}

\item \textbf{Read the sensor and compute the two angles.} The register sequence is the one from P01: 0x2D to enter measure mode, 0x31 for full resolution, then six bytes from 0x32.
\begin{pycode}
import math, time, smbus
bus = smbus.SMBus(1); A = 0x53
bus.write_byte_data(A, 0x2D, 0x08)   # measure
bus.write_byte_data(A, 0x31, 0x0B)   # full-res +/-16g

def g(lo, hi):
    v = (hi << 8) | lo
    if v & 0x8000:
        v -= 65536
    return v / 256.0

while True:
    b = bus.read_i2c_block_data(A, 0x32, 6)
    ax, ay, az = g(b[0], b[1]), g(b[2], b[3]), g(b[4], b[5])
    pitch = math.degrees(math.atan2(ax, math.sqrt(ay*ay + az*az)))
    roll  = math.degrees(math.atan2(ay, math.sqrt(ax*ax + az*az)))
    draw_two_bars(pitch, roll)       # on the framebuffer
    time.sleep(0.05)
\end{pycode}

\item \textbf{Draw two bars on the framebuffer.} One bar for pitch, one for roll, each with a marked centre and marked ends. The panel is 480x320, so a horizontal bar per angle across most of the width leaves room for a numeric readout if you want one.

\item \textbf{Check that nothing else is running.} No MQTT client, no publisher, no uplink of any kind. That is part of the acceptance test rather than a tidiness preference.
\begin{shellcode}
pgrep -a mosquitto_sub mosquitto_pub || echo "nothing publishing, good"
\end{shellcode}

\item \textbf{Tear the LCD down} when the lab is accepted, so the Pi 3 is free for P06 in a later sitting.
\end{steps}

\subsection*{Acceptance test}

\begin{itemize}
\item Rotate the assembly 90\textdegree{} on one axis: that bar travels full scale.
\item The other bar stays near zero during the same rotation.
\item No MQTT client is running.
\end{itemize}

\subsection*{Practices}

Keep the instrument offline for the whole sitting. The value of this lab is that a display and a sensor can be a complete product with no network behind them, and adding an uplink to see the numbers somewhere else removes the only thing this chapter owns.

One 40-pin HAT per host, as everywhere in this book. The 3.5 inch LCD is the HAT here, so the Pi 3 cannot also carry the MCC 118, a cellular HAT or the Explorer700, and it cannot be running P06 at the same time. Power off before unseating the panel, and put the ADXL345 back in the bin with its jumpers rather than leaving it dangling from a stored board.

The sensor is the same SEN0032 used in P01 and P09, and it is the same address and the same register sequence in all three. When it appears at 0x53 here, that is a part behaving as it has behaved twice already, and when it does not, the fault is in how the bus was reached rather than in the chip.

\subsection*{Pitfalls}

\begin{itemize}
\item \textbf{Following a 2016 flexfb recipe.} On Bookworm the Waveshare 3.5 inch (A) uses the current vendor overlay. The old recipe produces a panel that stays dark and a stack of changes to undo.
\item \textbf{Writing drawing code before checking for a framebuffer.} Confirm \texttt{/dev/fb0} or a DRM card first, exactly as P06 does.
\item \textbf{Long jumpers to the sensor.} The assembly gets rotated during its own test. A sensor swinging on loose leads measures its own swing.
\item \textbf{5 V on the accelerometer.} The SEN0032 goes to pin 1, never pin 2.
\item \textbf{Assuming a pass-through.} The source says one usually exists, and this panel has none, which is what blocked the lab. Read the board's own pinout before planning any wiring that depends on reaching a pin underneath it.
\item \textbf{Trying to run P06 on the same Pi in the same sitting.} Both labs want the LCD on the 40-pin header. Tear one down first.
\item \textbf{Leaving a subscriber running from an earlier lab.} The acceptance test says no MQTT client is running, and a forgotten \texttt{mosquitto\_sub} in another terminal is still a client.
\item \textbf{Using the single-argument arctangent.} The sign of the angle is then lost across half the range, and one bar travels the wrong way.
\end{itemize}

\subsection*{Sources}

\begin{itemize}
\item Waveshare 3.5 inch RPi LCD (A) wiki, and the current note for Bookworm, \url{https://www.waveshare.com/wiki/3.5inch_RPi_LCD_(A)}
\item DFRobot SEN0032 ADXL345 product wiki, \url{https://wiki.dfrobot.com/}
\item Analog Devices ADXL345 data sheet, register map 0x2D, 0x31, 0x32 and the SDO address strap
\item Raspberry Pi device tree and overlay documentation, \url{https://www.raspberrypi.com/documentation/computers/configuration.html}
\item Linux framebuffer and DRM documentation, for the difference between \texttt{/dev/fb0} and a DRM card
\item \texttt{i2c-tools} manual pages, \texttt{i2cdetect(8)}
\item Python \texttt{math} module documentation, \texttt{atan2} and its sign convention
\end{itemize}


\part{Rails, sidecars and the fleet}
\project{15}{POW-BB mixed-voltage discipline}{SBC-POW-BB + R/Y/G LEDs + any Pi or NanoPi host}{the electrical safety lab}

\begin{unique}
Owns the electrical safety lab. No radio, no IMU. Every other lab in this book assumes the rails are labelled and the logic level is settled; this is where that happens.
\end{unique}

\subsection*{Intent}

Label the rails. Prove that a 5 V LED circuit does not share a data pin with a 3.3 V system-on-chip input. The POW-BB is a labelled rail splitter, so the product of this lab is a bench where the 5 V rail and the 3.3 V rail are told apart by reading, not by guessing, and where the three numbers that say so are written in the lab book.

Red, yellow and green become a semaphore driven from 3.3 V GPIO, through the resistors already fitted on the LED module or through a series resistor on the breadboard. The lab sequence puts this lab first and says do not skip it. The reason is arithmetic rather than caution: Pi, NanoPi and Nucleo I/O are 3.3 V, and a 5 V pull-up sits above that limit whatever the rest of the circuit does.

\begin{kit}
Any Pi or NanoPi host, SBC-POW-BB, breadboard, jumpers, red, yellow and green LEDs.
\end{kit}

\begin{note}{Where this lab sits in the order}
The lab sequence in the appendix opens with P15 rails, and adds: do not skip. Everything after it assumes a labelled 5 V rail, a labelled 3.3 V rail and a ground that is common to the host, the breadboard and the supply. The jig of P14 and the sidecars of P16 and P17 are the next three labs, and each of them puts 3.3 V logic on a wire you can reach.
\end{note}

\subsection*{System architecture}

\diagram{p15_arch}{Two domains that meet only at ground. The power domain carries no data, and the data domain carries no 5 V. The multimeter is the only instrument in the lab, and it is the acceptance test.}

There are two domains on this bench and exactly one point where they meet. The power domain runs from the supply through the POW-BB into a labelled 5 V rail and a labelled 3.3 V rail. The data domain runs from a blink script through \texttt{gpiochip0} to three header pins. The rails feed LED anodes, the header pins sink LED cathodes, and the only shared conductor is ground.

Read the architecture figure as a statement about what is missing. Nothing connects the 5 V rail to a header pin. That absence is the lab, and the negative test at the end of the procedure is there to show what the absent connection would have measured.

There is no microcontroller in this chapter, which makes it the plainest example of the rule that Linux owns the system. The host holds the script, the state and the record. Every later lab adds a peripheral to this picture, whether it is a modem on USB, a Nucleo on USB-CDC or an ESP on a UART, and none of them changes the voltage that arrives at a header pin.

\subsection*{Wiring and schematic}

\diagram{p15_schematic}{The rails, the semaphore and the negative test. LED anodes sit on the 3.3 V rail through a series resistor and the cathodes sink into GPIO, so a pin driven low lights its LED. The dashed fragment at the bottom is the negative test, and it lives on the breadboard.}

\begin{table}[H]\centering\small
\begin{tabular}{p{40mm}p{24mm}p{20mm}p{62mm}}
\toprule
Signal & Header pin & BCM GPIO & Note \\
\midrule
POW-BB 5 V rail & none & none & Measured, labelled, and connected to no header pin \\
POW-BB 3V3 rail & none & none & From the board regulator, not a guess. Feeds the LED anodes \\
Ground & 6 & GND & Common to POW-BB, breadboard and host \\
Green LED cathode & 13 & GPIO27 & Sinks when the pin is driven low \\
Yellow LED cathode & 15 & GPIO22 & Sinks when the pin is driven low \\
Red LED cathode & 16 & GPIO23 & Sinks when the pin is driven low \\
Series resistor & none & none & On the breadboard, unless the LED module already carries one \\
\bottomrule
\end{tabular}
\caption{Wiring for a Raspberry Pi host. The GPIO numbers are the source's BCM 27/22/23; the pin numbers are the matching positions on the 40-pin header.}
\end{table}

On a NanoPi NEO Air the same circuit uses pin 1 SYS\_3V3 and pin 6 GND on the 24-pin header. The NEO Air GPIO chosen for each cathode is not fixed by this book: pick three free pins and confirm them against the silkscreen before wiring. The one number that does not change is the logic level. The NEO Air is an Allwinner H3 board with 3.3 V I/O, and its 3.3 V pin is named SYS\_3V3 rather than 3V3, which is the kind of difference that a labelled rail prevents from becoming a mistake.

Two details in the schematic carry the whole chapter. The first is the direction of the LED: the anode is on the 3.3 V rail and the cathode goes to the pin, so the pin sinks current and a low output is the lit state. The second is that the 5 V rail ends at a probe and goes nowhere else. Neither is difficult to build. Both are easy to build the other way round by accident, which is why they are drawn.

\begin{asciiart}
PSU 5V ---- POW-BB 5V rail   (label it)
POW-BB 3V3 rail labelled     (from the board regulator, not a guess)
GND common
LED anodes on 3V3 via series R, cathodes to GPIO as sinks
  OR LED modules already 3.3 V tolerant on the data pin
Never tie a 5V pull-up to a Pi GPIO.
\end{asciiart}

\subsection*{Bench layout}

\diagram{p15_bench}{Bench layout. The POW-BB sits between the supply and the breadboard, the host contributes only ground and three sinking pins, and the multimeter stays on the bench for the whole sitting.}

Put the POW-BB between the supply and the breadboard, not beside the host. That placement is a reminder of where the current comes from. Keep the multimeter out for the whole sitting: the three rail readings are taken at the start, and the negative test at the end uses the same probes without rewiring anything on the host side.

The host contributes two things to this bench and nothing else: a ground and three pins that can sink current. Nothing on the host side of the bench is a power source in this lab, and no jumper runs from the 5 V rail toward the header. A bench that looks like this one is a bench where the next lab can be wired without rechecking the rails.

\subsection*{Software design (UML)}

\diagram{p15_uml}{The procedure as an activity. The negative test is a branch that ends in a dismantle step, never in a connection to a header pin.}

The software in this lab is three output pins and a loop, so the interesting design is the procedure rather than the program. Measure, label, wire, blink, then run the negative test and take it apart. The dismantle step is part of the procedure and not an afterthought: a pull-up left on the breadboard is a 5 V source sitting within jumper reach of a 3.3 V input for the rest of the week.

The blink script itself sets every pin high before it does anything else, because high is the unlit state when the cathode sinks. It also restores the pins on the way out. A semaphore that is left in an arbitrary state at the end of a run is a semaphore the next lab cannot read, and the same three colours carry meaning in P01, P02 and P03, where green means registered, yellow means a default route and red means an event.

\subsection*{Data flow (ASCII)}

The current path and the information path are drawn together below. Current leaves the 3.3 V rail, passes the series resistor and the LED, and returns through the pin into the common ground. Information travels the other way: the state of the script decides whether the pin sinks, and the colour on the bench is the only output of the program.

\begin{asciiart}
   PSU 5 V
      |
      v
  +------------------+   5V rail  ---> multimeter probe only, no header pin
  |   SBC-POW-BB     |
  |  rail splitter   |   3V3 rail ---> series R ---> LED anode
  +------------------+                                  |
      |                                                 | cathode
      | GND                                             v
      |                                        Pi header pin 13 / 15 / 16
      |                                                 |
      +--- common ground -----------------------------> + GPIO low = LED on
                                                          GPIO high = LED off

  Negative test (breadboard only):
      5V rail --- pull-up R --- test node --- multimeter --- GND
      test node ---> a GPIO:  never.  Read it, then take it apart.
\end{asciiart}

\subsection*{Steps}

\begin{steps}
\item \textbf{Measure and write down three numbers.} With the multimeter: the 5 V rail is 5 V, the 3.3 V rail is 3.3 V, and GND is common. Write the three numbers in the lab book. A rail that reads something else is a wiring question to answer now, not after the LEDs are fitted.
\begin{plaincode}
lab book, P15
  5V rail        measured   ____ V
  3V3 rail       measured   ____ V
  GND continuity POW-BB to breadboard to host   ____
  note: no 5 V on header pins
\end{plaincode}
The ground line is a continuity check rather than a voltage: the POW-BB ground, the breadboard ground rail and the host ground are one node or the rest of the lab means nothing.

\item \textbf{Blink the semaphore from 3.3 V logic.} Drive green, yellow and red from BCM 27/22/23 on a Pi, without ever connecting those GPIOs to the 5 V rail. The anodes sit on the 3.3 V rail through the series resistor, so the pin sinks the current and a low output is the lit state.
\begin{pycode}
import time
import RPi.GPIO as GPIO

PINS = {"grn": 27, "yel": 22, "red": 23}      # BCM numbering, cathodes sink

GPIO.setmode(GPIO.BCM)
for p in PINS.values():
    GPIO.setup(p, GPIO.OUT, initial=GPIO.HIGH)   # HIGH = no sink = LED off

try:
    while True:
        for name, p in PINS.items():
            GPIO.output(p, GPIO.LOW)             # LOW = sink = LED on
            time.sleep(0.4)
            GPIO.output(p, GPIO.HIGH)
finally:
    GPIO.cleanup()
\end{pycode}

\item \textbf{Run the negative test on the breadboard only, not on a GPIO.} Build a 5 V pull-up on the breadboard, measure it, and show that it sits above 3.3 V. Then take it apart. Nothing in this step touches a header pin, and the pull-up does not stay on the board afterwards.
\end{steps}

\begin{rulebox}{The 3.3 V rule in the rest of the book}
\begin{itemize}
\item Pi, NanoPi and Nucleo I/O are 3.3 V. The POW-BB is a labelled rail splitter.
\item Never 5 V TTL into ESP pins (P14).
\item Do not put 5 V onto SYS\_3V3 on the NanoPi NEO Air (P09).
\item Never exceed $\pm$10.1 V on an MCC 118 input, and Pi GPIO is not an analog source (P05).
\item Never source a Cat-4 power amplifier from the PPK2 (P10).
\end{itemize}
\end{rulebox}

\subsection*{Acceptance test}

\begin{itemize}
\item Three measured rail voltages are written in the lab book: the 5 V rail, the 3.3 V rail and a ground that is common to the POW-BB, the breadboard and the host.
\item A blink script runs the three LEDs as a semaphore from BCM 27/22/23.
\item The lab book carries a sentence that says "no 5 V on header pins".
\item The breadboard is clear of the negative-test pull-up when the sitting ends.
\item No jumper runs between the 5 V rail and any header pin, on any host, at any point in the lab.
\end{itemize}

\subsection*{Practices}

The source gives this lab no practices paragraph of its own, so the practices are its own sentences. Label the rails on the POW-BB and treat the label as the instrument: it is a rail splitter, not a supply with a display. Take the three measurements before the LEDs go in, because a mislabelled rail found afterwards means unwiring work that was already done.

Keep every I/O pin on this bench at 3.3 V: Pi, NanoPi and Nucleo I/O are 3.3 V, and that number is the same in every other lab in the book. The negative test exists to be seen once and then removed. Build one lab at a time, and do not skip this one, which is why it is first in the lab sequence.

\subsection*{Pitfalls}

\begin{itemize}
\item \textbf{A 5 V pull-up on a GPIO.} The symptom ranges from an input that reads high and never changes to a pin that stops responding altogether. This is the one connection the lab exists to prevent.
\item \textbf{An unlabelled rail.} Two red jumpers, two rails, one guess. Label the rails and the guess disappears.
\item \textbf{No series resistor.} An LED straight across a rail and a pin draws whatever the rail can deliver. Confirm whether the LED module already carries a resistor before adding another.
\item \textbf{Ground not common.} If the POW-BB ground and the host ground are separate, the LED either stays dark or lights in a way that does not follow the script, and the measured rail numbers are meaningless.
\item \textbf{LED fitted the wrong way round.} With anodes on 3.3 V and cathodes on GPIO, a reversed LED never lights and looks like a software fault.
\item \textbf{The negative test left standing.} A 5 V pull-up that stays on the breadboard is a hazard to the next lab that borrows a jumper from the same row.
\item \textbf{Two red jumpers from the same strip.} Colour is a convention, not a measurement. The label on the POW-BB and the reading in the lab book are what settle which rail a wire came from.
\item \textbf{Assuming the semaphore is inverted software.} With cathodes on GPIO, low is lit. A script written for the opposite convention runs perfectly and shows the wrong state on every LED.
\end{itemize}

\subsection*{Sources}

\begin{itemize}
\item JOY-iT SBC-POW-BB breadboard power supply documentation, \url{https://joy-it.net/en/products/SBC-POW-BB}
\item Raspberry Pi hardware documentation, 40-pin header and GPIO electrical limits, \url{https://www.raspberrypi.com/documentation/computers/raspberry-pi.html}
\item FriendlyElec NanoPi NEO Air wiki, 24-pin header pinout, \url{https://wiki.friendlyelec.com/wiki/index.php/NanoPi_NEO_Air}
\item RPi.GPIO documentation for the BCM numbering used by the blink script, \url{https://sourceforge.net/projects/raspberry-gpio-python/}
\item libgpiod and the character-device GPIO interface, for hosts where \texttt{RPi.GPIO} is not available, \url{https://git.kernel.org/pub/scm/libs/libgpiod/libgpiod.git/}
\item JOY-iT LED module documentation, to confirm whether a series resistor is already fitted, \url{https://joy-it.net/en/}
\end{itemize}

\project{14}{USB-TTL provisioning jig}{Raspberry Pi 3 + Renkforce USB-TTL + ESP32 and ESP8266-PROG}{the Renkforce cable as a factory tool}

\begin{unique}
Owns the Renkforce cable as a factory tool, plus both ESP boards as targets. Not a sensor lab. The ESP32 becomes a product in P16 and the ESP8266 becomes a modem in P17; here they are units passing through a jig.
\end{unique}

\subsection*{Intent}

The Pi 3 flashes and AT-provisions the ESP32 and the ESP8266 without either board becoming the product. The output of the lab is not a running device, it is a repeatable procedure and a written record: a chip that answers, an image that was written deliberately, and a checklist line that says which unit it was.

The jig is a Linux host, a 3.3 V serial cable and one target at a time. That last clause is the discipline of the chapter. The Renkforce cable is one UART path and the ESP8266-PROG carries its own USB-UART, so a bench that has both connected has two drivers for one set of pins.

\begin{kit}
Pi 3, Renkforce USB/TTL cable, SBC-NodeMCU-ESP32, SBC-ESP8266-PROG, red, yellow and green LEDs, breadboard.
\end{kit}

\begin{note}{One UART path}
The ESP8266-PROG carries an onboard USB-UART. The Renkforce cable is a second one. Connect one of them at a time, never both at once. The lab sequence puts this lab second, after the rails of P15 and before the sidecars of P16 and P17, because both of those labs start with a board that was provisioned here.
\end{note}

\subsection*{System architecture}

\diagram{p14_arch}{The jig. One Linux host, one serial character device, one 3.3 V cable, and one target at a time. The ESP8266-PROG carries a second USB-UART of its own, drawn as the path that stays unplugged.}

Linux owns the system here in the plainest possible way: the target has no storage, no clock and no network of its own during a provisioning run, and every artifact of the run lands in a file on the Pi. The tools are \texttt{esptool.py} for the flash side and \texttt{minicom} for the AT or console side, and both of them speak to the same character device.

Both tools are on the Linux side of the boundary, and so is every file the run produces. The target contributes an answer and nothing else. That is the same boundary the rest of the book draws between a Linux host and a microcontroller, stated here in its smallest form: the microcontroller is a unit under test, and the host is the instrument.

The kernel binds a USB serial driver to the Renkforce cable and presents \texttt{/dev/ttyUSB0}. Confirm which driver claims it with \texttt{dmesg} on the first plug: the fleet script in P20 looks for the same class of adapter by its USB vendor identifiers, \texttt{1a86} and \texttt{10c4}. When two adapters are plugged in at once the numbering is no longer predictable, which is the practical reason for one target at a time.

\subsection*{Wiring and schematic}

\diagram{p14_schematic}{Three signal wires to the target, with transmit crossed to receive. The red lead is 5 V and stays out of the circuit; the target takes its power from its own USB socket. GPIO0 is held low for an ESP32 download and released for a normal boot.}

\begin{table}[H]\centering\small
\begin{tabular}{p{34mm}p{34mm}p{26mm}p{52mm}}
\toprule
Renkforce USB-TTL & ESP target & Direction & Note \\
\midrule
Red & nothing & none & 5 V from the host USB rail. Leave it unconnected \\
Black, GND & GND & common & Fit this wire before the others, remove it last \\
Green, TX & RX & to the target & Transmit crossed to receive \\
White, RX & TX & from the target & Receive crossed to transmit \\
Target power & its own USB & from the bench & Not from this cable \\
(none) & GPIO0 & held low & ESP32 download mode only \\
(none) & EN & reset & Release GPIO0, then reset to run the image \\
\bottomrule
\end{tabular}
\caption{The three-wire mapping, plus the two ESP32 pins that select download mode. The cable is a sealed PL2303HX with four flying leads and no voltage selector of any kind, so there is nothing to set before connecting; what there is to do is leave the red lead alone. The ESP8266-PROG has an onboard USB-UART: use one UART path, not both at once.}
\end{table}

\begin{asciiart}
USB-TTL red    -> nothing        5 V from the host USB rail
USB-TTL black  -> ESP GND        fit first, remove last
USB-TTL green  -> ESP RX         TX crossed to RX
USB-TTL white  <- ESP TX         RX crossed to TX
ESP power      <- its own USB socket, not this cable
ESP32: hold GPIO0 low for download, EN reset
ESP8266-PROG: onboard USB-UART may already exist, so use ONE
              UART path, not both at once
\end{asciiart}

\begin{note}{What this cable is, and the one thing about it that is not measured}
It is a sealed moulding over a PL2303HX with four flying leads: red is 5 V taken straight from the host USB rail, black is ground, green is the adapter's transmit and white its receive. There is no jumper, no switch and no voltage marking on the shell, so an instruction to select 3.3 V on it cannot be followed. An earlier draft of this chapter carried one, along with a pitfall about leaving it at 5 V, and both described a part that does not exist.

The signal lines are documented as 3.3 V logic, by the source this volume follows and by the vendor's own description. Neither is a measurement, and this bench has no voltmeter to take one, so the claim is recorded as a claim. It matters in one direction only: a transmit line that is really 5 V damages whatever 3.3 V receive pin it meets.

That is a good reason to prefer the path that does not depend on it. Both targets carry their own USB-UART, so flash and converse over the board's own socket and keep this cable for a board that has none, which in this book is the NanoPi of P09. When the cable is the only way in, fit ground first and accept the documented level knowingly rather than by not having asked.
\end{note}

The LK-LED10 modules sit on the breadboard as jig status, driven from the host exactly as the semaphore of P15 is driven: the header pin feeds the module's signal pin, which is its anode side, and the module returns to the common ground rail, so a pin driven high is the lit state and each module's resistor is already inside it. The source gives no dedicated pin assignment for this lab, so reuse the P15 wiring rather than inventing a second one, and settle the polarity on one module first, as P15 does, because which way round the LED sits is not printed on the part.

\subsection*{Bench layout}

\diagram{p14_bench}{Bench layout. One target is wired to the Renkforce cable. The other target waits on the bench, and the ESP8266-PROG USB cable stays out of the Pi while the TTL path is in use.}

Keep the jig small: the Pi, the cable, the breadboard and one target. Put the second target on the bench, not on the breadboard, so that the question "which board answered" never has to be asked. The ESP8266-PROG USB cable belongs beside its board, unplugged, whenever the Renkforce cable is the active path.

The Pi 3 wears no 40-pin HAT in this lab, so the header is free for the semaphore. That is worth noticing, because it is what makes this jig portable: the same procedure runs on whichever Linux host is free, and nothing about it competes for the exclusive header that P02, P05, P06, P11 and P19 each need.

\subsection*{Software design (UML)}

\diagram{p14_uml}{A provisioning run. It begins with an identification read and ends in the checklist file, not in a running application.}

A provisioning run is a fixed sequence with a written result. Identify the chip, write an image you can rebuild, reset the board, confirm it boots, and record the unit. The checklist is the deliverable: the MAC address, the flash size, a unique client identifier and the fact that the Wi-Fi credentials were written to NVS. A run that ends without a checklist line has produced a board nobody can tell apart from the next one.

The last exchange in the diagram is the one that makes the run a test rather than a hope. Pull the power, let the board come back on its own, and confirm that it boots the image that was written. A target that only runs while the serial cable is attached has not been provisioned, it has been driven.

\subsection*{Data flow (ASCII)}

\begin{asciiart}
  Pi 3                                Renkforce USB/TTL       one target at a time
  +-------------------------+         +-----------------+     +--------------------+
  | esptool.py / minicom    |         | PL2303HX, sealed|     |  ESP32 NodeMCU     |
  |   |                     |  USB    | 3.3 V logic     |green-> RX                |
  |   v                     |<=======>| (documented,    |white<- TX                |
  | /dev/ttyUSB0            |         |  not measured)  |black-- GND GPIO0 low     |
  |   |                     |         | red: 5 V, unused|     |  EN reset          |
  |   v                     |         +-----------------+     +--------------------+
  | p14_checklist.txt       |                                 power from its own USB
  |   MAC, flash size,      |         ESP8266-PROG carries its own USB-UART.
  |   client_id, NVS ok     |         Use ONE UART path, never both at once.
  +-------------------------+
\end{asciiart}

\subsection*{Steps}

\begin{steps}
\item \textbf{Install the tools and identify each chip.} One target at a time on \texttt{/dev/ttyUSB0}.
\begin{shellcode}
sudo apt-get install -y python3-serial esptool minicom
esptool.py --chip esp32 --port /dev/ttyUSB0 chip_id
esptool.py --chip esp8266 --port /dev/ttyUSB0 chip_id
\end{shellcode}

\item \textbf{Flash a known AT or blink image.} Confirm the yellow LED script you wrote, not a random vendor blob you cannot reset. For the ESP32 the sequence is: hold GPIO0 low, reset with EN, run the write, release GPIO0, reset again.

When the image is an AT firmware, the same cable is the console. Open it and confirm that the board answers before the unit leaves the jig. The join sequence itself belongs to P17, where the ESP8266 is an AT modem rather than a unit under test.
\begin{atcode}
sudo minicom -D /dev/ttyUSB0 -b 115200
AT
\end{atcode}

\item \textbf{Record a provisioning checklist.} MAC, flash size, a unique \texttt{client\_id}, and Wi-Fi credentials written to NVS. One line per unit, in a file on the Pi.
\begin{plaincode}
unit   chip      mac                 flash    client_id      nvs_wifi
0001   esp32     xx:xx:xx:xx:xx:xx   ____     p16-led-01     written
0002   esp8266   xx:xx:xx:xx:xx:xx   ____     p17-at-01      written
\end{plaincode}
\end{steps}

The four fields are not arbitrary. Each of them answers a question that comes up later:

\begin{description}
\item[MAC] the only identifier the board carries before you give it one. It is what a broker log or an access point list will show.
\item[Flash size] what \texttt{esptool.py} read from the chip, not what a listing claimed. An image that does not fit is a flash that was assumed rather than measured.
\item[client\_id] the unique name the unit will use on a broker. P16 subscribes an ESP32 to \texttt{lab/p16/led/\#}, and two boards with one identifier make that topic unreadable.
\item[NVS Wi-Fi] whether credentials were written to non-volatile storage, so the board joins without a console attached.
\end{description}

\subsection*{Acceptance test}

\begin{itemize}
\item Both chips report \texttt{chip\_id} through \texttt{esptool.py}, each on its own run.
\item A second run after a power pull still boots the image you flashed.
\item The checklist file has one line per unit, with a MAC address, a flash size and a unique client identifier.
\item Only one UART path is connected while either of those runs is in progress.
\end{itemize}

\subsection*{Practices}

3.3 V logic. Never 5 V TTL into ESP pins, and note that this cable offers nothing to set: its level is fixed by its internal design and its red lead is 5 V whatever else is true, so the discipline is to leave that lead unconnected and fit ground first. Flash an image you can rebuild, so that a board can always be returned to a known state. Keep the two targets apart on the bench: the ESP32 goes on to be a peer in P16 and the ESP8266 goes on to be an AT slave in P17, and a mislabelled unit turns both of those labs into guesswork.

\subsection*{Pitfalls}

\begin{itemize}
\item \textbf{Two UART paths at once.} The Renkforce cable and the onboard USB-UART of the ESP8266-PROG drive the same pins. The symptom is a sync that fails at random, or a console that prints half the characters.
\item \textbf{The red lead connected to anything.} It is 5 V from the host USB rail, and the one thing on this cable that can damage a 3.3 V part. There is no jumper to get wrong, so this is the connection to leave out rather than the setting to check.
\item \textbf{Powering an ESP32 from the cable.} Even at the right voltage it is the wrong source: an association peak draws more than a serial cable's supply lead is meant to carry, and the symptom is a brownout in the middle of a write rather than a clean failure.
\item \textbf{Transmit wired to transmit.} Nothing is damaged and nothing answers. \texttt{esptool.py} reports a failed connection and the cause is the pair of wires that were not crossed.
\item \textbf{GPIO0 not held low.} The ESP32 boots the existing image instead of the serial bootloader, so the write never starts.
\item \textbf{Two adapters plugged in.} \texttt{/dev/ttyUSB0} is then whichever enumerated first. Unplug one, or read \texttt{dmesg} and use the right node.
\item \textbf{A vendor blob nobody can rebuild.} The board works until it does not, and there is no known image to return to.
\item \textbf{No checklist line.} Two identical modules with no record is the same as one unknown module.
\item \textbf{Ground fitted last.} Signal wires connected before the ground reference produce failures that look intermittent. Fit ground first, remove it last.
\item \textbf{Treating a target as a product.} The ESP32 becomes a peer in P16 and the ESP8266 becomes an AT slave in P17. Neither role starts here, and a jig that grows an application stops being a jig.
\end{itemize}

\subsection*{Sources}

\begin{itemize}
\item Espressif esptool documentation, chip\_id and flash commands, \url{https://docs.espressif.com/projects/esptool/}
\item Espressif ESP32 boot mode selection, GPIO0 and EN, \url{https://docs.espressif.com/projects/esptool/en/latest/esp32/advanced-topics/boot-mode-selection.html}
\item JOY-iT SBC-NodeMCU-ESP32 product page and manual, \url{https://joy-it.net/en/products/SBC-NodeMCU-ESP32}
\item JOY-iT SBC-ESP8266-PROG product page and manual, \url{https://joy-it.net/en/products/SBC-ESP8266-PROG}
\item Renkforce USB to TTL serial cable documentation, 3.3 V logic level
\item Linux USB serial drivers, \url{https://www.kernel.org/doc/html/latest/usb/usb-serial.html}
\end{itemize}

\project{16}{ESP32 Wi-Fi/BLE sidecar}{Raspberry Pi 4 + SBC-NodeMCU-ESP32}{the sidecar pattern and LED actuation}

\begin{unique}
Owns the sidecar pattern: Linux keeps Ethernet and LTE, the ESP32 keeps the 2.4 GHz radio personality. The actuator LEDs from the P12 broker live here.
\end{unique}

\begin{blueprint}
Draft P14 mixed the glass and the actuator. Split them. The panel, the keyboard and the broker stay in P12. The radio and the LEDs are this lab.
\end{blueprint}

\subsection*{Intent}

The ESP32 joins the Pi 4 access point or the LAN, subscribes to \texttt{lab/p16/led/\#}, and drives red, yellow and green. The Pi never bit-bangs those LEDs. That last sentence is the pattern: an actuator that is reached through a topic can be moved to the other side of the room without changing the publisher.

The division of labour is the teaching. Linux keeps the parts of the system that need storage, time and a routable network: the broker, the log, and whatever wide-area path P01 or P03 provided. The ESP32 keeps the 2.4 GHz radio personality and three output pins. Neither side reaches into the other, and the contract between them is a topic name.

\begin{kit}
Pi 4, with the P12 broker welcome but not required, SBC-NodeMCU-ESP32, SBC-POW-BB, red, yellow and green LEDs, USB for power and serial.
\end{kit}

\begin{note}{What the ESP32 arrives with}
The board that lands on this bench came through the jig of P14, so it already has a known image, a recorded MAC and a unique client identifier. Provision it there, not here. In P17 the other ESP board takes the opposite role, an AT slave on a UART rather than a peer on the LAN.
\end{note}

\subsection*{System architecture}

\diagram{p16_arch}{The sidecar. Linux keeps Ethernet and the cellular path, the ESP32 keeps the 2.4 GHz radio, and the only thing crossing between them is an MQTT topic. The path from the Pi header to the LEDs is the one that is deliberately absent.}

There are two radios in this picture and they do not compete. Whatever wide-area or wired path the Pi holds stays with Linux, and the 2.4 GHz personality lives on the ESP32. A publisher on the Pi writes to a topic; a subscriber on the ESP32 reads it and moves a pin. The broker in the middle is either the Mosquitto instance of P12 or a public test broker, and nothing in the design changes when that choice changes.

Note what the figure does not contain: a wire from the Pi 40-pin header to the LEDs. The Pi could drive them, as it does in P15, and that is exactly the design this lab rejects. Once the actuator is behind a topic it can be re-homed, duplicated or replaced by a second ESP32 without the publisher learning anything new.

\subsection*{Wiring and schematic}

\diagram{p16_schematic}{The sidecar wiring. USB carries power and the serial console, the POW-BB 3.3 V rail feeds the LED anodes through series resistors, and three ESP32 outputs sink the cathodes. The pin choice is yours: confirm it against the silkscreen.}

\begin{table}[H]\centering\small
\begin{tabular}{p{38mm}p{30mm}p{80mm}}
\toprule
Signal & Where & Note \\
\midrule
USB & ESP32 to the Pi & Power and serial console in one cable \\
POW-BB 3V3 rail & LED anodes & Through a series resistor, as in P15 \\
Ground & Common & POW-BB, breadboard and ESP32 \\
Three ESP32 outputs & LED cathodes & Three free GPIOs, confirm them against the silkscreen \\
Wi-Fi & 2.4 GHz & The ESP32 radio, not the Pi radio \\
\bottomrule
\end{tabular}
\caption{The sidecar has four wires and one radio link. The source names three GPIOs but not which three: pick them on the board in front of you and record the choice with the unit.}
\end{table}

The source for this lab gives a kit and a behaviour, not a pin map. Take the electrical convention from P15 without inventing anything else: anodes on the 3.3 V rail through a series resistor, cathodes sinking into an output pin, one common ground. ESP pins are 3.3 V, which is the same rule P14 states for the jig.

\begin{asciiart}
POW-BB 3V3 rail --- series R --- LED anode
LED cathode     --- ESP32 output pin        (three of them, low = lit)
POW-BB GND      --- ESP32 GND               (common)
ESP32 USB       --- Pi USB-A                (power and serial console)
Wi-Fi           --- 2.4 GHz to the broker   (the ESP32 radio, not the Pi radio)
Never 5 V into an ESP pin.
\end{asciiart}

\subsection*{Bench layout}

\diagram{p16_bench}{Bench layout. The ESP32 sits on the breadboard with the LEDs and the POW-BB rail. The only connections to the Pi are a USB cable for power and console, and a radio link that carries the actual work.}

Put the ESP32 far enough from the Pi that the USB cable is doing something visible. The point of the layout is that the cable could be replaced by any 5 V source and the lab would still work, because the traffic arrives over the air. If the P12 broker is on the glass beside the bench, the publish and the LED are visible in one glance, which is the fastest way to see the latency figure.

\subsection*{Software design (UML)}

\diagram{p16_uml}{One publish, end to end. The budget is 200 ms from the command on the Pi to the pin moving on the ESP32, on a quiet LAN.}

The sequence has four hops and one deadline. A publish enters the broker, the broker delivers to the subscriber, the subscriber parses one byte of payload, and the pin moves. Everything in that chain is small, so when the number is missed the cause is almost always the network rather than the code: a busy channel, a retry, or an access point that put the client to sleep.

Keep the payload trivial. The topic carries the colour and the message carries the state, exactly as the source's publish line shows. A subscriber that has to parse JSON to turn on a light has been given a second job it does not need.

\subsection*{Data flow (ASCII)}

\begin{asciiart}
  Pi 4  (the P12 broker is welcome)        2.4 GHz         ESP32 sidecar
  +-------------------------------+                    +--------------------+
  | mosquitto_pub -t              |                    | Wi-Fi STA          |
  |   lab/p16/led/red -m 1        |                    | MQTT client        |
  |        |                      |                    |   sub lab/p16/led/#|
  |        v                      |    TCP over        |        |           |
  | mosquitto broker :1883        |====  Wi-Fi  ======>|        v           |
  |                               |                    | three output pins  |
  | eth0 / usb0 stay with Linux   |                    +--------|-----------+
  +-------------------------------+                             v
                                                        R / Y / G LEDs
  budget: publish to LED flip in under 200 ms on a quiet LAN
  the Pi 40-pin header drives nothing in this lab
\end{asciiart}

\subsection*{Steps}

\begin{steps}
\item \textbf{Write the sidecar sketch.} An ESP32 Arduino or ESP-IDF sketch: Wi-Fi STA, MQTT client, three GPIOs. Subscribe to \texttt{lab/p16/led/\#} and map the last topic element to a pin. If the broker is the P12 instance, it has \texttt{allow\_anonymous false} and a password file, so the client carries a user and a password.
\begin{ccode}
/* outline only: Wi-Fi STA, MQTT client, three outputs */
void on_message(const char *topic, const char *payload, int len)
{
    int pin = pin_for_colour(topic);   /* lab/p16/led/red -> the red pin */
    if (pin < 0)
        return;
    gpio_set_level(pin, payload[0] == '1' ? 0 : 1);   /* low = lit */
}
\end{ccode}

\item \textbf{Publish from the Pi.} One topic, one byte of payload.
\begin{shellcode}
mosquitto_pub -h localhost -t lab/p16/led/red -m 1
\end{shellcode}

\item \textbf{Optional BLE.} The ESP32 advertises a lab UUID; the Pi uses \texttt{bluetoothctl} only as a presence check. Do not build a phone app.
\end{steps}

\subsection*{Acceptance test}

\begin{itemize}
\item A topic flip toggles the matching LED in under 200 ms on a quiet LAN.
\item All three colours respond to their own topic, and no other colour moves.
\item No process on the Pi holds a GPIO for these LEDs. The header is empty of them.
\item With BLE enabled, \texttt{bluetoothctl} sees the advertised lab UUID. Nothing else is built on top of it.
\end{itemize}

\subsection*{Practices}

The source gives this lab no practices paragraph of its own, so the practices are its own sentences. Linux keeps Ethernet and LTE; the ESP32 keeps the 2.4 GHz radio personality. The Pi never bit-bangs those LEDs. Keep the BLE work to a presence check and do not build a phone app. Measure the 200 ms figure on a quiet LAN, because that is the condition the acceptance test names, and record which broker answered.

Provision the board in P14 before it arrives here, and leave the other ESP board for P17. The two labs use the same family of silicon in opposite roles, and the difference between a peer and an AT slave is the thing worth keeping straight.

\subsection*{Pitfalls}

\begin{itemize}
\item \textbf{The Pi driving the LEDs "just to test".} It works, and it removes the entire point of the lab. The actuator is on the far side of a topic or it is not a sidecar.
\item \textbf{An anonymous client against the P12 broker.} That broker sets \texttt{allow\_anonymous false}. The symptom is a client that connects and is immediately dropped.
\item \textbf{A busy channel during the measurement.} The 200 ms budget is quoted for a quiet LAN. A scan, a large transfer or a crowded access point makes the number meaningless.
\item \textbf{Power saving on the station.} An ESP32 that is allowed to sleep between beacons answers late and inconsistently, which looks exactly like a broker problem.
\item \textbf{5 V into an ESP pin.} The same rule as P14 and P15: ESP pins are 3.3 V.
\item \textbf{A payload that grew.} One byte is enough. A JSON document for a light means the subscriber now has a parser to get wrong.
\item \textbf{Two boards with one client identifier.} The broker disconnects one when the other connects. This is why P14 records a unique identifier per unit.
\end{itemize}

\subsection*{Sources}

\begin{itemize}
\item Eclipse Mosquitto documentation, \texttt{mosquitto\_pub} and broker configuration, \url{https://mosquitto.org/documentation/}
\item Espressif ESP-IDF MQTT client component, \url{https://docs.espressif.com/projects/esp-idf/en/latest/esp32/api-reference/protocols/mqtt.html}
\item Espressif ESP-IDF Wi-Fi station and power-save modes, \url{https://docs.espressif.com/projects/esp-idf/en/latest/esp32/api-guides/wifi.html}
\item Arduino core for the ESP32, \url{https://github.com/espressif/arduino-esp32}
\item JOY-iT SBC-NodeMCU-ESP32 product page and manual, \url{https://joy-it.net/en/products/SBC-NodeMCU-ESP32}
\item BlueZ \texttt{bluetoothctl} documentation, for the presence check only, \url{http://www.bluez.org/}
\end{itemize}

\project{17}{ESP8266 AT modem on the NanoPi, onboard Wi-Fi off}{NanoPi NEO Air + SBC-ESP8266-PROG, AP6212 down}{the air-gap radio}

\begin{unique}
Owns the air-gap radio. The NanoPi's AP6212 is down and the ESP8266 speaks AT over UART1. Different from P16, where the ESP32 is a peer, not an AT slave.
\end{unique}

\subsection*{Intent}

Take the radio away from the Linux host and give it to a serial device. The AP6212 that the NEO Air carries for Wi-Fi and Bluetooth is switched off for the whole session, and the only network interface that reaches an access point is the one inside the ESP8266. The host reaches it through four wires and a command language.

The split is the product of the lab. When it is built correctly, \texttt{nmcli} on the NanoPi reports no Wi-Fi at all while \texttt{AT+CIFSR} on the same bench returns an IP address. An operating system with no radio and a radio with no operating system, joined by a UART: that is the shape a great many field devices actually have, and the acceptance test is the demonstration that the two halves really are separate.

\begin{kit}
NanoPi NEO Air with FriendlyElec Ubuntu, SBC-ESP8266-PROG, four jumpers to the 24-pin header, 3V3 from SYS\_3V3. The source names no other part for this lab.
\end{kit}

\begin{note}{Three serial ports and three jobs}
UART1 on pins 8 and 10 is the modem path in this lab. The 4-pin debug header of UART0 is the console, as it is in P09, and the ESP8266-PROG onboard USB-UART is the path from P14 that stays unplugged here. Three serial ports on one bench, each with one job.
\end{note}

\subsection*{System architecture}

\diagram{p17_arch}{The air gap. The AP6212 is switched off on the Linux side, and the IP address lives on the far side of a UART. Nothing in the host network stack knows the access point exists.}

Read this architecture next to P16 and the difference is the direction of authority. There the ESP32 is a peer on the LAN with its own subscription and its own opinion about when to move a pin. Here the ESP8266 is an AT slave: it answers commands, it holds an address the host cannot see in \texttt{ip addr}, and it originates nothing.

That has a consequence worth stating plainly. Anything the host wants to send leaves through a command, not through a socket. There is no route, no interface and no DNS resolver on the Linux side taking part in this path, so the usual tools report a machine with no network, which is exactly what the acceptance test asks you to show.

\subsection*{Wiring and schematic}

\diagram{p17_schematic}{Four wires and one switched-off radio. Transmit on pin 8 goes to the ESP receive pin, receive on pin 10 comes from the ESP transmit pin, and SYS\_3V3 on pin 1 is the supply. The AP6212 stays down for the whole session.}

\begin{table}[H]\centering\small
\begin{tabular}{p{36mm}p{28mm}p{22mm}p{58mm}}
\toprule
NEO Air 24-pin & Signal & ESP8266 & Note \\
\midrule
pin 1 & SYS\_3V3 & 3V3 & Supply. Do not put 5 V onto SYS\_3V3 \\
pin 6 & GND & GND & Common ground \\
pin 8 & UART1\_TX, PG6 & RX & Host transmit to modem receive \\
pin 10 & UART1\_RX, PG7 & TX & Host receive from modem transmit \\
onboard & AP6212 Wi-Fi and BT & & Switched off with \texttt{nmcli radio wifi off} \\
\bottomrule
\end{tabular}
\caption{The NEO Air is a 24-pin board and cannot take a 40-pin HAT. Confirm pin 1 against the silkscreen before applying power.}
\end{table}

\begin{asciiart}
NEO Air UART1 TX (pin 8 / GPIOG6)  --> ESP8266 RX
NEO Air UART1 RX (pin 10 / GPIOG7) --> ESP8266 TX
GND common, 3V3 from SYS_3V3
nmcli radio wifi off            # AP6212 down
\end{asciiart}

\subsection*{Bench layout}

\diagram{p17_bench}{Bench layout. Two small boards, four jumpers and one antenna that belongs to the ESP. The console arrives on the UART0 debug header or over a wired session, never over the radio that is switched off.}

The bench is deliberately sparse. Keep the console question settled before the radio goes down: on a headless NEO Air the console is the UART0 debug header of P09 or a wired login, because a session over the onboard Wi-Fi will end the moment the first step succeeds. That is the most common way this lab appears to fail when it has in fact worked.

\subsection*{Software design (UML)}

\diagram{p17_uml}{The AT sequence, with the host radio off for the whole run. Each command is a line on one UART, and the address that comes back belongs to the ESP.}

There is no program in this lab, only a session. The sequence is short: switch the host radio off, open \texttt{/dev/ttyS1} at 115200, check that the module answers, set station mode, join an access point, and ask for the address. Each step is a single AT line, and each answer is read by a human before the next line is typed.

Automating it later is a small job, and the shape of that job is the one P03 describes for a very different module: a serial port, one command at a time, and a parser that reads the answer rather than assuming it. Nothing here needs to be automated before the split has been demonstrated once by hand.

\subsection*{Data flow (ASCII)}

\begin{asciiart}
  NanoPi NEO Air, FriendlyElec Ubuntu          ESP8266-PROG
  +-----------------------------------+        +----------------------+
  | minicom -D /dev/ttyS1 -b 115200   |        | AT firmware          |
  |        |                          |        |                      |
  |        v                          | pin 8  |                      |
  | UART1: PG6 TX, PG7 RX  -----------+--------> RX                   |
  |                        <----------+--------- TX                   |
  |                                   | pin 10 |                      |
  | pin 1 SYS_3V3 --------------------+--------> 3V3                  |
  | pin 6 GND ------------------------+--------- GND                  |
  |                                   |        |          |           |
  | AP6212 Wi-Fi and BT: OFF          |        +----------|-----------+
  | nmcli radio wifi off              |                   v
  | ip addr shows no wireless         |          2.4 GHz to the access point
  +-----------------------------------+          AT+CIFSR returns the address

  The only IP address on this bench is inside the ESP8266.
\end{asciiart}

\subsection*{Steps}

\begin{steps}
\item \textbf{Put the onboard radio down.} The AP6212 carries both Wi-Fi and Bluetooth on this board, and the lab needs the Wi-Fi side switched off before anything else happens. Settle the console first: use the UART0 debug header as in P09, or a wired login.
\begin{shellcode}
nmcli radio wifi off
nmcli radio
iwconfig
\end{shellcode}

\item \textbf{Open UART1 and bring the modem up.} One command per line, reading each answer.
\begin{atcode}
sudo minicom -D /dev/ttyS1 -b 115200
AT
AT+CWMODE=1
AT+CWJAP="YOURAP","..."
AT+CIFSR
\end{atcode}

\item \textbf{Show the split.} Ask the host what network it has, then ask the module. The two answers disagree, and that disagreement is the lab.
\end{steps}

\subsection*{Acceptance test}

\begin{itemize}
\item \texttt{iwconfig} or \texttt{nmcli} on the NanoPi shows Wi-Fi off.
\item \texttt{AT+CIFSR} returns an IP on the ESP, in the same session.
\item That split is the lab: no interface on the Linux side holds the address.
\item \texttt{AT} answers before \texttt{AT+CWJAP} is attempted, so a join failure is a join failure and not a wiring fault.
\end{itemize}

\subsection*{Practices}

The source gives this lab no practices paragraph of its own, so the practices are its own sentences. Keep the AP6212 down for the whole session: the acceptance test is a comparison, and a host radio that comes back up quietly makes the comparison meaningless. Keep the roles apart, too. Here the ESP8266 is an AT slave on a UART; in P16 the ESP32 is a peer on the LAN with its own subscription, and the two are not interchangeable.

The NEO Air is a 24-pin board. It cannot accept any 40-pin HAT, so no cellular modem, analog HAT, console board or display belongs on this host. Its supply pin is SYS\_3V3, and P09 states the rule for it: do not put 5 V onto SYS\_3V3.

\subsection*{Pitfalls}

\begin{itemize}
\item \textbf{Working over the onboard Wi-Fi.} Step 1 ends the session. Settle the console on the UART0 debug header or a wired login before switching the radio off.
\item \textbf{Transmit wired to transmit.} Pin 8 is the host transmit and goes to the module receive. Nothing answers when the pair is not crossed, and the symptom looks like a dead module.
\item \textbf{The ESP8266-PROG USB cable still plugged in.} That is the second UART path of P14. Here it competes with UART1 for the same pins.
\item \textbf{5 V onto SYS\_3V3.} The NEO Air is a 3.3 V board. The 5 V pin is pin 2 and it is not the supply for this module.
\item \textbf{Expecting the address in \texttt{ip addr}.} It is not there and it is not meant to be. \texttt{AT+CIFSR} is the only place the address appears.
\item \textbf{Skipping the bare \texttt{AT}.} Without that answer, a failed join cannot be told apart from a wiring fault or a wrong baud rate.
\item \textbf{Reaching for a 40-pin HAT on this host.} The NEO Air header is 24-pin. SIM7020, SIM7070, SIM7600, MCC 118, Explorer700 and the 3.5 inch LCD do not seat on it.
\end{itemize}

\subsection*{Sources}

\begin{itemize}
\item FriendlyElec NanoPi NEO Air wiki, 24-pin header and UART1 on PG6 and PG7, \url{https://wiki.friendlyelec.com/wiki/index.php/NanoPi_NEO_Air}
\item Espressif ESP8266 AT command set, \texttt{AT+CWMODE}, \texttt{AT+CWJAP}, \texttt{AT+CIFSR}, \url{https://docs.espressif.com/projects/esp-at/en/latest/esp32/AT_Command_Set/}
\item JOY-iT SBC-ESP8266-PROG product page and manual, \url{https://joy-it.net/en/products/SBC-ESP8266-PROG}
\item NetworkManager \texttt{nmcli} manual, \texttt{nmcli radio wifi off}, \url{https://networkmanager.dev/docs/api/latest/nmcli.html}
\item Allwinner H3 UART documentation as published by FriendlyElec for the NEO Air
\end{itemize}

\project{20}{Four-host fleet health}{all four Linux hosts over SSH}{orchestration, no new physics}

\begin{unique}
Owns orchestration across the Pi 4, the Pi 3B+, the Pi 3 and the NanoPi. No new physics.
\end{unique}

\subsection*{Intent}

From P12, or from any Pi that has SSH keys to the others, a script of about forty lines reaches all four hosts and records what each one is: \texttt{uname}, free RAM, which 40-pin HAT EEPROM answers, and which USB gadgets are present. The output goes to \texttt{/var/log/lab-fleet.txt} on a timer, and that file is the acceptance artifact.

There is no new hardware and no new measurement in this chapter. Every board in it has already been brought up by an earlier lab, and this one only asks each board what it currently is. That is why the source puts it last in the working order: the fleet file is worth having when it reflects a known sitting, and worth very little when it reflects a bench that is half torn down.

\begin{kit}
The four Linux hosts already built in the earlier labs: Pi 4, Pi 3B+, Pi 3 and NanoPi NEO Air, on one LAN, with SSH keys from the operator host to the other three.
\end{kit}

\begin{rulebox}{This lab changes nothing}
No step here reconfigures a modem, seats a HAT, or edits a device tree. It reads and it reports. A HAT that is on the shelf is reported absent, honestly.
\end{rulebox}

\subsection*{System architecture}

\diagram{p20_arch}{The four hosts and the operator Pi. Every arrow is an SSH session that runs a read-only command and returns text, and the only thing this lab writes is one file on the operator host.}

The shape is one operator and four targets, and the operator is usually the Pi 4 from P12, because it already has a keyboard and a panel to watch the sweep on. Nothing in the design requires that host specifically: any Pi with keys to the other three will do, including a Pi that is currently wearing a HAT for a different lab, since the sweep neither needs the header nor touches it.

Three of the targets are Raspberry Pi hosts running Bookworm and one is the NanoPi NEO Air running FriendlyElec Ubuntu. That difference is visible in the output rather than in the script: \texttt{vcgencmd} exists on the Pi hosts and not on the NanoPi, and the header bus that P09 established on the NanoPi is 0 rather than 1. The script is written with error output redirected so that it stays quiet where a command does not apply, and so that the sweep completes rather than stopping at the first host that lacks a tool.

\begin{table}[H]\centering\small
\begin{tabular}{p{24mm}p{34mm}p{34mm}p{58mm}}
\toprule
Name in the script & Board & Operating system & What the sweep expects to find \\
\midrule
\texttt{pi4} & Raspberry Pi 4 & Bookworm & The broker from P12, USB gadgets, often an empty header \\
\texttt{pi3bp} & Raspberry Pi 3B+ & Bookworm & The Explorer700 from P11 if it is still seated \\
\texttt{pi3} & Raspberry Pi 3 & Bookworm & The 3.5 inch LCD from P06 or P19 if it is still seated \\
\texttt{air} & NanoPi NEO Air & FriendlyElec Ubuntu & About 512 MB of RAM, no HAT, and no \texttt{vcgencmd} \\
\bottomrule
\end{tabular}
\caption{The fleet. The four names in the loop are the four hosts of the book, and nothing else is ever added to that list.}
\end{table}

\subsection*{Wiring and schematic}

\diagram{p20_schematic}{The only wiring in this lab is the network. Each host has an address the operator can reach and a key that lets the operator in without a password prompt.}

\begin{table}[H]\centering\small
\begin{tabular}{p{34mm}p{34mm}p{82mm}}
\toprule
Connection & Medium & Note \\
\midrule
Operator to \texttt{pi4} & LAN & SSH key, no password prompt during the sweep \\
Operator to \texttt{pi3bp} & LAN & SSH key \\
Operator to \texttt{pi3} & LAN & SSH key \\
Operator to \texttt{air} & LAN over Wi-Fi & The NanoPi joined its access point in P09 \\
40-pin headers & & Read, never changed \\
\bottomrule
\end{tabular}
\caption{Connections. There are no flying leads, no jumpers and no HAT work in this chapter.}
\end{table}

\begin{asciiart}
operator Pi (P12, or any Pi with keys)
   |
   +-- ssh pi4    -> Raspberry Pi 4,    Bookworm
   +-- ssh pi3bp  -> Raspberry Pi 3B+,  Bookworm
   +-- ssh pi3    -> Raspberry Pi 3,    Bookworm
   +-- ssh air    -> NanoPi NEO Air,    FriendlyElec Ubuntu, joined in P09

Every session runs read-only commands and returns text.
No step in this lab reconfigures a modem.
\end{asciiart}

\subsection*{Bench layout}

\diagram{p20_bench}{Bench layout for a sweep. Whatever is seated on a header stays seated: the point of the file is to record the bench as it actually is.}

Run this lab last in a sitting, with the bench in the state you want recorded. If a HAT came off an hour ago, the file will say the header is empty, and it will be right. The failure mode to avoid is the opposite one: tearing down after the sweep and then treating the file as current.

There is very little to lay out. Four boards, a switch or an access point, and the operator's keyboard and panel. The only physical requirement is that all four hosts are reachable at the same moment, which for the NanoPi means the Wi-Fi connection from P09 is still saved and still joining on its own.

\subsection*{Software design (UML)}

\diagram{p20_uml}{One sweep. The operator opens a session per host in turn, each session answers with a block of text, and the four blocks are appended to \texttt{/var/log/lab-fleet.txt}.}

The sweep is a loop with no state carried between iterations. Each host is asked the same questions, each answer is text, and the operator concatenates them. That is the entire design, and it is deliberately not a monitoring system: there is no agent on the targets, nothing to install, and nothing that keeps running after the session closes.

The one piece of judgement in the script is what to do when a command is missing. \texttt{vcgencmd} is not on the NanoPi and \texttt{i2cdetect} may not be installed everywhere, so both are written with error output discarded. The sweep then produces a shorter block for that host rather than an error, and a shorter block is a true statement about that host.

\begin{note}{The file is the product}
Everything else in this chapter exists to produce \texttt{/var/log/lab-fleet.txt}. There is no dashboard, no database and no alerting. When the acceptance test is run, it is run against that file, and the question asked of it is whether it matches the bench a person can see.
\end{note}

\subsection*{Data flow (ASCII)}

\begin{asciiart}
  operator Pi                       four hosts                      artifact
  +-----------------+               +-----------------+             +----------------------+
  | for h in        |  ssh, key     | hostname        |             |                      |
  |   pi4 pi3bp     |-------------->| uptime          |             | /var/log/            |
  |   pi3 air       |               | free -h         |             |   lab-fleet.txt      |
  |                 |<--------------| vcgencmd temp   |  text       |                      |
  | append the text |   one block   | i2cdetect head  |------------>| four blocks,         |
  | to the log      |   per host    | lsusb | grep    |             | one per host         |
  +-----------------+               | /proc/.../hat   |             +----------------------+
         ^                          +-----------------+
         |
         +-- a timer runs the sweep again; the file is the acceptance artifact

  read-only everywhere: no modem is reconfigured, no HAT is seated or removed
\end{asciiart}

\subsection*{Steps}

\begin{steps}
\item \textbf{Put SSH keys on the three targets} from the operator Pi, so that the sweep runs without a password prompt. A sweep that stops to ask for a password is a sweep that cannot run on a timer.
\begin{shellcode}
ssh-keygen -t ed25519
ssh-copy-id pi4
ssh-copy-id pi3bp
ssh-copy-id pi3
ssh-copy-id air
ssh pi4 true && ssh pi3bp true && ssh pi3 true && ssh air true && echo "all four reachable"
\end{shellcode}

\item \textbf{Write the sweep.} This is the outline from the source, and it is about as long as the script needs to be.
\begin{shellcode}
for h in pi4 pi3bp pi3 air; do
  ssh $h 'hostname; uptime; free -h; vcgencmd measure_temp 2>/dev/null
          i2cdetect -y 1 2>/dev/null | head
          lsusb | grep -iE "st|sim|nordic|1a86|10c4"'
done
\end{shellcode}
The redirects are load bearing. \texttt{vcgencmd} does not exist on the NanoPi, and the header bus there is 0 rather than 1, as P09 established. Both facts show up as a shorter block for that host, which is the honest result.

\item \textbf{Detect the HAT EEPROM} at the ID bus, and report \texttt{none} when the header is empty.
\begin{shellcode}
ssh pi3bp 'test -d /proc/device-tree/hat && \
  cat /proc/device-tree/hat/product 2>/dev/null || echo none'
\end{shellcode}
The honest report is the point. A header with nothing on it should produce \texttt{none} in the file, not a blank line that could mean anything.

\item \textbf{Write the output to the log} and put the sweep on a timer. That file is the acceptance artifact.
\begin{shellcode}
sudo touch /var/log/lab-fleet.txt
sudo chown "$USER" /var/log/lab-fleet.txt
{ date -Is; ./fleet.sh; } >> /var/log/lab-fleet.txt
\end{shellcode}

\item \textbf{Read the file and check it against the bench.} Four host names, and a HAT named where one is seated and reported absent where one is not.
\begin{shellcode}
tail -n 60 /var/log/lab-fleet.txt
\end{shellcode}
\end{steps}

\subsection*{Acceptance test}

\begin{itemize}
\item Four host names appear in \texttt{/var/log/lab-fleet.txt}.
\item A HAT that is physically seated is named.
\item A HAT that is on the shelf is reported absent.
\item No step in this lab reconfigures a modem.
\end{itemize}

\subsection*{Practices}

Run P20 last, so the fleet file reflects a known sitting. That is the source's own working order, and it is the only reason the file means anything: a sweep taken in the middle of a teardown records a bench that existed for ten minutes.

Keep the lab read-only. The commands in the sweep report state and do not change it, and the moment one of them starts fixing something the file stops being a record and becomes a side effect. If a host needs work, do the work in that host's own lab and run the sweep again afterwards.

Keep the host list at four names. The fleet of this book is the Pi 4, the Pi 3B+, the Pi 3 and the NanoPi NEO Air, and every one of them was brought up by a chapter that came earlier. A fifth entry is either a duplicate or a machine this book has not built, and in both cases the file becomes harder to trust rather than more complete.

\subsection*{Pitfalls}

\begin{itemize}
\item \textbf{No SSH keys.} The sweep stops at the first password prompt, which also means it can never run on a timer.
\item \textbf{Expecting \texttt{vcgencmd} on the NanoPi.} It is a Raspberry Pi tool. The redirect is there so that its absence produces a shorter block rather than an error.
\item \textbf{Expecting bus 1 on the NanoPi.} The header bus there is 0, as P09 established. The \texttt{i2cdetect -y 1} line in the sweep is written to stay quiet where it does not apply.
\item \textbf{A blank line where a HAT should be named.} Report \texttt{none} explicitly. A blank line reads as a broken script rather than as an empty header.
\item \textbf{Running the sweep in the middle of a teardown.} The file then describes a bench that no longer exists.
\item \textbf{Adding a fifth name to the loop.} The fleet is four hosts. A fifth name is a host this book has not brought up.
\item \textbf{Letting the sweep fix things.} A repair inside the sweep makes the artifact a record of what the script did rather than of what the bench was.
\end{itemize}

\subsection*{Sources}

\begin{itemize}
\item OpenSSH documentation, \texttt{ssh(1)}, \texttt{ssh-keygen(1)} and \texttt{ssh-copy-id(1)}
\item Raspberry Pi documentation, the HAT ID EEPROM and \texttt{/proc/device-tree/hat}, \url{https://www.raspberrypi.com/documentation/computers/raspberry-pi.html}
\item Raspberry Pi \texttt{vcgencmd} documentation, for \texttt{measure\_temp}
\item \texttt{i2c-tools} manual pages, \texttt{i2cdetect(8)} and its bus argument
\item \texttt{lsusb(8)} and the USB vendor identifier list, for the 1a86 and 10c4 bridges
\item \texttt{systemd.timer(5)}, for running the sweep on a schedule
\end{itemize}


\appendix
\clearpage
\section{Lab sequence and teardown}

Work in this order if you are starting from bare boards.

\begin{enumerate}
\item P15 rails (do not skip).
\item P14 jig, then P16 and P17 sidecars.
\item P09 NanoPi bring-up.
\item P11 Explorer700 RTC, then take the HAT off.
\item P12 operator glass.
\item P05 MCC 118, P19 LCD tilt, P06 ToF kiosk, each tears the 40-pin down after acceptance.
\item P07 then P08, swapping the Nucleo shield, then P18 recorder, then P04 STWIN.
\item P01, P03, P02 cellular, then P13 failover.
\item P10 energy numbers on whatever DUT you still trust.
\item P20 last, so the fleet file reflects a known sitting.
\end{enumerate}

\begin{rulebox}{Teardown}
Power off before unseating a HAT. Discharge curiosity, not the modem PA: refit antennas before the next transmission. Put IKS4A1 and IKS5A1 back in separate bags, they look similar from across the bench.
\end{rulebox}

\section{Reference AT and addresses}

Three modules, three command sets. The full transcripts live in the labs that own them; this list is what to reach for at the bench.

\begin{description}
\item[SIM7600E-H, P01 and P13] \texttt{AT}, \texttt{AT+CPIN?}, \texttt{AT+CSQ}, \texttt{AT+COPS?}, \texttt{AT+CEREG?}, \texttt{AT+CGPS=1}, \texttt{AT+CGPSINFO}
\item[SIM7020E, P02 and P13] \texttt{AT+CBAND}, \texttt{AT+CEREG?}, \texttt{AT+CPSMS}, \texttt{AT+CEDRXS}, \texttt{AT+CSOC}, \texttt{AT+CSOCON}, \texttt{AT+CSOSEND}
\item[SIM7070G, P03] \texttt{AT+CNMP=38}, \texttt{AT+CMNB=1}, \texttt{AT+CPSI?}, \texttt{AT+CGNSPWR}, \texttt{AT+CGNSINF}, \texttt{AT+SMCONF}, \texttt{AT+SMCONN}, \texttt{AT+SMPUB}
\end{description}

The ESP8266 AT set used in P17 is a fourth list and a much shorter one: \texttt{AT}, \texttt{AT+CWMODE=1}, \texttt{AT+CWJAP}, \texttt{AT+CIFSR}.

The 7-bit I2C addresses for every sensor in the bin are in the default I2C map in the front matter. Two of them are worth repeating here because they decide a lab: the ADXL345 answers at 0x53 with SDO floating, and on the IKS5A1 the ISM6HG256X is at 0x6A while the ISM330IS is at 0x6B.

\section{Raspberry Pi 40-pin header reference}

The numbering below is the one used by every wiring table in this book. Pins 4, 7, 17, 27 and 28 are not named by any lab chapter; they follow the standard Raspberry Pi header and should be confirmed against the Raspberry Pi hardware documentation before they are used.

\begin{table}[H]\centering\small
\begin{tabular}{p{14mm}p{42mm}p{14mm}p{42mm}}
\toprule
Pin & Function & Pin & Function \\
\midrule
1 & 3V3 & 21 & GPIO9, MISO \\
2 & 5V & 22 & GPIO25 \\
3 & GPIO2, SDA1 & 23 & GPIO11, SCLK \\
4 & 5V & 24 & GPIO8, CE0 \\
5 & GPIO3, SCL1 & 25 & GND \\
6 & GND & 26 & GPIO7, CE1 \\
7 & GPIO4 & 27 & ID\_SD \\
8 & GPIO14, TXD0 & 28 & ID\_SC \\
9 & GND & 29 & GPIO5 \\
10 & GPIO15, RXD0 & 30 & GND \\
11 & GPIO17 & 31 & GPIO6 \\
12 & GPIO18 & 32 & GPIO12 \\
13 & GPIO27 & 33 & GPIO13 \\
14 & GND & 34 & GND \\
15 & GPIO22 & 35 & GPIO19 \\
16 & GPIO23 & 36 & GPIO16 \\
17 & 3V3 & 37 & GPIO26 \\
18 & GPIO24 & 38 & GPIO20 \\
19 & GPIO10, MOSI & 39 & GND \\
20 & GND & 40 & GPIO21 \\
\bottomrule
\end{tabular}
\caption{The 40-pin header. Ground is on pins 6, 9, 14, 20, 25, 30, 34 and 39.}
\end{table}

Three groups of these pins carry most of the traffic in this book. Pins 1, 3, 5 and 6 are the I2C corner. P01 reaches it through the SIM7600E-H's full 40-pin pass-through, which leaves pins 3 and 5 untouched, and that has been checked. P19 cannot reach it at all: the 3.5 inch LCD (A) mates with the first 26 pins and carries no pass-through, so pins 1, 3, 5 and 6 sit under the panel. Since pin 1 and pin 17 are the only 3.3 V pins on the header and both are inside that covered range, a sensor on that host needs its supply from somewhere other than the header.

Pins 19, 21, 23, 24 and 26 are SPI0, which is why the MCC 118, the 3.5 inch LCD and the Explorer700 cannot share a host. The MCC 118 also claims GPIO12, GPIO13 and GPIO26 as board-address pins, and the ID EEPROM lives on pins 27 and 28 whatever is seated.

\textbf{There is no universal semaphore triple, and the draft's one was the worst available choice.} It put the indicators on pins 13, 15 and 16, which on a bare Pi is fine and is what P15 uses. Under the modem HATs it is not: the SIM7600E-H drives GPIO27, GPIO22 and GPIO23 as ring indicator, data terminal ready and clear to send, and the SIM7020E drives GPIO27 as the module's power key, so an indicator there would switch the modem off and on as it reported. P01, P02, P03 and P13 therefore use BCM 16, 20 and 21 on pins 36, 38 and 40, chosen because nothing in this book documents a claim on them, with grounds on pin 34 or 39. Confirm them against the pinout of the board in front of you, because every one of these is a fact about a HAT and not about the Raspberry Pi.

\section{NanoPi NEO Air 24-pin header}

The NEO Air is a 24-pin board. It cannot accept any 40-pin HAT. Only the pins that the labs actually use are listed here; the rest of the header is outside the scope of this book.

\begin{table}[H]\centering\small
\begin{tabular}{p{16mm}p{44mm}p{86mm}}
\toprule
Pin & Signal & Used by \\
\midrule
1 & SYS\_3V3 & ADXL345 VCC in P09, ESP8266 supply in P17. Do not put 5 V on this pin \\
2 & VDD\_5V & Board supply, 5 V at 2 A. Never a sensor or module supply \\
3 & I2C0\_SDA & ADXL345 SDA in P09 \\
5 & I2C0\_SCL & ADXL345 SCL in P09 \\
6 & GND & Common ground in P09 and P17 \\
8 & UART1\_TX, PG6 & To the ESP8266 RX pin in P17 \\
10 & UART1\_RX, PG7 & From the ESP8266 TX pin in P17 \\
\bottomrule
\end{tabular}
\caption{The 24-pin header pins used in this book. I2C0, UART1, SPI0 and GPIO are what this header offers, and no 40-pin HAT seats on it.}
\end{table}

\begin{table}[H]\centering\small
\begin{tabular}{p{16mm}p{44mm}p{86mm}}
\toprule
Pin & Signal & Note \\
\midrule
1 & GND & \\
2 & 5V & \\
3 & TX & To the USB-TTL RX line \\
4 & RX & From the USB-TTL TX line \\
\bottomrule
\end{tabular}
\caption{The separate 4-pin debug header, UART0. Renkforce USB-TTL at 3.3 V logic, 115200 8N1. This is the first login in P09 and the console of choice in P17.}
\end{table}

Confirm every one of these against the silkscreen before applying power. The NEO Air is a small board with two serial headers and a 3.3 V supply pin sitting next to a 5 V one.

\section{Which lab owns which part}

One physics per lab, one exclusive part per lab. This table is the quickest way to see whether two labs can share a sitting: if they name the same part, they cannot.

\begin{table}[H]\centering\small
\begin{tabular}{p{40mm}p{16mm}p{90mm}}
\toprule
Exclusive part & Owner & Also appears in \\
\midrule
SIM7600E-H & P01 & P13 reuses the P01 attachment as the primary path \\
SIM7070G & P03 & Nowhere else. The same HAT cannot serve P02 \\
SIM7020E & P02 & P13 reuses the P02 attachment as the backup path \\
MCC 118 & P05 & Nowhere else. It is alone on the 40-pin header \\
RB-Explorer700 & P11 & Nowhere else. Do not move it onto a cellular Pi \\
3.5 inch LCD (A) & P19 & P06 uses it as the ToF heat map, in a different sitting \\
Official DSI panel & P12 & P04 and P05 use the glass. It is not a HAT \\
STEVAL-STWINBX1 & P04 & Nowhere else. There is no second STWIN product lab \\
X-NUCLEO-IKS4A1 & P07 & P18 may stack it for the optional EKF session \\
X-NUCLEO-IKS5A1 & P08 & Nowhere else. Not interchangeable with IKS4A1 \\
X-NUCLEO-53L8A1 & P06 & Nowhere else. One Arduino shield on the Nucleo \\
SEN0032 ADXL345 & P01 & P09 on I2C0 and P19 as the tilt source \\
nRF PPK2 & P10 & Nowhere else. Never sources a Cat-4 PA \\
SBC-NodeMCU-ESP32 & P16 & P10 as the DUT, P14 as a provisioning target \\
SBC-ESP8266-PROG & P17 & P14 as a provisioning target \\
NanoPi NEO Air & P09 & P17 as the air-gap host, P20 as the fourth host \\
\bottomrule
\end{tabular}
\caption{Exclusive parts and the labs that own them. The Nucleo-H7A3ZI-Q is shared by P06, P07, P08 and P18, one shield at a time.}
\end{table}

\section{Glossary}

\begin{description}
\item[AP6212] The Wi-Fi and Bluetooth part on the NanoPi NEO Air. It shares one radio between both functions, and P17 switches it off for the whole session.
\item[APN] Access point name. The operator string a cellular module needs before it can carry data. P02 takes it from the SIM sleeve.
\item[AT command] The text control language of a cellular or Wi-Fi module. Sent on a serial port, one line at a time, with an answer read before the next line.
\item[Cat-4] The LTE category of the SIM7600E-H. It carries a default IPv4 route in P01 and is the only such route in the book.
\item[Cat-M] Also eMTC. The low-power LTE mode locked in P03 with \texttt{AT+CMNB=1}. A module that falls back to GPRS fails that lab.
\item[CDC-ACM] The USB communications device class that presents a serial port to Linux as \texttt{/dev/ttyACM*}. It exists only after the firmware declares it, which is the correction P04 makes to the draft.
\item[CEREG] The AT query that reports EPS network registration. It drives the green LED in P01 and P02.
\item[CoAP] A compact request and response protocol over UDP, used in P02 because an NB-IoT node cannot afford a session-heavy stack.
\item[daqhats] The MCC userspace library that talks to the MCC 118 over SPI. It is the supported path in P05, and there is no mainline IIO driver for that board.
\item[dead-band] The window around 1 g in P01 inside which no event is published. It turns a sample stream into an event stream.
\item[DSI] The display serial interface the official Raspberry Pi touchscreen uses. It does not consume the 40-pin header, so a DSI panel may coexist with a HAT.
\item[DTO] Device tree overlay. The mechanism that binds a board-level device such as the DS3231 at boot. P11 is the overlay lab because the chips are there.
\item[eDRX] Extended discontinuous reception. A negotiated listening schedule that lets an NB-IoT module sleep between paging windows. Requested in P02 with \texttt{AT+CEDRXS}.
\item[eMMC] The soldered flash on the NanoPi NEO Air. P09 discusses its wear and keeps append-heavy logs off the rootfs.
\item[F2FS] A flash-friendly filesystem. In P09 it is an option for a data partition you created yourself, never for the rootfs you just booted.
\item[GNSS] Satellite positioning. P03 stamps its payloads with a GNSS time and powers the receiver down before transmitting.
\item[HAT] A board that seats on the Raspberry Pi 40-pin header. Six of them exist in this bin and exactly one may be fitted at a time.
\item[hwclock] The Linux tool that reads and writes a hardware real-time clock. In P11 the DS3231 on the Explorer700 is that clock, and the lab proves it by pulling the power.
\item[I2C] The two-wire bus on pins 3 and 5 of the Pi header and on pins 3 and 5 of the NanoPi header. Every address used in this book is in the default I2C map.
\item[ID EEPROM] The small memory on pins 27 and 28 that identifies a seated HAT. P20 reports it honestly, including reporting none.
\item[IIO] The Linux industrial I/O subsystem. Named here mainly to record that the MCC 118 does not use it.
\item[isolcpus] The kernel parameter that keeps a CPU free of the general scheduler. P04 uses it with \texttt{taskset} so the FFT thread is not migrated.
\item[ISPU] The intelligent sensor processing unit inside the ISM330IS on the IKS5A1. A programmable core in the sensor, not a machine learning core.
\item[MLC] Machine learning core. It lives in the ISM330DHCX on the STWIN.box, and the audit rejects any lab that puts it on the IKS5A1.
\item[MQTT] The publish and subscribe protocol used by P01, P03, P12, P13 and P16. In P16 the topic is the entire contract between the Pi and the ESP32.
\item[NB-IoT] The narrowband cellular mode of the SIM7020E. It is not a phone network, and P02 measures a duty cycle rather than a ping.
\item[PSM] Power saving mode. The state in which a cellular module stays registered but stops listening. Requested in P02 with \texttt{AT+CPSMS}.
\item[QMI] Qualcomm MSM interface. With \texttt{qmi\_wwan} it is the supported alternative to RNDIS for the P01 user plane, and it is never run beside PPP.
\item[Qvar] The electric charge variation sensing channel on the IKS4A1. P07 enables it only after the IMU stream is clean.
\item[RNDIS] A USB networking class that presents the modem as an Ethernet-like interface, typically \texttt{usb0}. The first thing P01 tries.
\item[source-measure] Supplying a device and measuring its current on the same terminals. The PPK2 does this in P10, which is why the DUT USB cable must come out.
\item[SPI] The serial bus on pins 19, 21, 23, 24 and 26. The MCC 118, the 3.5 inch LCD and the Explorer700 all want it, which is one more reason they cannot share a Pi.
\item[ToF] Time of flight ranging. The VL53L8CX in P06 returns an 8 by 8 grid of millimetre values, streamed to Linux over USB-CDC.
\item[UART] An asynchronous serial port. Three of them appear in P17 alone: the NanoPi UART1 modem path, the UART0 debug console and the USB-UART on the ESP8266-PROG.
\item[USB-CDC] The USB serial path from a microcontroller to a Linux host. It is how the Nucleo and the STWIN.box speak in P04, P06, P07, P08 and P18.
\end{description}

\end{document}
